% Sensibilisation sur l'influence des enzymes sur les réactions chimiques indispensables au renouvellement moléculaire — SVT 1ère D — Cours signé OnBuch+
% Programme officiel MINESEC. Compiler avec : tectonic ND04-influence-enzymes-reactions-chimiques.tex
\newcommand{\DOCMATIERE}{SVT}
\newcommand{\DOCNIVEAU}{1ère D}
\newcommand{\DOCMODULE}{Module 1 : Le Monde Vivant}
\newcommand{\DOCLECON}{ND04}
\newcommand{\DOCTITRE}{Sensibilisation sur l'influence des enzymes sur les réactions chimiques indispensables au renouvellement moléculaire}
% OnBuch+ — préambule « cours signés » ENRICHI (compilé avec tectonic / XeLaTeX)
% Système visuel « Pop » d'OnBuch+ : fond crème (jamais blanc), bordures encre
% épaisses, ombres « dures » décalées, titres Archivo Black en capitales.
% Version MATHÉMATIQUES : graphes (pgfplots/tikz), boîtes pédagogiques
% (définition, propriété, méthode, attention, le savais-tu, exercice résolu,
% exercice, TP, bilan), page de garde et signature OnBuch+ ancrée en bas.
%
% Macros attendues dans main.tex AVANT \input{preamble} :
%   \DOCMATIERE  \DOCNIVEAU  \DOCMODULE  \DOCLECON (numéro)  \DOCTITRE
\documentclass[11pt]{article}

\usepackage{fontspec}
\usepackage[french]{babel}
\usepackage[a4paper,margin=2cm,top=3.4cm,bottom=2.8cm,headheight=1.4cm,footskip=1.3cm]{geometry}
\usepackage{amsmath,amssymb,amsfonts}
\usepackage{mathtools} % \xmapsto, \coloneqq... souvent utilisés par les modèles
\usepackage{cancel} % simplifications barrées (bilans de forces)
% \abs{z} (module, valeur absolue) et \norm{u} (norme), utilisés par les modèles
\providecommand{\abs}{}\let\abs\relax\DeclarePairedDelimiter{\abs}{\lvert}{\rvert}
\providecommand{\norm}{}\let\norm\relax\DeclarePairedDelimiter{\norm}{\lVert}{\rVert}
\usepackage[table]{xcolor} % [table] : fournit aussi \rowcolors (alternance de lignes)
\usepackage{pagecolor}
\usepackage{tikz}
\usetikzlibrary{arrows.meta,positioning,calc,shapes.geometric,shapes.misc,shapes.arrows,decorations.pathmorphing,decorations.pathreplacing,decorations.markings,patterns,shadows,mindmap,trees,fit,backgrounds,matrix,intersections,angles,quotes}
\usepackage{pgfplots}
\usepackage[version=4]{mhchem} % équations nucléaires \ce{^{235}_{92}U}
\usepackage[european,siunitx]{circuitikz} % schémas électriques
\pgfplotsset{compat=1.18}
\usepgfplotslibrary{fillbetween,groupplots} % aires entre courbes (calcul intégral)
\usepackage{enumitem}
\usepackage{fancyhdr}
% [most] charge toutes les bibliothèques tcolorbox (listings, minted, raster,
% documentation...) et gonfle inutilement la mémoire TeX sur un document long
% (~300 boîtes) : on ne charge que ce qui est réellement utilisé.
\usepackage[skins,breakable]{tcolorbox}
\usepackage{graphicx}
\usepackage{colortbl}
\usepackage{booktabs}
\usepackage{tabularx}
\usepackage{diagbox} % case d'en-tête diagonale des tableaux à double entrée
% Colonnes extensibles centrée / gauche / droite (C, L, R), souvent utilisées par les modèles
\newcolumntype{C}{>{\centering\arraybackslash}X}
\newcolumntype{L}{>{\raggedright\arraybackslash}X}
\newcolumntype{R}{>{\raggedleft\arraybackslash}X}
\usepackage{array}
\usepackage{multicol}
\usepackage{siunitx}
% parse-numbers=false : accepte aussi \SI{2,5 \times 10^{-3}}{...}
\sisetup{locale=FR,output-decimal-marker={,},group-minimum-digits=4,parse-numbers=false}
\DeclareSIUnit\ohm{\ensuremath{\symup{\Omega}}} % U+2126 absent de la police
\DeclareSIUnit\division{div} % réglages d'oscilloscope (ms/div, V/div)
\DeclareSIUnit\tr{tr} % tours (vitesse de rotation, ex. \SI{1500}{\tr\per\minute})
\DeclareSIUnit\ch{ch} % cheval-vapeur (puissance moteur, unité usuelle hors SI)
\DeclareSIUnit\franccfa{FCFA} % franc CFA (coûts, contexte camerounais)
\DeclareSIUnit\dioptre{\ensuremath{\delta}} % vergence d'une lentille (dioptrie)
\usepackage{qrcode}
% \degree : commande fréquemment utilisée par les modèles pour un angle ou une
% température en dehors de \SI{}{} (non fournie par les paquets chargés).
\providecommand{\degree}{^{\circ}}
% \qed (fin de démonstration), utilisé par les modèles sans amsthm
\providecommand{\qed}{\hfill$\square$}
% \barre : double barre des valeurs interdites dans un tableau de variations (tkz-tab)
\providecommand{\barre}{\ensuremath{\|}}
% environnement proof (amsthm non chargé) utilisé par les modèles
\ifdefined\proof\else\newenvironment{proof}[1][Démonstration]{\par\noindent\textit{#1.}\ }{\qed\par}\fi
\usepackage{lastpage}
\usepackage{hyperref}
\hypersetup{hidelinks}

% ============================================================
% Palette « Pop » OnBuch+ — mêmes hex que la classe P (plus_ui.dart)
% ============================================================
\definecolor{poporange}{HTML}{F59321}
\definecolor{popdark}{HTML}{E07A0C}
\definecolor{popcream}{HTML}{FFF6E8}
\definecolor{popink}{HTML}{1C1714}
\definecolor{poppurple}{HTML}{7A5AE0}
\definecolor{popgreen}{HTML}{2E9E6A}
\definecolor{poppink}{HTML}{E0457B}
\definecolor{popblue}{HTML}{2F7DE1}
\definecolor{popgold}{HTML}{C98A12}
\definecolor{popmuted}{HTML}{8A7F76}
\definecolor{popline}{HTML}{EADFD2}
\definecolor{popgray}{HTML}{8A7F76}
\colorlet{popgrey}{popgray}
% Alias tolérants (le modèle nomme parfois les couleurs différemment)
\colorlet{popwhite}{white}
\colorlet{popblack}{popink}
\colorlet{popred}{poppink}
\colorlet{popyellow}{popgold}
% Teintes claires pour fonds de tableaux / zones
\colorlet{popblueL}{popblue!10!white}
\colorlet{poporangeL}{poporange!14!white}
\colorlet{popgreenL}{popgreen!12!white}
\colorlet{poppurpleL}{poppurple!12!white}
\colorlet{poppinkL}{poppink!12!white}
\colorlet{popgoldL}{popgold!14!white}

\pagecolor{popcream}

% ============================================================
% Typographie
% ============================================================
\defaultfontfeatures{Ligatures=TeX}
\setmainfont{PlusJakartaSans-Medium.ttf}[Path=../../../fonts/,BoldFont=PlusJakartaSans-Bold.ttf]
% Maths en sans-serif (Fira Math) avec chiffres et lettres droites de Plus
% Jakarta, pour que les formules se fondent dans le texte.
\usepackage[math-style=ISO,bold-style=ISO]{unicode-math}
\setmathfont{FiraMath-Regular.otf}[Path=../../../fonts/]
\setmathfont{PlusJakartaSans-Medium.ttf}[Path=../../../fonts/,range={up/num,up/{Latin,latin},\mathup}]
\usepackage{icomma} % virgule décimale sans espace en mode math
% Caractères mathématiques Unicode absents des polices de texte (Plus Jakarta,
% Archivo Black) : rendus par leur équivalent mathématique, en texte comme en
% formule — sinon ils s'affichent en carré vide (ex. « Arithmétique dans ℤ »).
\usepackage{newunicodechar}
\newunicodechar{ℝ}{\ensuremath{\mathbb{R}}}
\newunicodechar{ℤ}{\ensuremath{\mathbb{Z}}}
\newunicodechar{ℂ}{\ensuremath{\mathbb{C}}}
\newunicodechar{ℕ}{\ensuremath{\mathbb{N}}}
\newunicodechar{ℚ}{\ensuremath{\mathbb{Q}}}
\newunicodechar{𝔻}{\ensuremath{\mathbb{D}}}
\newunicodechar{α}{\ensuremath{\alpha}}
\newunicodechar{β}{\ensuremath{\beta}}
\newunicodechar{γ}{\ensuremath{\gamma}}
\newunicodechar{δ}{\ensuremath{\delta}}
\newunicodechar{ε}{\ensuremath{\varepsilon}}
\newunicodechar{θ}{\ensuremath{\theta}}
\newunicodechar{λ}{\ensuremath{\lambda}}
\newunicodechar{μ}{\ensuremath{\mu}}
\newunicodechar{σ}{\ensuremath{\sigma}}
\newunicodechar{τ}{\ensuremath{\tau}}
\newunicodechar{φ}{\ensuremath{\varphi}}
\newunicodechar{ψ}{\ensuremath{\psi}}
\newunicodechar{ω}{\ensuremath{\omega}}
\newunicodechar{Σ}{\ensuremath{\Sigma}}
% majuscules produites par \MakeUppercase dans les titres (α-aminés -> Α)
\newunicodechar{Α}{\ensuremath{\alpha}}
\newunicodechar{Β}{\ensuremath{\beta}}
\newunicodechar{Γ}{\ensuremath{\Gamma}}
\newunicodechar{Δ}{\ensuremath{\Delta}}
\newunicodechar{Ε}{\ensuremath{\varepsilon}}
\newunicodechar{Θ}{\ensuremath{\Theta}}
\newunicodechar{Λ}{\ensuremath{\Lambda}}
\newunicodechar{Μ}{\ensuremath{\mu}}
\newunicodechar{Τ}{\ensuremath{\tau}}
\newunicodechar{Φ}{\ensuremath{\Phi}}
\newunicodechar{Ψ}{\ensuremath{\Psi}}
\newunicodechar{Ω}{\ensuremath{\Omega}}
\newunicodechar{↦}{\ensuremath{\mapsto}}
\newunicodechar{∈}{\ensuremath{\in}}
\newunicodechar{∉}{\ensuremath{\notin}}
\newunicodechar{∧}{\ensuremath{\wedge}}
\newunicodechar{⇔}{\ensuremath{\Leftrightarrow}}
\newunicodechar{⟺}{\ensuremath{\Longleftrightarrow}}
\newunicodechar{⇒}{\ensuremath{\Rightarrow}}
\newunicodechar{⟹}{\ensuremath{\Longrightarrow}}
\newunicodechar{∘}{\ensuremath{\circ}}
\newunicodechar{∪}{\ensuremath{\cup}}
\newunicodechar{∩}{\ensuremath{\cap}}
\newunicodechar{≡}{\ensuremath{\equiv}}
\newunicodechar{⊂}{\ensuremath{\subset}}
\newunicodechar{⊥}{\ensuremath{\perp}}
\newunicodechar{∃}{\ensuremath{\exists}}
\newunicodechar{∀}{\ensuremath{\forall}}
\newunicodechar{∛}{\ensuremath{\sqrt[3]{\,}}}
\newunicodechar{∜}{\ensuremath{\sqrt[4]{\,}}}
\newunicodechar{⃗}{\ensuremath{{}^{\to}}}
\newunicodechar{ⁿ}{\ifmmode^{n}\else\textsuperscript{\ensuremath{n}}\fi}
\newunicodechar{ˣ}{\ifmmode^{x}\else\textsuperscript{\ensuremath{x}}\fi}
\newunicodechar{ᵇ}{\ifmmode^{b}\else\textsuperscript{\ensuremath{b}}\fi}
\newunicodechar{ʳ}{\ifmmode^{r}\else\textsuperscript{\ensuremath{r}}\fi}
\newunicodechar{ᵘ}{\ifmmode^{u}\else\textsuperscript{\ensuremath{u}}\fi}
\newunicodechar{ʸ}{\ifmmode^{y}\else\textsuperscript{\ensuremath{y}}\fi}
\newunicodechar{ᵃ}{\ifmmode^{a}\else\textsuperscript{\ensuremath{a}}\fi}
\newunicodechar{ᵉ}{\ifmmode^{e}\else\textsuperscript{\ensuremath{e}}\fi}
\newunicodechar{ᵏ}{\ifmmode^{k}\else\textsuperscript{\ensuremath{k}}\fi}
\newunicodechar{ᵖ}{\ifmmode^{p}\else\textsuperscript{\ensuremath{p}}\fi}
\newunicodechar{ᴺ}{\ifmmode^{N}\else\textsuperscript{\ensuremath{N}}\fi}
\newunicodechar{ᶜ}{\ifmmode^{c}\else\textsuperscript{\ensuremath{c}}\fi}
\newunicodechar{ʰ}{\ifmmode^{h}\else\textsuperscript{\ensuremath{h}}\fi}
\newunicodechar{ᵠ}{\ifmmode^{q}\else\textsuperscript{\ensuremath{q}}\fi}
\newunicodechar{ₙ}{\ifmmode_{n}\else\textsubscript{\ensuremath{n}}\fi}
\newunicodechar{ₐ}{\ifmmode_{a}\else\textsubscript{\ensuremath{a}}\fi}
\newunicodechar{ᵢ}{\ifmmode_{i}\else\textsubscript{\ensuremath{i}}\fi}
\newunicodechar{ᵣ}{\ifmmode_{r}\else\textsubscript{\ensuremath{r}}\fi}
\newunicodechar{ₖ}{\ifmmode_{k}\else\textsubscript{\ensuremath{k}}\fi}
\newunicodechar{ᵦ}{\ifmmode_{\beta}\else\textsubscript{\ensuremath{\beta}}\fi}
\newunicodechar{ₓ}{\ifmmode_{x}\else\textsubscript{\ensuremath{x}}\fi}
\newfontfamily\popdisplay{ArchivoBlack-Regular.ttf}[Path=../../../fonts/]
\newfontfamily\poplogo{SpaceGrotesk-Bold.ttf}[Path=../../../fonts/]
\newfontfamily\popsemi{PlusJakartaSans-SemiBold.ttf}[Path=../../../fonts/]
\newfontfamily\popxbold{PlusJakartaSans-ExtraBold.ttf}[Path=../../../fonts/]
\newfontfamily\popxboldls{PlusJakartaSans-ExtraBold.ttf}[Path=../../../fonts/,LetterSpace=3.0]

\setlist[enumerate]{leftmargin=1.7em,itemsep=5pt,topsep=4pt}
\setlist[itemize]{leftmargin=1.7em,itemsep=4pt,topsep=4pt,label={\color{poporange}\textbullet}}
\setlength{\parindent}{0pt}
\setlength{\parskip}{5pt}
\renewcommand{\arraystretch}{1.25}

% pgfplots : style Pop par défaut
\pgfplotsset{
  every axis/.append style={axis line style={popink,line width=1pt},tick style={popink},
    grid style={popline,line width=0.5pt},label style={font=\small},tick label style={font=\footnotesize}},
  popcurve/.style={poporange,line width=1.8pt,no markers,smooth},
}
% Même style disponible dans un \draw TikZ simple (hors pgfplots)
\tikzset{popcurve/.style={poporange,line width=1.8pt}}
\tikzset{popgrid/.style={popline,very thin}}
\colorlet{popcurve}{poporange} % le nom du style est parfois utilisé comme couleur

% ============================================================
% Pill / squiggle
% ============================================================
\newcommand{\poppill}[3]{%
  \tikz[baseline=-0.55ex]\node[draw=popink,line width=1.2pt,rounded corners=2.6pt,%
    fill=#2,inner xsep=6.5pt,inner ysep=3pt]{%
    \popxboldls\fontsize{7}{7}\selectfont\color{#3}\MakeUppercase{#1}};%
}
\newcommand{\popsquiggle}[1]{%
  \tikz[baseline=-0.4ex]{
    \draw[#1,line width=2.6pt,line cap=round]
      plot [smooth] coordinates {(0,0.09) (0.34,0.01) (0.68,0.21) (1.02,0.01) (1.36,0.10)};
  }%
}
% Mot-clé surligné (marqueur orange sous le texte)
\newcommand{\cle}[1]{{\popxbold\color{popdark}#1}}

% ============================================================
% En-tête / pied
% ============================================================
\pagestyle{fancy}
\fancyhf{}
\renewcommand{\headrulewidth}{0pt}
\renewcommand{\footrulewidth}{0pt}
\lhead{\raisebox{-0.15ex}{\poplogo\fontsize{13}{13}\selectfont\color{popink}OnBuch\color{poporange}+}}
\rhead{\poppill{\DOCMATIERE\ \textbullet\ \DOCNIVEAU\ \textbullet\ Leçon \DOCLECON}{white}{popink}}
\lfoot{\popxbold\fontsize{7.6}{7.6}\selectfont\color{popmuted}\MakeUppercase{Cours signé OnBuch+}\ \textbullet\ {\popsemi\DOCTITRE}}
\rfoot{\tikz[baseline=-0.6ex]\node[rounded corners=8.5pt,fill=popink,minimum height=17pt,minimum width=17pt,inner xsep=5pt,inner sep=0]{\popxbold\fontsize{8}{8}\selectfont\color{popcream}\thepage\,/\,\pageref*{LastPage}};}

% ============================================================
% Titres
% ============================================================
\newcommand{\chaptitre}[2]{%
  \begin{tcolorbox}[enhanced,colback=poporange,colframe=popink,boxrule=2.6pt,arc=11pt,
      drop shadow=popink,left=15pt,right=15pt,top=12pt,bottom=11pt,before skip=3pt,after skip=18pt]
    \poppill{#2}{popink}{popcream}\par\vspace{9pt}
    {\raggedright\hyphenpenalty=10000\exhyphenpenalty=10000\popdisplay\fontsize{22}{24}\selectfont\color{popcream}\MakeUppercase{#1}\par}
  \end{tcolorbox}
}
% Section numérotée : pastille orange avec le numéro + titre Archivo
\newcounter{popsec}
\newcommand{\coursec}[1]{%
  \stepcounter{popsec}\setcounter{popsub}{0}%
  \par\vspace{16pt}\noindent
  \begin{minipage}{\linewidth}
  \tikz[baseline=-0.7ex]\node[draw=popink,line width=1.6pt,fill=poporange,rounded corners=4pt,minimum size=22pt,inner sep=0]{\popdisplay\fontsize{12}{12}\selectfont\color{popcream}\thepopsec};%
  \hspace{8pt}{\popdisplay\fontsize{15}{17}\selectfont\color{popink}\MakeUppercase{#1}}\par
  \vspace{4pt}\noindent{\color{poporange}\rule{36pt}{3.6pt}}
  \end{minipage}\par\nopagebreak\vspace{8pt}
}
\newcounter{popsub}
\newcommand{\courssub}[1]{%
  \stepcounter{popsub}%
  \par\vspace{9pt}\noindent{\popxbold\fontsize{11.5}{13}\selectfont\color{popink}{\color{poporange}\thepopsec.\thepopsub}\hspace{6pt}#1}\par\nopagebreak\vspace{4pt}
}

% ============================================================
% Boîtes pédagogiques — même recette : fond blanc, bordure épaisse colorée,
% ombre dure encre, badge plein. Titre optionnel [#1].
% ============================================================
\tcbset{popbase/.style={breakable,enhanced,colback=white,boxrule=2.3pt,arc=9pt,
  shadow={3pt}{-3pt}{0pt}{popink},left=14pt,right=14pt,top=11pt,bottom=11pt,before skip=11pt,after skip=15pt,
  fontupper=\popsemi\fontsize{10.6}{14.5}\selectfont\color{popink},
  fontlower=\popsemi\fontsize{10.6}{14.5}\selectfont\color{popink}}}
% Coche dessinée (le glyphe ✓ n'existe pas dans les polices du document)
\newcommand{\popcheck}[1]{\tikz[baseline=-0.5ex]\draw[#1,line width=1.6pt,line cap=round,line join=round](-0.13,0.0)--(-0.04,-0.09)--(0.13,0.1);}
\providecommand{\coche}{\popcheck{popgreen}}
\providecommand{\resultat}[1]{\par\smallskip\noindent\fcolorbox{popgreen}{popgreenL}{\parbox{\dimexpr\linewidth-2\fboxsep-2\fboxrule\relax}{\textbf{Résultat.} #1}}\par\smallskip}
\newcommand{\popboxhead}[3]{\poppill{#1}{#2}{white}\ifx\relax#3\relax\else\hspace{6pt}{\popxbold\fontsize{10.6}{12}\selectfont\color{popink}#3}\fi\par\vspace{7pt}}

\newtcolorbox{aretenir}[1][]{popbase,colframe=popgreen,before upper={\popboxhead{À retenir}{popgreen}{#1}}}
\newtcolorbox{exemplebox}[1][]{popbase,colframe=poporange,before upper={\popboxhead{Exemple}{poporange}{#1}}}
\newtcolorbox{definition}[1][]{popbase,colframe=popblue,before upper={\popboxhead{Définition}{popblue}{#1}}}
\newtcolorbox{propriete}[1][]{popbase,colframe=poppurple,before upper={\popboxhead{Propriété}{poppurple}{#1}}}
\newtcolorbox{methode}[1][]{popbase,colframe=popgold,before upper={\popboxhead{Méthode}{popgold}{#1}}}
\newtcolorbox{attention}[1][]{popbase,colframe=poppink,before upper={\popboxhead{Attention}{poppink}{#1}}}
\newtcolorbox{experience}[1][]{popbase,colframe=popblue,colback=popblueL,before upper={\popboxhead{Expérience}{popblue}{#1}}}
% « Le savais-tu ? » : carte violette pleine (contexte, culture, Cameroun)
\newtcolorbox{savaistu}[1][]{popbase,colback=poppurple,colframe=popink,
  fontupper=\popsemi\fontsize{10.4}{14.2}\selectfont\color{white},
  before upper={\poppill{Le savais-tu\,?}{white}{poppurple}\ifx\relax#1\relax\else\hspace{6pt}{\popxbold\color{white}#1}\fi\par\vspace{7pt}}}
% Exercice résolu : énoncé en haut, \tcblower puis la solution sur fond crème
\newcounter{popexo}\newcounter{popexr}
\newtcolorbox{exoresolu}[1][]{popbase,colframe=poporange,colbacklower=popcream,
  segmentation style={popink,line width=1.2pt,dashed},
  before upper={\stepcounter{popexr}\popboxhead{Exercice résolu \thepopexr}{poporange}{#1}},
  before lower={\poppill{Solution}{popink}{popcream}\par\vspace{6pt}}}
% Exercice (non résolu en place) — la correction est dans la partie Corrigés
\newtcolorbox{exercice}[1][]{popbase,colframe=popink,
  before upper={\stepcounter{popexo}\popboxhead{Exercice \thepopexo}{popink}{#1}}}
% Corrigé
\newtcolorbox{corrige}[1][]{popbase,colframe=popgreen,colback=popgreenL,
  before upper={\popboxhead{Corrigé}{popgreen}{#1}}}
% Objectifs / prérequis / bilan
\newtcolorbox{objectifs}{popbase,colframe=popink,colback=poporangeL,
  before upper={\popboxhead{Objectifs de la leçon}{popink}{}}}
\newtcolorbox{prerequis}{popbase,colframe=popink,colback=popblueL,
  before upper={\popboxhead{Prérequis}{popblue}{}}}
\newtcolorbox{bilan}{popbase,colframe=popink,colback=white,boxrule=2.8pt,
  before upper={\popboxhead{Fiche bilan}{poporange}{L'essentiel à connaître pour l'examen}}}
% Figure encadrée avec légende
\newtcolorbox{popfigure}[1][]{enhanced,breakable=false,colback=white,colframe=popline,boxrule=1.4pt,arc=8pt,
  left=10pt,right=10pt,top=10pt,bottom=8pt,before skip=10pt,after skip=12pt,halign=center,
  lower separated=false,halign lower=center,
  fontlower=\popsemi\fontsize{9}{11.5}\selectfont\color{popmuted},
  #1}
\newcommand{\legende}[1]{\par\vspace{6pt}{\popsemi\fontsize{9}{11.5}\selectfont\color{popmuted}#1}}

% ============================================================
% Signature OnBuch+
% ============================================================
\newcommand{\poplogoinline}[1]{{\poplogo\fontsize{#1}{#1}\selectfont\color{popink}OnBuch\color{poporange}+}}
\newcommand{\signatureonbuch}{%
  \begin{tcolorbox}[enhanced,colback=popink,colframe=popink,boxrule=2.2pt,arc=10pt,
      drop shadow=poporange,left=14pt,right=14pt,top=10pt,bottom=10pt,before skip=0pt,after skip=0pt]
    \tikz[baseline=-0.7ex]\node[circle,fill=popgreen,draw=popcream,line width=1.3pt,minimum size=22pt,inner sep=0]{\popcheck{white}};%
    \hspace{9pt}\begin{minipage}[c]{0.62\linewidth}
      {\popxbold\fontsize{12}{13}\selectfont\color{popcream}Signé\ \textbullet\ {\poplogo OnBuch\color{poporange}+}}\par\vspace{2pt}
      {\raggedright\popsemi\fontsize{8}{10.5}\selectfont\color{popline}\MakeUppercase{Équipe pédagogique OnBuch+}\\\MakeUppercase{Conforme au programme officiel MINESEC}\par}
    \end{minipage}\hfill
    {\poplogo\fontsize{22}{22}\selectfont\color{popcream}OnBuch\color{poporange}+}
  \end{tcolorbox}%
}

% ============================================================
% Page de garde
% #1 titre · #2 matière · #3 niveau · #4 module · #5 n° leçon · #6 durée ·
% #7 description / objectifs (texte ou itemize)
% ============================================================
\newcommand{\pagegarde}[7]{%
  \thispagestyle{empty}
  \newgeometry{a4paper,margin=1.7cm,top=1.6cm,bottom=1.4cm}%
  \begin{tikzpicture}[remember picture,overlay]
    \fill[poporange!18] ([xshift=-3.2cm,yshift=-3.6cm]current page.north east) circle (4.6cm);
    \fill[poppurple!14] ([xshift=2.4cm,yshift=7.6cm]current page.south west) circle (3.8cm);
  \end{tikzpicture}%
  {\poplogo\fontsize{30}{30}\selectfont\color{popink}OnBuch\color{poporange}+}\hfill
  \raisebox{0.6ex}{\poppill{Cours signé \textbullet\ Premium}{popink}{popcream}}\par
  \vspace{1pt}\popsquiggle{poppurple}\par
  \vspace{8pt}
  \begin{tcolorbox}[enhanced,colback=poporange,colframe=popink,boxrule=2.8pt,arc=13pt,
      drop shadow=popink,left=18pt,right=18pt,top=12pt,bottom=13pt,after skip=10pt]
    \poppill{#2 \textbullet\ #3}{popink}{popcream}\hspace{5pt}\poppill{#4}{white}{popink}\par\vspace{8pt}
    {\popdisplay\fontsize{14}{14}\selectfont\color{popink}LEÇON #5}\par\vspace{4pt}
    {\raggedright\hyphenpenalty=10000\exhyphenpenalty=10000\popdisplay\fontsize{25}{27}\selectfont\color{popcream}\MakeUppercase{#1}\par}
    \vspace{8pt}
    {\popxbold\fontsize{9.5}{9.5}\selectfont\color{popink}\MakeUppercase{Durée conseillée : #6}}
  \end{tcolorbox}
  \begin{tcolorbox}[enhanced,colback=white,colframe=popink,boxrule=2.2pt,arc=10pt,
      drop shadow=popink,left=16pt,right=16pt,top=10pt,bottom=10pt,after skip=8pt]
    \poppill{Dans cette leçon}{poppurple}{white}\par\vspace{5pt}
    \popsemi\fontsize{9.3}{12.6}\selectfont\color{popink}#7
  \end{tcolorbox}
  \noindent\begin{minipage}[t]{0.487\linewidth}
    \begin{tcolorbox}[enhanced,colback=popgreen,colframe=popink,boxrule=2.2pt,arc=10pt,
        drop shadow=popink,left=11pt,right=11pt,top=10pt,bottom=10pt,height=3.1cm,valign=center]
      \tcbox[colback=white,colframe=popink,boxrule=1.3pt,arc=5pt,boxsep=0pt,nobeforeafter,
        left=3pt,right=3pt,top=3pt,bottom=3pt,valign=center]{\qrcode[height=1.9cm]{https://play.google.com/store/apps/details?id=com.luvvix.onbuch&hl=fr}}%
      \hspace{8pt}\begin{minipage}[c]{0.5\linewidth}
        {\popdisplay\fontsize{11}{12.5}\selectfont\color{white}\MakeUppercase{Télécharge OnBuch}}\par\vspace{4pt}
        {\raggedright\popsemi\fontsize{8.4}{11}\selectfont\color{white}Gratuit \textbullet\ Google Play\\Tuteur IA, annales, concours, résultats\par}
      \end{minipage}
    \end{tcolorbox}
  \end{minipage}\hfill
  \begin{minipage}[t]{0.487\linewidth}
    \begin{tcolorbox}[enhanced,colback=poppurple,colframe=popink,boxrule=2.2pt,arc=10pt,
        drop shadow=popink,left=11pt,right=11pt,top=10pt,bottom=10pt,height=3.1cm,valign=center]
      \tikz[baseline=-0.7ex]\node[circle,fill=white,draw=popink,line width=1.3pt,minimum size=1.6cm,inner sep=0]{\popdisplay\fontsize{24}{24}\selectfont\color{poppurple}+};%
      \hspace{8pt}\begin{minipage}[c]{0.6\linewidth}
        {\popdisplay\fontsize{11}{12.5}\selectfont\color{white}\MakeUppercase{Passe à OnBuch+}}\par\vspace{4pt}
        {\raggedright\popsemi\fontsize{8.4}{11}\selectfont\color{white}Tous les cours signés, Tuteur IA illimité, fiches PDF — onglet OnBuch+ de l'app\par}
      \end{minipage}
    \end{tcolorbox}
  \end{minipage}
  \vspace{8pt}
  \popcontenu
  \vfill
  \signatureonbuch
  \restoregeometry
  \newpage
}

% Rangée « Ce cours contient » (page de garde)
\newcommand{\poptile}[3]{% #1 couleur #2 gros texte #3 légende
  \begin{tcolorbox}[enhanced,colback=#1,colframe=popink,boxrule=2pt,arc=9pt,shadow={2.5pt}{-2.5pt}{0pt}{popink},
      left=6pt,right=6pt,top=9pt,bottom=9pt,halign=center,valign=center,height=1.75cm,width=\linewidth,nobeforeafter]
    {\popdisplay\fontsize{12.5}{13}\selectfont\color{white}\MakeUppercase{#2}}\par\vspace{3pt}
    {\centering\popsemi\fontsize{7.8}{9.5}\selectfont\color{white}#3\par}
  \end{tcolorbox}}
\newcommand{\popcontenu}{%
  \par\noindent{\popxboldls\fontsize{7.5}{7.5}\selectfont\color{popmuted}CE COURS SIGNÉ CONTIENT}\par\vspace{7pt}
  \noindent\begin{minipage}[t]{0.235\linewidth}\poptile{popblue}{Cours}{complet, illustré, pas à pas}\end{minipage}\hfill
  \begin{minipage}[t]{0.235\linewidth}\poptile{poporange}{Méthodes}{et exercices résolus}\end{minipage}\hfill
  \begin{minipage}[t]{0.235\linewidth}\poptile{poppink}{Exercices}{type examen + corrigés}\end{minipage}\hfill
  \begin{minipage}[t]{0.235\linewidth}\poptile{popgreen}{Fiche bilan}{l'essentiel à retenir}\end{minipage}\par}

% Sommaire
\newtcolorbox{sommaire}{enhanced,colback=white,colframe=popink,boxrule=2.2pt,arc=10pt,
  drop shadow=popink,left=18pt,right=18pt,top=13pt,bottom=13pt,before skip=6pt,after skip=20pt,
  before upper={\poppill{Sommaire}{poppurple}{white}\par\vspace{9pt}\popsemi\fontsize{10.6}{14.5}\selectfont\color{popink}}}

% Signature de fin de cours — poussée tout en bas de la dernière page
\newcommand{\findecours}{%
  \par\vfill\nopagebreak
  \begin{center}\popsquiggle{poporange}\end{center}\vspace{4pt}
  \signatureonbuch
}
\AtBeginDocument{\def\llbracket{\mathopen{[\mkern-3.5mu[}}\def\rrbracket{\mathclose{]\mkern-3.5mu]}}} % intervalles d'entiers (glyphe ⟦ absent de la police)
% Glyphes absents de Fira Math : redéfinis à partir de glyphes disponibles
\AtBeginDocument{%
  \def\setminus{\mathbin{\backslash}}\let\smallsetminus\setminus
  \def\vdots{\mathord{\vbox{\baselineskip3.2pt\lineskiplimit0pt\kern4pt\hbox{.}\hbox{.}\hbox{.}}}}%
  \def\ddots{\mathinner{\mkern1mu\raise6.5pt\hbox{.}\mkern2mu\raise3.6pt\hbox{.}\mkern2mu\raise0.7pt\hbox{.}\mkern1mu}}%
  \def\triangle{\mathord{\scalebox{0.9}{$\symup{\Delta}$}}}%
  \def\nabla{\mathord{\raisebox{\depth}{\scalebox{1}[-1]{$\symup{\Delta}$}}}}%
  \def\checkmark{\popcheck{popdark}}%
  \def\star{\mathbin{\ast}}\def\bigstar{\mathord{\ast}}%
  \def\diamond{\mathbin{\scalebox{0.6}{$\Diamond$}}}%
  \def\bigcup{\mathop{\mathchoice{\raisebox{-0.3em}{\scalebox{1.7}{$\cup$}}}{\scalebox{1.25}{$\cup$}}{\cup}{\cup}}\displaylimits}%
  \def\bigcap{\mathop{\mathchoice{\raisebox{-0.3em}{\scalebox{1.7}{$\cap$}}}{\scalebox{1.25}{$\cap$}}{\cap}{\cap}}\displaylimits}%
  \def\lesseqqgtr{\mathrel{\lessgtr}}\def\bigoplus{\mathop{\oplus}\displaylimits}%
}
\providecommand{\ssi}{si et seulement si} % abréviation fréquente
\AtBeginDocument{\providecommand{\wideparen}{\overparen}} % arc de cercle

%% FIGFIX-BEGIN v4 — correctifs globaux des figures (docs/onbuch-generation/tools/figfix)
\usepackage{adjustbox}\usepackage{etoolbox}
% toute figure TikZ/pgfplots plus large que la ligne est réduite à la largeur disponible
\BeforeBeginEnvironment{tikzpicture}{\begin{adjustbox}{max width=\linewidth}}
\AfterEndEnvironment{tikzpicture}{\end{adjustbox}}
% légendes pgfplots lisibles par-dessus les courbes
\pgfplotsset{every axis/.append style={legend style={fill=white,fill opacity=0.92,text opacity=1,draw=black!15,inner sep=3pt}}}
% étiquettes d'arêtes (syntaxe edge["..."]) avec fond blanc
\tikzset{every edge quotes/.append style={fill=white,inner sep=1.2pt}}
%% FIGFIX-END
\begin{document}

\pagegarde{Sensibilisation sur l'influence des enzymes sur les réactions chimiques indispensables au renouvellement moléculaire}{SVT}{1ère D}{Module 1 : Le Monde Vivant}{ND04}{12 h}{Cette leçon introduit la notion d'enzyme, catalyseur biologique de nature protéique, et explique comment les enzymes accélèrent les réactions chimiques du métabolisme cellulaire. Elle montre comment la température, le pH et les concentrations influencent leur activité, et relie les déficits enzymatiques à des pathologies métaboliques concrètes.
\vspace{4pt}
\begin{multicols}{2}\small
\begin{itemize}
\item À la fin de cette leçon, je sais définir une enzyme, un substrat et un produit de réaction enzymatique.
\item À la fin de cette leçon, je sais expliquer la spécificité enzymatique à l'aide du modèle clé-serrure et du site actif.
\item À la fin de cette leçon, je sais exploiter une courbe d'activité enzymatique en fonction de la température pour déterminer la température optimale.
\item À la fin de cette leçon, je sais exploiter une courbe d'activité enzymatique en fonction du pH pour déterminer le pH optimal.
\item À la fin de cette leçon, je sais interpréter une expérience mettant en évidence la spécificité d'une enzyme (ex. amylase sur amidon vs cellulose).
\item À la fin de cette leçon, je sais distinguer le rôle des enzymes dans les réactions de synthèse (anabolisme) et de dégradation (catabolisme).
\end{itemize}\end{multicols}}

\begin{sommaire}
\begin{multicols}{2}
\begin{enumerate}
\item Notion d'enzyme et de catalyse enzymatique
\item Spécificité enzymatique et site actif
\item Influence de la température sur l'activité enzymatique
\item Influence du pH et des concentrations
\item Rôle des enzymes dans le renouvellement moléculaire
\item Déficits enzymatiques et pathologies métaboliques
\item Activité d'intégration
\item Méthodes et exercices résolus
\item Exercices
\item Corrigés des exercices
\item Fiche bilan
\end{enumerate}
\end{multicols}
\end{sommaire}

\begin{objectifs}
\begin{itemize}
\item À la fin de cette leçon, je sais définir une enzyme, un substrat et un produit de réaction enzymatique.
\item À la fin de cette leçon, je sais expliquer la spécificité enzymatique à l'aide du modèle clé-serrure et du site actif.
\item À la fin de cette leçon, je sais exploiter une courbe d'activité enzymatique en fonction de la température pour déterminer la température optimale.
\item À la fin de cette leçon, je sais exploiter une courbe d'activité enzymatique en fonction du pH pour déterminer le pH optimal.
\item À la fin de cette leçon, je sais interpréter une expérience mettant en évidence la spécificité d'une enzyme (ex. amylase sur amidon vs cellulose).
\item À la fin de cette leçon, je sais distinguer le rôle des enzymes dans les réactions de synthèse (anabolisme) et de dégradation (catabolisme).
\item À la fin de cette leçon, je sais relier un déficit ou un dérèglement enzymatique à une pathologie métabolique (intolérance au lactose, phénylcétonurie, albinisme).
\item À la fin de cette leçon, je sais interpréter l'effet de la concentration en substrat ou en enzyme sur la vitesse de réaction.
\end{itemize}
\end{objectifs}

\begin{prerequis}
\begin{itemize}
\item Notion de réaction chimique et d'équation-bilan (physique-chimie, classe de Seconde)
\item Structure de la cellule et notion de métabolisme (leçons précédentes du Module 1)
\item Notion de protéine et de liaison peptidique (constituants de la matière vivante)
\item Lecture et exploitation de graphiques cartésiens (mathématiques)
\item Notion de catalyseur vue en physique-chimie (accélération d'une réaction sans être consommé)
\end{itemize}
\end{prerequis}

\newpage

\coursec{Notion d'enzyme et de catalyse enzymatique}

\courssub{Situation d'accroche et problématique}

Imagine Aminata, élève en Première D au lycée de Yaoundé. Un après-midi, elle prépare un jus de goyave frais pour sa famille. Elle en verse une partie dans un verre qu'elle laisse sur la table, et l'autre partie elle la fait chauffer quelques minutes avant de la mettre au réfrigérateur. Le lendemain, elle constate avec surprise : le jus laissé à température ambiante a bruni et a un goût fermenté, tandis que le jus chauffé est resté clair et conserve son goût sucré. Sa grand-mère, observant la scène, lui dit : « Le feu a tué ce qui gâtait le jus. »

Par ailleurs, au centre de santé du quartier, on signale plusieurs cas de nourrissons qui vomissent et ont des diarrhées violentes après chaque tétée. Le diagnostic tombe : \textbf{intolérance au lactose} due à un déficit en \textbf{lactase}, l'enzyme qui digère le sucre du lait.

\begin{center}
\textbf{Problématique :} Comment de simples molécules, les enzymes, contrôlent-elles la vitesse des réactions chimiques indispensables à la vie ? Pourquoi la chaleur « tue » l'action qui fait brunir le jus, et pourquoi l'absence d'une seule enzyme rend-elle le lait toxique pour certains bébés ?
\end{center}

Cette section va nous permettre de définir ce qu'est une enzyme, de comprendre son rôle de catalyseur biologique et de poser les bases de son mode d'action.

\courssub{Mise en évidence expérimentale : l'amylase salivaire et l'amidon}

Commençons par une expérience classique, réalisable au laboratoire comme à la maison, qui met en évidence l'action d'une enzyme.

\begin{experience}[Mise en évidence de l'activité de l'amylase salivaire sur l'amidon]
\textbf{Matériel :} tubes à essai, amidon (fécule de maïs ou manioc), eau distillée, salive humaine (ou extrait de salive), solution d'iode (réactif de Lugol), bain-marie à 37\,°C.

\textbf{Protocole :}
\begin{enumerate}
\item Préparer une solution d'amidon à 1\,\% (1\,g dans 100\,mL d'eau bouillie, refroidie).
\item Répartir la solution dans trois tubes :
\begin{itemize}
\item Tube T1 (témoin négatif) : 5\,mL d'amidon + 1\,mL d'eau distillée.
\item Tube T2 (expérimental) : 5\,mL d'amidon + 1\,mL de salive fraîche.
\item Tube T3 (témoin chaleur) : 5\,mL d'amidon + 1\,mL de salive \textbf{chauffée à 100\,°C pendant 5 min} (salive dénaturée).
\end{itemize}
\item Incuber les trois tubes au bain-marie à 37\,°C pendant 15 min.
\item Ajouter 3 gouttes de solution d'iode dans chaque tube et observer la coloration.
\end{enumerate}

\textbf{Observations :}
\begin{itemize}
\item T1 : coloration \textbf{bleu-noir} intense (présence d'amidon).
\item T2 : coloration \textbf{jaune-orangé} (absence d'amidon, formation de sucres réducteurs).
\item T3 : coloration \textbf{bleu-noir} (amidon non hydrolysé).
\end{itemize}

\textbf{Interprétation :} La salive contient une substance (l'\textbf{amylase salivaire} ou \textbf{ptyaline}) qui transforme l'amidon en sucres simples (maltose, glucose) à 37\,°C. Cette substance est \textbf{thermosensible} : la chaleur la détruit (dénaturation), abolissant son activité. L'eau seule ne transforme pas l'amidon.
\end{experience}

\begin{popfigure}
\begin{tikzpicture}[
    box/.style={rectangle, draw=popdark, thick, fill=popcream, rounded corners, minimum width=2.5cm, minimum height=1.2cm, align=center},
    arrow/.style={-Stealth, thick, popdark},
    label/.style={font=\small, text=popdark}
]
% Tubes
\node[box, align=center] (T1) at (0,0){Tube T1\\Amidon + Eau};
\node[box, align=center] (T2) at (5,0){Tube T2\\Amidon + Salive};
\node[box, align=center] (T3) at (10,0){Tube T3\\Amidon + Salive bouillie};

% Arrows down to results
\draw[arrow] (T1.south) -- ++(0,-1.2) node[below, label, align=center]{+ Iode $\rightarrow$ \textbf{Bleu-noir}\\(Amidon présent)};
\draw[arrow] (T2.south) -- ++(0,-1.2) node[below, label, align=center]{+ Iode $\rightarrow$ \textbf{Jaune-orangé}\\(Amidon hydrolysé)};
\draw[arrow] (T3.south) -- ++(0,-1.2) node[below, label, align=center]{+ Iode $\rightarrow$ \textbf{Bleu-noir}\\(Amidon intact)};

% Top bar with label (replaces brace decoration)
\draw[thick, popblue] (T1.north west) -- (T3.north east) node[midway, above=10pt, font=\bfseries, text=popblue] {Incubation à 37\,°C pendant 15 min};
\draw[thick, popblue] (T1.north west) -- ++(0,0.3) (T3.north east) -- ++(0,0.3);
\end{tikzpicture}
\legende{Expérience classique mettant en évidence l'activité hydrolytique de l'amylase salivaire sur l'amidon. Le témoin T1 montre que l'amidon ne s'hydrolyse pas spontanément. Le tube T3 prouve que l'agent actif de la salive est thermolabile (protéine).}
\end{popfigure}

\courssub{Définition et nature chimique des enzymes}

L'expérience ci-dessus révèle l'existence de \textbf{catalyseurs biologiques} produits par les cellules vivantes.

\begin{definition}[Enzyme]
Une \textbf{enzyme} est un \textbf{catalyseur biologique de nature protéique} (protéine globulaire, sauf rares ARN catalytiques ou \emph{ribozymes}), synthétisée par les cellules vivantes, qui accélère une réaction chimique précise sans être consommée ni modifiée définitivement au cours de cette réaction.
\end{definition}

\begin{propriete}[Propriétés fondamentales d'un catalyseur (admises, justifiées expérimentalement)]
Un catalyseur (minéral ou enzymatique) possède trois propriétés cardinales :
\begin{enumerate}
\item \textbf{Accélération :} Il augmente la vitesse d'atteinte de l'équilibre chimique (sans déplacer cet équilibre).
\item \textbf{Non-consommation :} Il n'apparaît pas dans le bilan global de la réaction ; il est retrouvé intact à la fin.
\item \textbf{Faible quantité :} Il agit en quantités infinitésimales par rapport aux substrats (numéro de tour élevé).
\end{enumerate}
\end{propriete}

\begin{attention}[Erreur fréquente : « L'enzyme est consommée par la réaction »]
\textbf{FAUX.} L'enzyme \textbf{se lie temporairement} au substrat, le transforme, puis \textbf{se libère intacte} et peut recommencer un nouveau cycle catalytique. Elle n'est « perdue » que si elle est \textbf{dénaturée} (chaleur, pH extrême, agents chimiques) ou dégradée par des protéases. Au Probatoire, écrire « l'enzyme est consommée » est une erreur éliminatoire.
\end{attention}

\courssub{Vocabulaire de la réaction enzymatique}

Toute réaction enzymatique met en jeu des acteurs précis dont il faut maîtriser les noms.

\begin{definition}[Substrat, Produit, Réaction enzymatique]
\begin{itemize}
\item \textbf{Substrat (S) :} Molécule(s) sur laquelle (lesquelles) agit l'enzyme. C'est le réactif spécifique.
\item \textbf{Produit (P) :} Molécule(s) formée(s) à l'issue de la transformation catalysée.
\item \textbf{Réaction enzymatique :} Transformation du substrat en produit(s) catalysée par une enzyme E.
\end{itemize}
On note schématiquement : $\ce{E + S <=> ES -> E + P}$ où \textbf{ES} est le \textbf{complexe enzyme-substrat} (intermédiaire transitoire).
\end{definition}

\begin{popfigure}
\begin{tikzpicture}[
    node/.style={circle, draw=popdark, thick, minimum size=1.2cm, fill=popblueL, font=\bfseries},
    complex/.style={rectangle, draw=poporange, thick, fill=poporangeL, rounded corners, minimum width=2.5cm, minimum height=1.5cm, align=center},
    arrow/.style={-Stealth, thick, popdark},
    label/.style={font=\small, text=popdark}
]
% Reactants
\node[node, align=center] (E) at (0,0){E\\Enzyme};
\node[node, align=center] (S) at (0,-2.5){S\\Substrat};

% Complex
\node[complex, align=center] (ES) at (4.5,-1.25){Complexe\\Enzyme-Substrat\\(ES)\\$\downarrow$\\État de transition};

% Products
\node[node, align=center] (E2) at (9,0){E\\Enzyme\\libre};
\node[node, align=center] (P) at (9,-2.5){P\\Produit};

% Arrows
\draw[arrow] (E.east) -- (ES.west) node[midway, above, label] {Reconnaissance};
\draw[arrow] (S.east) -- (ES.west) node[midway, below, label] {Liaison};
\draw[arrow] (ES.east) -- (E2.west) node[midway, above, label] {Libération};
\draw[arrow] (ES.east) -- (P.west) node[midway, below, label] {Libération};

% Energy profile sketch (small)
\begin{scope}[shift={(12,-1.25)}, scale=0.6]
\draw[->] (-1,0) -- (4,0) node[right] {Coordonnée de réaction};
\draw[->] (0,-1) -- (0,2.5) node[above] {Énergie};
\draw[thick, popdark] (-0.5,1.5) .. controls (1,2) and (2,2) .. (3,0.5);
\draw[thick, poporange, dashed] (-0.5,1.5) .. controls (1,0.5) and (2,0.5) .. (3,0.5);
\node[label] at (1.5,2.2) {Sans enzyme};
\node[label, text=poporange] at (1.5,0.2) {Avec enzyme};
\node[label] at (-0.5,1.7) {E+S};
\node[label] at (3,0.7) {E+P};
\end{scope}
\end{tikzpicture}
\legende{Schéma du cycle catalytique : l'enzyme (E) se lie au substrat (S) pour former le complexe ES, abaisse l'énergie d'activation, transforme S en P, puis libère P et se retrouve libre pour un nouveau cycle. Le profil d'énergie (en haut à droite) illustre l'abaissement de l'énergie d'activation.}
\end{popfigure}

\courssub{Nomenclature et classification des enzymes}

La nomenclature officielle (UICN) classe les enzymes en 6 classes principales selon le type de réaction catalysée (Oxydoréductases, Transférases, Hydrolases, Lyases, Isomérases, Ligases). Au niveau Probatoire, on retient la règle pratique du \textbf{suffixe « -ase »}.

\begin{propriete}[Règle de nommage usuelle (à connaître)]
Le nom de l'enzyme se construit généralement en ajoutant le suffixe \textbf{-ase} au nom du \textbf{substrat} ou au \textbf{type de réaction} :
\begin{center}
\begin{tabular}{|l|l|l|}\hline
\rowcolor{popblueL}
\textbf{Enzyme} & \textbf{Substrat / Réaction} & \textbf{Exemple physiologique} \\\hline
\cle{Amylase} & Amidon (hydrolyse) & Salive, pancréas (digestion glucides) \\\hline
\cle{Lipase} & Lipides (triglycérides) & Suc pancréatique, jus gastrique (digestion lipides) \\\hline
\cle{Protéase} / \cle{Peptidase} & Protéines / Peptides & Pepsine (estomac), Trypsine (intestin) \\\hline
\cle{Lactase} & Lactose (sucre du lait) & Bordure en brosse intestinale \\\hline
\cle{Saccharase} / \cle{Invertase} & Saccharose & Levures, intestin \\\hline
\cle{Catalase} & Eau oxygénée ($\ce{H2O2}$) & Peroxysomes (détotoxication) \\\hline
\cle{Transaminase} & Transfert groupe amino & Foie, muscle (métabolisme acides aminés) \\\hline
\cle{Kinase} & Transfert phosphate (ATP) & Glycolyse, signalisation cellulaire \\\hline
\end{tabular}
\end{center}
\end{propriete}

\begin{savaistu}[Contexte camerounais : le lait et la lactase]
Au Cameroun, comme dans une grande partie de l'Afrique subsaharienne, la \textbf{persistance de la lactase} à l'âge adulte est fréquente chez les populations d'éleveurs (Peuls, Mbororos du Nord et de l'Adamaoua) mais plus rare chez les populations traditionnellement non-éleveuses (zones forestières du Sud et de l'Est). C'est pourquoi la consommation de lait frais peut provoquer ballonnements et diarrhées chez certains adultes : l'enzyme \textbf{lactase} n'est plus produite en quantité suffisante pour hydrolyser le lactose, qui fermente alors dans le côlon. Le lait fermenté (pendé, kindirmo) est mieux toléré car les bactéries lactiques ont déjà hydrolysé le lactose.
\end{savaistu}

\courssub{Localisation : enzymes intracellulaires et extracellulaires}

Les enzymes agissent là où le métabolisme en a besoin.

\begin{propriete}[Localisation et fonction]
\begin{itemize}
\item \textbf{Enzymes intracellulaires :} Agissent \emph{à l'intérieur} de la cellule qui les a synthétisées. Elles assurent le \textbf{métabolisme intermédiaire} (glycolyse, cycle de Krebs, synthèse des protéines, réplication ADN, détoxication). Exemples : hexokinase, citrate synthase, ADN polymérase, catalase.
\item \textbf{Enzymes extracellulaires :} Sécrétées \emph{hors} de la cellule (souvent sous forme inactive = \textbf{zymogène} ou proenzyme) pour agir dans un milieu externe (lumière digestive, sang, espace interstitiel). Elles assurent la \textbf{digestion} ou la \textbf{coagulation}. Exemples : amylase salivaire/pancréatique, trypsine (sécrétée sous forme \textbf{trypsinogène}), pepsine (sécrétée sous forme \textbf{pepsinogène}), thrombine.
\end{itemize}
\end{propriete}

\begin{attention}[Zymogène vs Enzyme active]
Ne pas confondre \textbf{zymogène} (précurseur inactif, ex. \emph{trypsinogène}, \emph{pepsinogène}) et \textbf{enzyme active} (ex. \emph{trypsine}, \emph{pepsine}). L'activation se fait souvent par clivage protéolytique (changement de pH, action d'une autre enzyme). C'est un mécanisme de protection : la cellule ne se digère pas elle-même.
\end{attention}

\courssub{Comparaison : Enzyme vs Catalyseur minéral}

Bien que tous deux soient des catalyseurs, les enzymes surpassent les catalyseurs minéraux sur plusieurs points clés pour la vie.

\begin{popfigure}
\begin{center}
\begin{tabularx}{\linewidth}{|l|X|X|}\hline
\rowcolor{popblueL}
\textbf{Critère} & \textbf{Catalyseur minéral (ex. Pt, $\ce{MnO2}$, $\ce{Fe^{3+}}$)} & \textbf{Enzyme (catalyseur biologique)} \\\hline
\textbf{Nature chimique} & Minéral (métal, oxyde, ion) & \textbf{Protéique} (chaîne polypeptidique ± cofacteur) \\\hline
\textbf{Spécificité} & Faible (un catalyseur pour plusieurs réactions) & \textbf{Absolue / très élevée} (une enzyme $\leftrightarrow$ un substrat ou groupe proche) \\\hline
\textbf{Efficacité (vitesse)} & Modérée (facteur $10^2$ à $10^4$) & \textbf{Extrême} (facteur $10^8$ à $10^{17}$) \\\hline
\textbf{Conditions optimales} & Souvent T° élevée, pH extrême, pression & \textbf{Conditions physiologiques} : 37\,°C, pH 6--8 (sauf exceptions) \\\hline
\textbf{Régulation} & Difficile (température, concentration) & \textbf{Très fine} : inhibition, activation, synthèse/dégradation, compartimentation \\\hline
\textbf{Sensibilité} & Stable à la chaleur, aux solvants & \textbf{Thermolabile, sensible au pH, aux métaux lourds, aux agents dénaturants} \\\hline
\textbf{Exemple classique} & $\ce{MnO2}$ pour $\ce{2H2O2 -> 2H2O + O2}$ & \textbf{Catalase} pour la même réaction (millions de fois plus rapide) \\\hline
\end{tabularx}
\end{center}
\legende{Tableau comparatif : l'enzyme est un catalyseur « sur mesure » pour la vie, opérant dans la douceur et la précision, au prix d'une grande fragilité.}
\end{popfigure}

\courssub{Exemple concret résolu : la catalase, championne de l'efficacité}

\begin{exoresolu}[La catalase : une enzyme « recordwoman »]
\textbf{Énoncé :} La catalase est une enzyme présente dans quasi tous les organismes aérobies (peroxysomes). Elle catalyse la décomposition de l'eau oxygénée ($\ce{H2O2}$), sous-produit toxique du métabolisme aérobie, en eau et dioxygène : $\ce{2H2O2 ->[catalase] 2H2O + O2}$. Son nombre de tour ($k_{\text{cat}}$) est d'environ $4 \times 10^7\,\mathrm{s^{-1}}$ à 25\,°C. Une cellule hépatique contient environ $10^{10}$ molécules de catalase. Calculer le nombre maximal de molécules de $\ce{H2O2}$ qu'une seule cellule peut dégrader par seconde. Comparer à une réaction non catalysée (constante de vitesse d'ordre 1 $k_{\text{non cat}} \approx 10^{-6}\,\mathrm{s^{-1}}$ pour $[\ce{H2O2}] = 1\,\mathrm{mM}$, volume cellulaire $\approx 2 \times 10^{-12}\,\mathrm{L}$).

\textbf{Solution détaillée :}
\begin{enumerate}
\item \textbf{Vitesse maximale par enzyme :} $v_{\text{max}} = k_{\text{cat}} = 4 \times 10^7$ molécules $\ce{H2O2}$ par seconde par molécule de catalase.
\item \textbf{Vitesse totale pour la cellule :} $V_{\text{cell}} = (4 \times 10^7) \times (10^{10}) = \mathbf{4 \times 10^{17}}$ molécules $\ce{H2O2}\,\mathrm{s^{-1}}$.
\item \textbf{En moles :} $\dfrac{4 \times 10^{17}}{6,02 \times 10^{23}} \approx \mathbf{6,6 \times 10^{-7}\,\mathrm{mol\,s^{-1}}}$ (soit $\approx 0,66\,\mathrm{\mu mol\,s^{-1}}$).
\item \textbf{Comparaison sans enzyme :} 
Nombre de molécules de substrat dans la cellule : $N_S = [\ce{H2O2}] \times V \times N_A = (10^{-3}\,\mathrm{mol\,L^{-1}}) \times (2 \times 10^{-12}\,\mathrm{L}) \times (6,02 \times 10^{23}\,\mathrm{mol^{-1}}) \approx 1,2 \times 10^9$ molécules.
Vitesse non catalysée (ordre 1) : $v_{\text{non cat}} = k_{\text{non cat}} \times N_S = 10^{-6}\,\mathrm{s^{-1}} \times 1,2 \times 10^9 \approx \mathbf{1,2 \times 10^3}$ molécules/s.
\item \textbf{Ratio d'accélération :} $\dfrac{4 \times 10^{17}}{1,2 \times 10^3} \approx \mathbf{3 \times 10^{14}}$. L'enzyme accélère la réaction d'un facteur $\sim 10^{14}$ dans ces conditions.
\end{enumerate}

\textbf{Conclusion :} Une seule molécule de catalase traite des \textbf{dizaines de millions de substrats par seconde}. Sans elle, la $\ce{H2O2}$ s'accumulerait et détruirait la cellule en quelques secondes. C'est la puissance de la catalyse enzymatique.
\end{exoresolu}

\courssub{Synthèse de la section}

\begin{aretenir}[Points clés à retenir]
\begin{itemize}
\item Une \textbf{enzyme} est une \textbf{protéine globulaire} (sauf ribozymes) à activité \textbf{catalytique}, produite par les cellules.
\item Elle accélère une réaction précise sans être consommée : $\ce{E + S <=> ES -> E + P}$.
\item \textbf{Substrat} = réactif spécifique ; \textbf{Produit} = molécule formée ; \textbf{Complexe ES} = intermédiaire transitoire.
\item Nomenclature : suffixe \textbf{-ase} (amylase, lactase, catalase, kinase...). Classification UICN en 6 classes.
\item \textbf{Intracellulaires} (métabolisme) vs \textbf{Extracellulaires} (digestion, souvent sécrétées comme zymogènes).
\item Comparées aux catalyseurs minéraux : \textbf{spécificité absolue, efficacité énorme ($10^8$--$10^{17}$), conditions douces (37\,°C, pH neutre), régulation fine, grande fragilité (thermolabilité)}.
\end{itemize}
\end{aretenir}

\begin{methode}[Comment analyser une expérience mettant en évidence une activité enzymatique]
\begin{enumerate}
\item \textbf{Identifier} les tubes : expérimental (enzyme active) et témoins (sans enzyme, enzyme dénaturée, substrat seul).
\item \textbf{Noter} les observations (coloration, dégagement gazeux, variation pH, etc.).
\item \textbf{Comparer} expérimental vs témoins : la différence prouve l'action de l'enzyme.
\item \textbf{Conclure} : l'enzyme catalyse la transformation du substrat en produit(s) ; elle est thermolabile (si témoin bouilli négatif) ; elle agit en faible quantité.
\item \textbf{Utiliser} le vocabulaire précis : substrat, produit, hydrolyse/synthèse, dénaturation, zymogène.
\end{enumerate}
\end{methode}

\begin{exercice}[Application : Le jus de goyave d'Aminata]
Reprenons l'observation d'Aminata : jus cru $\rightarrow$ brunissement rapide ; jus chauffé $\rightarrow$ pas de brunissement.
\begin{enumerate}
\item Quelle enzyme est probablement responsable du brunissement enzymatique des fruits ? (Indice : elle oxyde des phénols en quinones).
\item Pourquoi la chaleur empêche-t-elle le brunissement ?
\item Si l'on ajoute du jus de citron (acide) sur le fruit coupé, le brunissement ralentit. Quel facteur enzymatique cela illustre-t-il ?
\end{enumerate}
\end{exercice}

\begin{corrige}[Exercice : Le jus de goyave d'Aminata]
\begin{enumerate}
\item \textbf{Polyphénol oxydase (PPO)} aussi appelée \textbf{tyrosinase} ou \textbf{catécholase}. Elle catalyse l'oxydation des polyphénols en quinones (brunes) en présence d'$\ce{O2}$.
\item La chaleur \textbf{dénature} la PPO (protéine) : perte de structure tertiaire/quaternaire $\rightarrow$ site actif détruit $\rightarrow$ activité nulle. L'enzyme est « tuée » (inactivée irrémédiablement).
\item L'acidité (pH bas) du citron \textbf{inhibe} la PPO (pH optimum souvent neutre/alcalin). Cela illustre la \textbf{sensibilité au pH} de l'activité enzymatique (voir section 4).
\end{enumerate}
\end{corrige}

\coursec{Spécificité enzymatique et site actif}

\courssub{Transition depuis la section précédente}

Dans la section précédente, nous avons découvert que les enzymes sont des catalyseurs biologiques protéiques qui accélèrent les réactions chimiques du métabolisme sans être consommées. Nous avons vu qu'elles abaissent l'énergie d'activation et qu'elles interviennent aussi bien dans les réactions de dégradation (catabolisme) que de synthèse (anabolisme). Mais une question fondamentale demeure : \textbf{comment une enzyme « choisit-elle » sa réaction parmi les milliers qui ont lieu dans la cellule ?} Pourquoi l'amylase hydrolyse-t-elle l'amidon mais ignore-t-elle la cellulose, alors que ces deux polymères sont tous deux constitués de glucose ? C'est la notion de \cle{spécificité enzymatique} que nous allons explorer maintenant.

\courssub{Observation initiale : l'amylase, l'amidon et la cellulose}

\begin{experience}[Spécificité de l'amylase salivaire (test à l'iode)]
\textbf{Matériel :} 4 tubes à essai, solution d'amylase salivaire (salive diluée 1/10), solution d'amidon (1 \%), solution de cellulose (1 \%), solution d'iode (révélateur d'amidon : complexe amidon-iode bleu-noir), bain-marie à 37 °C.

\textbf{Protocole :}
\begin{enumerate}
    \item Tube 1 : 2 mL d'amidon + 1 mL d'amylase.
    \item Tube 2 : 2 mL de cellulose + 1 mL d'amylase.
    \item Tube 3 : 2 mL d'amidon + 1 mL d'eau (témoin sans enzyme).
    \item Tube 4 : 2 mL de cellulose + 1 mL d'eau (témoin sans enzyme).
    \item Incuber 10 min à 37 °C.
    \item Ajouter 3 gouttes de solution d'iode dans chaque tube. Observer la couleur.
\end{enumerate}

\textbf{Observations attendues :}
\begin{center}
\begin{tabularx}{\linewidth}{|l|X|X|}\hline
Tube & Contenu & Couleur après ajout d'iode \\ \hline
1 & Amidon + Amylase & \textbf{Jaune-orangé} (absence d'amidon : hydrolyse complète) \\
2 & Cellulose + Amylase & \textbf{Incolore} (cellulose non hydrolysée \emph{et} non révélée par l'iode) \\
3 & Amidon + Eau & \textbf{Bleu-noir} (amidon intact, pas d'hydrolyse spontanée) \\
4 & Cellulose + Eau & \textbf{Incolore} (cellulose non révélée par l'iode) \\ \hline
\end{tabularx}
\end{center}

\textbf{Interprétation :}
\begin{itemize}
    \item Le tube 1 montre la disparition de la coloration bleu-noir : l'amidon a été hydrolysé en maltose/glucose (non révélés par l'iode) par l'amylase.
    \item Le tube 2 reste incolore : la cellulose n'a \emph{pas} été hydrolysée (sinon on verrait du glucose, non coloré par l'iode) ; de plus, l'iode ne colore pas la cellulose.
    \item Le tube 3 (bleu-noir) confirme que l'hydrolyse ne se fait pas sans enzyme.
    \item Le tube 4 (incolore) confirme que l'iode ne réagit pas avec la cellulose.
\end{itemize}

\textbf{Conclusion :} L'amylase agit \emph{uniquement} sur l'amidon, pas sur la cellulose. Pourtant, chimiquement, l'amidon (liaisons $\alpha$-1,4) et la cellulose (liaisons $\beta$-1,4) sont très proches : tous deux sont des polymères de glucose. La différence réside dans la \textbf{configuration spatiale} des liaisons osidiques. L'enzyme « reconnaît » cette différence stéréochimique.
\end{experience}

\begin{aretenir}[Leçon de l'expérience]
Même substrats chimiquement voisins (polymères de glucose) $\rightarrow$ une seule enzyme n'agit que sur l'un d'entre eux. La \cle{spécificité} enzymatique repose sur une reconnaissance \textbf{stéréochimique} (forme tridimensionnelle), pas seulement sur la nature chimique des atomes.
\end{aretenir}

\courssub{Le site actif : la « serrure » de l'enzyme}

\begin{definition}[Site actif]
Le \textbf{site actif} (ou centre actif) est la région précise de l'enzyme, généralement une cavité ou une fente à sa surface, dont la \textbf{forme tridimensionnelle} et la \textbf{nature chimique} (acides aminés constitutifs) sont complémentaires à celles du \textbf{substrat}. C'est là que se fixe le substrat et que s'effectue la catalyse.
\end{definition}

\begin{itemize}
    \item Le site actif est formé par des acides aminés qui peuvent être éloignés dans la séquence primaire mais rapprochés par le repliement tertiaire (et quaternaire) de la protéine.
    \item Il ne représente qu'une \textbf{petite fraction} de la surface totale de l'enzyme (souvent $< 5\%$).
    \item Les acides aminés du site actif portent des groupements fonctionnels (–OH, –SH, –COOH, –NH$_2$, imidazole, etc.) qui participent à la liaison du substrat (liaisons hydrogène, interactions ioniques, forces de van der Waals, interactions hydrophobes) et à la réaction chimique elle-même (catalyse acide-base, catalyse covalente, catalyse par approximation).
\end{itemize}

\begin{popfigure}
\begin{tikzpicture}[scale=0.9, every node/.style={font=\small}]
% Enzyme shape
\draw[popblue, line width=2pt, fill=popblue!10] 
    (0,0) .. controls (0,2.5) and (3,3) .. (5,2.5) .. controls (7,2) and (7,0.5) .. (5,0) .. controls (3,-0.5) and (0,-0.5) .. cycle;

% Active site: a cleft
\draw[poporange, line width=2pt, fill=poporange!30] 
    (2.5,1.2) .. controls (2.5,1.8) and (3.5,2) .. (4.2,1.8) .. controls (4.5,1.5) and (4.5,0.8) .. (3.8,0.6) .. controls (3.2,0.5) and (2.5,0.7) .. cycle;

% Label enzyme
\node[popblue, align=center] at (-1,1.5) {Enzyme\\(protéine repliée)};
\draw[popblue, ->, line width=1pt] (-0.3,1.3) -- (0.5,1.1);

% Label active site
\node[poporange, align=center] at (5.5,2.2) {Site actif\\(cavité complémentaire)};
\draw[poporange, ->, line width=1pt] (5.2,2) -- (4.3,1.8);

% Substrate 1: complementary (key)
\begin{scope}[shift={(8,1.5)}]
\draw[popgreen, line width=2pt, fill=popgreen!30] 
    (0,0) -- (0.8,0) .. controls (1.1,0) and (1.1,0.6) .. (0.8,0.6) -- (0,0.6) .. controls (-0.3,0.6) and (-0.3,0) .. cycle;
\node[popgreen, align=center] at (0.4,-0.5) {Substrat\\complémentaire};
\end{scope}

% Arrow binding
\draw[popdark, ->, thick] (5.5,1.2) -- (7.2,1.5);
\node[popdark, align=center] at (6.3,1.8) {Fixation};

% Substrate 2: non-complementary (rejected)
\begin{scope}[shift={(8,-1.5)}]
\draw[poppink, line width=2pt, fill=poppink!30] 
    (0,0) .. controls (0,0.6) and (0.8,0.6) .. (0.8,0) .. controls (0.8,-0.6) and (0,-0.6) .. cycle;
\node[poppink, align=center] at (0.4,-1) {Substrat\\non complémentaire};
\end{scope}

% Arrow rejection
\draw[poppink, ->, thick, dashed] (5.5,0.5) -- (7.2,-1);
\node[poppink, align=center] at (6.3,-0.5) {Pas de fixation};

% Lock and key labels
\node[popdark, align=center, font=\bfseries] at (3.5,3.5) {Modèle « Clé-Serrure » (Fischer, 1894)};
\node[popdark, align=center] at (3.5,-2.5) {Seule la « bonne clé » (substrat) entre dans la « serrure » (site actif).};
\end{tikzpicture}
\legende{Modèle clé-serrure : le site actif de l'enzyme a une forme rigide complémentaire à celle du substrat spécifique. Un substrat de forme différente ne peut s'y fixer.}
\end{popfigure}

\begin{propriete}[Spécificité et complémentarité de forme (Modèle clé-serrure)]
La \cle{spécificité enzymatique} résulte de la \textbf{complémentarité stéréochimique} entre le site actif de l'enzyme et le substrat. Selon le modèle de \textbf{Fischer (1894)}, le site actif possède une conformation rigide, préformée, qui ne convient qu'à un seul substrat (ou à un groupe de substrats très voisins), comme une serrure ne s'ouvre qu'avec sa clé. Ce modèle explique pourquoi l'amylase hydrolyse l'amidon (liaisons $\alpha$) mais pas la cellulose (liaisons $\beta$) : la géométrie du site actif ne permet l'ajustement qu'avec la conformation $\alpha$.
\end{propriete}

\courssub{Complément pour le Probatoire : le modèle de l'ajustement induit (Koshland, 1958)}

Le modèle clé-serrure, bien que fondamental, est incomplet. Des études cinétiques et structurales (cristallographie aux rayons X) ont montré que le site actif n'est pas parfaitement rigide.

\begin{propriete}[Modèle de l'ajustement induit (Koshland)]
Lors de la fixation du substrat, le site actif subit une \textbf{déformation conformationnelle} (changement de forme) qui lui permet d'épouser parfaitement le substrat. Cette flexibilité améliore la catalyse en :
\begin{itemize}
    \item rapprochant les groupes catalytiques des liaisons à rompre/former,
    \item excluant l'eau du site actif (milieu hydrophobe favorable),
    \item déstabilisant le substrat vers l'état de transition.
\end{itemize}
L'enzyme est donc une protéine \textbf{dynamique}, pas une statue rigide.
\end{propriete}

\begin{attention}[Piège fréquent au Probatoire]
Ne pas opposer les deux modèles : le modèle \textbf{clé-serrure} explique la \textbf{spécificité} (reconnaissance initiale), le modèle \textbf{d'ajustement induit} explique l'\textbf{efficacité catalytique} (adaptation pour la réaction). Les deux sont compatibles et souvent enseignés ensemble. Au Probatoire, citez \textbf{Fischer} pour la spécificité, \textbf{Koshland} pour le mécanisme catalytique.
\end{attention}

\courssub{Formation du complexe enzyme-substrat (ES) : étape intermédiaire obligatoire}

La catalyse enzymatique suit un mécanisme en \textbf{trois étapes} successives, matérialisé par la formation d'un \textbf{complexe enzyme-substrat (ES)} transitoire.

\begin{popfigure}
\begin{tikzpicture}[scale=0.9, every node/.style={font=\small}]
% Styles
\tikzset{
    enzyme/.style={draw=popblue, line width=2pt, fill=popblue!10, rounded corners=5pt, minimum width=2.5cm, minimum height=1.5cm},
    substrate/.style={draw=popgreen, line width=2pt, fill=popgreen!30, rounded corners=3pt, minimum width=0.8cm, minimum height=0.6cm},
    product/.style={draw=poporange, line width=2pt, fill=poporange!30, rounded corners=3pt, minimum width=0.6cm, minimum height=0.5cm},
    arrow/.style={->, thick, popdark},
    label/.style={popdark, align=center, font=\small},
    cataly/.style={popdark, decorate, decoration={snake, amplitude=1pt, segment length=3pt}}
}

% Step 1: Free enzyme + free substrate
\begin{scope}
\node[enzyme] (E1) at (1.25,0.75) {};
\node[label] at (1.25,2.3) {Étape 1 : Fixation};
\node[label] at (1.25,-0.8) {Enzyme libre + Substrat libre};
\node[popblue] at (1.25,1.6) {E};
% Active site cleft
\draw[poporange, fill=poporange!30] (0.8,0.7) .. controls (0.8,1.1) and (1.2,1.2) .. (1.5,1.1) .. controls (1.7,1) and (1.7,0.5) .. (1.3,0.4) .. controls (1.1,0.3) and (0.8,0.4) .. cycle;
% Substrate
\node[substrate] (S1) at (3.8,0.75) {};
\node[popgreen] at (3.8,1.2) {S};
\draw[arrow] (2.5,0.75) -- (3.4,0.75);
\end{scope}

% Arrow between step 1 and 2
\draw[arrow] (4.6,0.75) -- (5.7,0.75);
\node[label] at (5.15,1.3) {Formation du complexe ES};

% Step 2: ES Complex
\begin{scope}[shift={(6.5,0)}]
\node[enzyme] (E2) at (1.25,0.75) {};
\node[label] at (1.25,2.3) {Étape 2 : Catalyse};
\node[label] at (1.25,-0.8) {Complexe Enzyme-Substrat (ES)};
\node[popblue] at (1.25,1.6) {E};
% Substrate inside active site (induced fit)
\node[substrate, minimum width=0.7cm, minimum height=0.5cm] (ES) at (1.25,0.75) {};
\node[popgreen] at (1.25,1.15) {S};
% Catalytic groups
\draw[cataly] (1.0,0.9) -- (1.1,1.05);
\draw[cataly] (1.5,0.9) -- (1.4,1.05);
\node[label, font=\tiny, align=center] at (0.6,1.1){Groupes\\catalytiques};
\node[label, font=\tiny, align=center] at (1.9,1.1){Groupes\\catalytiques};
% Transition state symbol
\node[label, font=\bfseries, text=poppink] at (1.25,0.25) {État de transition $ES^\ddagger$};
\end{scope}

% Arrow between step 2 and 3
\draw[arrow] (10.7,0.75) -- (11.8,0.75);
\node[label] at (11.25,1.3) {Libération};

% Step 3: Free enzyme + products
\begin{scope}[shift={(12.5,0)}]
\node[enzyme] (E3) at (1.25,0.75) {};
\node[label] at (1.25,2.3) {Étape 3 : Libération};
\node[label] at (1.25,-0.8) {Enzyme libre + Produits};
\node[popblue] at (1.25,1.6) {E};
% Products
\node[product] (P1) at (3.2,1.1) {};
\node[poporange] at (3.2,1.45) {P$_1$};
\node[product] (P2) at (3.8,0.4) {};
\node[poporange] at (3.8,0.05) {P$_2$};
\draw[arrow] (2.5,1.1) -- (2.9,1.1);
\draw[arrow] (2.5,0.4) -- (3.5,0.4);
\end{scope}

% Overall equation
\node[label, font=\bfseries] at (7,-2.7) {\ce{E + S <=>[k_1][k_{-1}] ES ->[k_2] E + P}};
\end{tikzpicture}
\legende{Les trois étapes de la catalyse enzymatique : (1) fixation du substrat sur le site actif, (2) formation du complexe ES et transformation en état de transition, (3) libération des produits et régénération de l'enzyme libre. L'enzyme n'est pas modifiée au bilan global.}
\end{popfigure}

\begin{aretenir}[Mécanisme général]
\[\ce{E + S <=>[k_1][k_{-1}] ES ->[k_2] E + P}\]
\begin{itemize}
    \item $k_1$ : constante de formation du complexe ES.
    \item $k_{-1}$ : constante de dissociation du complexe ES (sans réaction).
    \item $k_2$ : constante de transformation du complexe ES en produits (étape limitante souvent).
    \item L'enzyme E est \textbf{régénérée} à la fin : elle peut catalyser un nouveau cycle.
\end{itemize}
\end{aretenir}

\courssub{Spécificité d'action : une enzyme = une réaction (ou un type de réaction)}

La spécificité ne concerne pas seulement la \textbf{reconnaissance} du substrat (spécificité de liaison), mais aussi la \textbf{réaction catalysée} (spécificité d'action).

\begin{definition}[Spécificité d'action]
Une enzyme donnée ne catalyse qu'une \textbf{seule réaction chimique} (ou un seul type de réaction chimique) sur son substrat spécifique. Elle oriente le substrat vers un unique produit (ou couple de produits) parmi les transformations théoriquement possibles.
\end{definition}

\begin{exemplebox}[La lactase : un exemple classique de spécificité absolue]
La \textbf{lactase} (β-galactosidase) est une enzyme qui hydrolyse \textbf{uniquement} le \textbf{lactose} (sucre du lait) en \textbf{glucose} + \textbf{galactose}.
\begin{itemize}
    \item Substrat : lactose (disaccharide : galactose-β-1,4-glucose).
    \item Produits : glucose + galactose.
    \item Elle n'agit \emph{pas} sur le saccharose (glucose-α-1,2-fructose), ni sur le maltose (glucose-α-1,4-glucose), ni sur la cellobiose (glucose-β-1,4-glucose).
\end{itemize}
\textbf{Conséquence physiologique :} Chez l'adulte, la diminution de l'activité lactase (non-persistance de la lactase) entraîne une \textbf{intolérance au lactose} : le lactose non hydrolysé arrive dans le côlon, où il est fermenté par la flore bactérienne $\rightarrow$ gaz, ballonnements, diarrhée. Au Cameroun, comme dans une grande partie de l'Afrique subsaharienne, la non-persistance de la lactase est fréquente chez l'adulte, ce qui explique pourquoi le lait frais est mal toléré, alors que le yaourt (lactose pré-hydrolysé par les ferments lactiques) ou le fromage (faible en lactose) le sont mieux.
\end{exemplebox}

\begin{exemplebox}[Spécificité relative : les protéases trypsine et chymotrypsine]
Certaines enzymes ont une \textbf{spécificité relative} (ou de groupe) : elles reconnaissent un type de liaison chimique dans un contexte précis.
\begin{itemize}
    \item \textbf{Trypsine} : coupe les liaisons peptidiques \emph{après} une lysine (Lys) ou une arginine (Arg) — acides aminés basiques.
    \item \textbf{Chymotrypsine} : coupe \emph{après} un acide aminé aromatique (Phe, Tyr, Trp) ou à chaîne latérale volumineuse (Leu, Met).
\end{itemize}
Elles sont toutes deux des protéases (spécificité de type : hydrolyse de liaison peptidique) mais leur spécificité de substrat diffère selon l'acide aminé en position P1. Cette complémentarité est exploitée en laboratoire pour le séquençage des protéines et en biotechnologie (production d'hydrolysats de protéines pour l'alimentation).
\end{exemplebox}

\courssub{Conséquence majeure : la diversité enzymatique reflète la diversité métabolique}

\begin{propriete}[Corrélation diversité enzymatique / diversité métabolique]
Puisque chaque réaction métabolique (ou type de réaction) nécessite une enzyme spécifique, le \textbf{nombre d'enzymes différentes} présentes dans une cellule est au moins égal au \textbf{nombre de réactions métaboliques distinctes} qu'elle effectue. Une cellule de levure possède $\sim$ 6 000 protéines codantes, dont $\sim$ 1 000 enzymes ; une cellule humaine en exprime $\sim$ 20 000, dont plusieurs milliers d'enzymes. Cette diversité est encodée dans le génome : \textbf{un gène $\rightarrow$ une enzyme (hypothèse de Beadle et Tatum, 1941, prix Nobel 1958)}.
\end{propriete}

\begin{savaistu}[Contexte camerounais : enzymes et agriculture]
Au Cameroun, la compréhension de la spécificité enzymatique a des applications concrètes :
\begin{itemize}
    \item \textbf{Amélioration de la farine de manioc :} l'ajout d'enzymes spécifiques (amylases, cellulases, pectinases) lors de la transformation du manioc en gari ou en fufu permet de réduire la teneur en cyanure (hydrolyse des glycosides cyanogènes par la linamarase endogène, activée par le broyage) et d'améliorer la texture.
    \item \textbf{Production de jus de fruits :} les pectinases (spécifiques de la pectine) et cellulases clarifient les jus d'ananas, de mangue, de goyave en hydrolysant les parois cellulaires.
    \item \textbf{Alimentation animale :} l'ajout de phytase (hydrolyse l'acide phytique, libérant le phosphore) dans l'alimentation des volailles et porcs réduit le besoin en phosphate minéral importé et diminue la pollution phosphatée des rejets.
\end{itemize}
Dans chaque cas, le choix de l'enzyme repose sur sa \textbf{spécificité} connue pour le substrat cible.
\end{savaistu}

\courssub{Méthode : interpréter une expérience de spécificité enzymatique}

\begin{methode}[Interpréter une expérience de spécificité]
Pour conclure à la spécificité d'une enzyme à partir de résultats expérimentaux (dosage de substrat résiduel ou de produit formé), suivez ces étapes :
\begin{enumerate}
    \item \textbf{Identifier les variables :} enzyme testée, substrats testés (au moins deux : le substrat supposé spécifique et un ou plusieurs substrats « témoins » chimiquement voisins), conditions (T°, pH, temps) identiques pour tous les tubes.
    \item \textbf{Comparer les tubes \emph{avec enzyme} :} seul le tube contenant le « bon » substrat montre une disparition du substrat ou une apparition du produit (dosage colorimétrique, polarimétrie, chromatographie, etc.).
    \item \textbf{Vérifier les témoins \emph{sans enzyme} :} ils confirment que la réaction ne se fait pas spontanément dans les conditions du test.
    \item \textbf{Vérifier le témoin \emph{enzyme dénaturée} (ex. ébouillantée) + substrat spécifique :} confirme que l'activité protéique est nécessaire.
    \item \textbf{Conclure :} « L'enzyme X hydrolyse / synthétise uniquement le substrat Y (ou le type de liaison Z) dans ces conditions. Elle est spécifique de Y. »
\end{enumerate}
\end{methode}

\courssub{Exercice résolu : analyse d'une expérience de spécificité}

\begin{exoresolu}[Spécificité de la maltase]
\textbf{Énoncé :} On étudie la spécificité de la maltase (enzyme intestinale). On prépare 5 tubes contenant chacun 2 mL d'une solution de substrat à 1 \% et 1 mL de solution de maltase purifiée. Incubation 30 min à 37 °C. On dose ensuite le glucose formé (réaction de Fehling). Résultats :

\begin{center}
\begin{tabular}{|c|c|c|}\hline
Tube & Substrat testé & Résultat Fehling (glucose) \\ \hline
1 & Maltose & +++ (fort dépôt rouge brique) \\
2 & Saccharose & - (bleu, pas de dépôt) \\
3 & Lactose & - (bleu, pas de dépôt) \\
4 & Amidon & - (bleu, pas de dépôt) \\
5 & Maltose + maltase ébouillantée & - (bleu, pas de dépôt) \\ \hline
\end{tabular}
\end{center}

\textbf{Questions :}
\begin{enumerate}
    \item Interpréter les résultats des tubes 1 à 4.
    \item Quel est le rôle du tube 5 ?
    \item Conclure sur la spécificité de la maltase.
\end{enumerate}

\tcblower

\textbf{Solution détaillée :}

\begin{enumerate}
    \item \textbf{Tubes 1 à 4 (comparaison substrats) :} Seul le tube 1 (maltose) donne un test de Fehling positif (+++), indiquant la formation de glucose (osides réducteurs). Les tubes 2 (saccharose), 3 (lactose) et 4 (amidon) restent négatifs : aucun glucose n'est détecté. \emph{Interprétation :} La maltase hydrolyse le maltose en deux molécules de glucose, mais n'hydrolyse ni le saccharose, ni le lactose, ni l'amidon.
    \item \textbf{Tube 5 (témoin enzyme dénaturée) :} La maltase a été ébouillantée $\rightarrow$ dénaturation thermique (perte de la structure tertiaire, destruction du site actif). Le test est négatif. \emph{Rôle :} Confirmer que l'hydrolyse observée au tube 1 est bien due à l'activité \emph{catalytique} de la protéine maltase native, et non à une hydrolyse acide spontanée ou à une contamination.
    \item \textbf{Conclusion :} La maltase est une enzyme \textbf{spécifique du maltose} (disaccharide $\alpha$-1,4). Elle ne reconnaît pas les liaisons $\alpha$-1,2 (saccharose), $\beta$-1,4 (lactose), ni les chaînes polymériques $\alpha$-1,4 (amidon). Sa spécificité porte sur la nature du substrat (maltose) et le type de liaison osidique ($\alpha$-1,4 terminale).
\end{enumerate}
\end{exoresolu}

\courssub{Synthèse de la section}

\begin{aretenir}[Points clés à retenir]
\begin{itemize}
    \item La \cle{spécificité enzymatique} est la propriété qu'a une enzyme de ne catalyser qu'une seule réaction (ou un type de réaction) sur un substrat précis (ou un groupe de substrats très voisins).
    \item Elle repose sur la \textbf{complémentarité de forme} entre le \textbf{site actif} (cavité de l'enzyme) et le substrat : modèle \textbf{clé-serrure} (Fischer) pour la reconnaissance, modèle \textbf{d'ajustement induit} (Koshland) pour l'optimisation catalytique.
    \item La catalyse passe obligatoirement par la formation d'un \textbf{complexe enzyme-substrat (ES)} transitoire : $\ce{E + S <=> ES -> E + P}$.
    \item La diversité des réactions métaboliques impose une diversité correspondante d'enzymes (hypothèse « un gène – une enzyme »).
    \item Pour prouver la spécificité expérimentalement : tester plusieurs substrats structuralement voisins dans les mêmes conditions, avec témoins sans enzyme et enzyme dénaturée.
\end{itemize}
\end{aretenir}

\begin{attention}[Erreurs à ne pas commettre]
\begin{itemize}
    \item \textbf{Confondre spécificité et affinité :} une enzyme peut avoir une forte affinité (fixer fort) pour un substrat sans le transformer efficacement (inhibiteur compétitif). La spécificité concerne le \textbf{turnover catalytique} ($k_{\text{cat}}$), pas seulement la liaison ($K_M$).
    \item \textbf{Dire « l'enzyme change de forme pour accepter n'importe quel substrat » :} l'ajustement induit affine l'ajustement pour \emph{le bon} substrat ; il ne permet pas d'accepter n'importe quoi.
    \item \textbf{Oublier les témoins} dans l'interprétation d'expérience : sans témoin sans enzyme et sans témoin enzyme dénaturée, on ne peut pas attribuer l'effet à l'activité enzymatique spécifique.
    \item \textbf{Écrire « l'enzyme est spécifique \emph{au} substrat » :} on dit « spécifique \emph{de} » ou « spécifique \emph{pour} » un substrat.
    \item \textbf{Confondre « un gène – une enzyme » et la réalité moderne :} l'hypothèse historique (Beadle \& Tatum) a été affinée en « un gène – un polypeptide » (certains gènes codent pour des sous-unités, d'autres pour des ARN non codants, le splicing alternatif existe). Au Probatoire, on cite l'hypothèse historique comme base conceptuelle.
\end{itemize}
\end{attention}

\coursec{Influence de la température sur l'activité enzymatique}

Dans la section précédente, nous avons découvert que chaque enzyme possède un \cle{site actif} dont la forme est complémentaire de celle de son \cle{substrat} : c'est ce qui explique la \cle{spécificité enzymatique}. Mais cette belle architecture tridimensionnelle, si précise, est aussi fragile. Pourquoi conserve-t-on les aliments au réfrigérateur ? Pourquoi une fièvre de 42 °C met-elle la vie en danger ? Pourquoi la cuisson tue-t-elle les microbes ? Toutes ces questions de la vie quotidienne ont la même réponse : la \cle{température} modifie profondément l'activité des enzymes. C'est ce que nous allons démontrer expérimentalement dans cette section.

\courssub{Situation-problème : pourquoi le réfrigérateur conserve-t-il les aliments ?}

\begin{experience}[Observation de départ]
Un morceau de viande laissé à l'air libre à Ngaoundéré en saison chaude (environ 35 °C) se décompose en quelques heures : il change de couleur et dégage une odeur nauséabonde. Le même morceau placé dans un réfrigérateur à 4 °C se conserve plusieurs jours. Congelé à $-18$ °C, il se conserve plusieurs mois. Or la décomposition de la viande est due à des réactions chimiques de dégradation catalysées par des \cle{enzymes} (celles des microbes et celles des tissus eux-mêmes).
\end{experience}

\textbf{Question scientifique.} Si le froid empêche la décomposition, c'est peut-être que la température agit directement sur l'activité des enzymes. Formulons une hypothèse testable : \emph{« l'activité d'une enzyme dépend de la température »}. Pour la vérifier, il faut une expérience contrôlée, où l'on fait varier uniquement la température et où l'on mesure l'activité enzymatique.

\courssub{Expérience de référence : l'activité de l'amylase à différentes températures}

\begin{experience}[Activité de l'amylase en fonction de la température]
\textbf{Matériel :} solution d'amidon (substrat), solution de salive diluée (source d'\cle{amylase}, enzyme qui hydrolyse l'amidon en sucres simples), tubes à essais, bains-marie réglés à différentes températures, eau iodée (réactif qui colore l'amidon en bleu-noir), glace.

\textbf{Protocole :}
\begin{enumerate}
\item On prépare 5 tubes identiques contenant chacun le même volume de solution d'amidon et le même volume de salive (même concentration en enzyme).
\item Chaque tube est placé dans un bain à une température différente : 0 °C (glace), 20 °C, 37 °C, 60 °C et 100 °C.
\item Après 10 minutes, on prélève une goutte de chaque tube et on ajoute de l'eau iodée.
\item Plus la coloration bleu-noir est faible, plus l'amidon a été dégradé, donc plus l'amylase a été active.
\end{enumerate}

\textbf{Observations :}
\begin{center}
\begin{tabularx}{\linewidth}{|l|X|X|}
\hline
\rowcolor{popblueL} \textbf{Température} & \textbf{Coloration à l'eau iodée} & \textbf{Interprétation} \\
\hline
0 °C & Bleu-noir intense & Amidon presque intact : activité très faible \\
\hline
20 °C & Bleu-noir moyen & Dégradation partielle : activité moyenne \\
\hline
37 °C & Jaune (pas de bleu) & Amidon totalement dégradé : activité maximale \\
\hline
60 °C & Bleu-noir marqué & Dégradation faible : activité très réduite \\
\hline
100 °C & Bleu-noir intense & Amidon intact : activité nulle \\
\hline
\end{tabularx}
\end{center}

\textbf{Conclusion :} l'activité de l'amylase augmente avec la température jusqu'à 37 °C environ, puis diminue brutalement et s'annule à haute température. L'hypothèse est validée : \emph{la température influence l'activité enzymatique}.
\end{experience}

\courssub{La courbe d'activité : une montée puis une chute brutale}

En mesurant quantitativement la vitesse de la réaction (par exemple la quantité d'amidon disparue par minute), on obtient une courbe en cloche caractéristique.

\begin{popfigure}
\begin{center}
\begin{tikzpicture}
\begin{axis}[
    width=0.9\linewidth,
    height=7cm,
    xlabel={Température (\si{\celsius})},
    ylabel={Activité enzymatique (\%)},
    xmin=0, xmax=100,
    ymin=0, ymax=115,
    xtick={0,20,37,60,80,100},
    ytick={0,20,40,60,80,100},
    grid=both,
    grid style={popline!40},
    axis lines=left,
    legend style={at={(0.03,0.97)},anchor=north west,draw=none,fill=none}
]
% Zone de dénaturation hachurée
\fill[pattern=north east lines, pattern color=poporange!70, opacity=0.5]
    (axis cs:50,0) rectangle (axis cs:100,115);
% Courbe d'activité : montée gaussienne, chute exponentielle asymétrique
\addplot[popcurve, domain=0:37, samples=200] {100*exp(-((x-37)^2)/(2*12^2))};
\addplot[popcurve, domain=37:100, samples=200] {100*exp(-(x-37)/8)};
\addlegendentry{Amylase humaine}
% Optimum
\draw[dashed, popgreen, thick] (axis cs:37,0) -- (axis cs:37,100);
\draw[dashed, popgreen, thick] (axis cs:0,100) -- (axis cs:37,100);
\fill[popgreen] (axis cs:37,100) circle (2.5pt);
\node[popgreen, anchor=south, font=\small\bfseries] at (axis cs:37,103) {Optimum : 37 °C};
% Annotations
\node[popblue, font=\small, align=center] at (axis cs:15,55) {Branche\\ascendante\\(activation)};
\node[poporange, font=\small\bfseries, align=center] at (axis cs:78,60) {Zone de\\dénaturation\\(irréversible)};
\draw[->, poporange, thick] (axis cs:70,50) -- (axis cs:58,30);
\end{axis}
\end{tikzpicture}
\end{center}
\legende{Courbe de l'activité de l'amylase humaine en fonction de la température. L'activité croît jusqu'à un optimum (37 °C), puis chute brutalement dans la zone de dénaturation (hachurée).}
\end{popfigure}

\begin{definition}[Température optimale]
La \cle{température optimale} d'une enzyme est la température à laquelle son activité catalytique est \textbf{maximale}. Pour la plupart des enzymes de l'organisme humain, elle est voisine de \textbf{37 °C}, c'est-à-dire la température normale du corps. Ce n'est pas un hasard : les enzymes sont adaptées aux conditions de leur milieu de vie.
\end{definition}

\courssub{Interprétation de la branche ascendante : l'agitation moléculaire}

Pourquoi l'activité augmente-t-elle quand la température s'élève, entre 0 °C et 37 °C ?

La température mesure le degré d'\cle{agitation moléculaire} : plus il fait chaud, plus les molécules se déplacent rapidement et de façon désordonnée. Or, pour qu'une réaction enzymatique ait lieu, il faut que le substrat \emph{rencontre} le site actif de l'enzyme.

\begin{itemize}
\item Quand la température augmente, les molécules d'enzyme et de substrat se déplacent plus vite ;
\item le nombre de \cle{chocs} (collisions) entre l'enzyme et son substrat augmente ;
\item la formation du complexe enzyme-substrat est donc plus fréquente, et la vitesse de la réaction augmente.
\end{itemize}

\begin{attention}[Le froid ne « tue » pas les enzymes !]
\textbf{Erreur fréquente au Probatoire :} écrire que « le froid détruit les enzymes » ou que « le froid tue les microbes ». C'est faux ! À basse température (0 °C, 4 °C), l'enzyme est simplement \textbf{ralentie} : son activité est très faible, mais sa structure spatiale est \emph{intacte}. Si l'on réchauffe le milieu, l'enzyme retrouve pleinement son activité. L'effet du froid est donc \textbf{réversible}. C'est exactement le principe du réfrigérateur : les enzymes des microbes continuent d'exister, mais elles travaillent au ralenti, ce qui retarde la décomposition des aliments.
\end{attention}

\courssub{Interprétation de la chute brutale : la dénaturation thermique}

Au-delà de la température optimale, l'activité s'effondre très rapidement. Que se passe-t-il ?

Rappelons qu'une enzyme est une \cle{protéine} : une longue chaîne d'acides aminés repliée sur elle-même en une forme tridimensionnelle précise, maintenue par des liaisons chimiques fragiles (liaisons hydrogène, interactions hydrophobes, ponts disulfure). C'est cette forme spatiale qui crée le site actif.

Lorsque la température devient trop élevée, l'agitation thermique intense \textbf{rompt ces liaisons fragiles} : la protéine se déplie, perd sa forme tridimensionnelle, et le site actif est \textbf{déformé}. Le substrat ne peut plus s'y fixer : l'enzyme est inactive. Ce phénomène s'appelle la \cle{dénaturation}.

\begin{propriete}[Irréversibilité de la dénaturation thermique (admise)]
La \cle{dénaturation thermique} d'une enzyme est \textbf{irréversible} : elle détruit définitivement la structure tridimensionnelle de la protéine. Une enzyme dénaturée par la chaleur ne retrouve \emph{jamais} son activité, même si l'on ramène la température à 37 °C. Cette propriété se démontre expérimentalement : si l'on chauffe de l'amylase à 100 °C puis qu'on la refroidit à 37 °C, elle reste incapable de dégrader l'amidon.
\end{propriete}

\begin{popfigure}
\begin{center}
\begin{tikzpicture}[scale=1, every node/.style={font=\small}]
% Enzyme native
\begin{scope}
\draw[thick, popblue, fill=popblueL, rounded corners=6pt]
    (0,0) .. controls (0.3,1.2) and (1.2,1.6) .. (2.2,1.3)
    .. controls (2.6,1.2) and (2.7,0.9) .. (2.5,0.6)
    .. controls (2.3,0.3) and (2.3,0.1) .. (2.5,-0.2)
    .. controls (2.8,-0.6) and (2.5,-1.1) .. (2.0,-1.2)
    .. controls (1.0,-1.5) and (0.2,-1.0) .. (0,0) -- cycle;
% Substrat complémentaire
\draw[thick, popgreen, fill=popgreenL]
    (2.9,0.15) -- (3.6,0.15) -- (3.6,0.55) -- (3.2,0.55) -- (3.2,0.9) -- (2.9,0.9) -- cycle;
\node[popblue, font=\small\bfseries] at (1.2,-2.0) {Enzyme native};
\draw[thin, popdark] (2.45,0.35) -- (3.6,1.6) node[right, popdark] {Site actif intact};
\draw[thin, popdark] (3.25,0.9) -- (3.25,1.3);
\node[popgreen] at (3.25,-0.5) {Substrat};
\node[popgreen, align=center] at (1.2,-2.7) {Le substrat s'emboîte\\dans le site actif :\\réaction possible};
\end{scope}
% Flèche chauffage
\draw[->, very thick, poporange] (4.4,0.2) -- (5.8,0.2) node[midway, above, poporange, font=\small\bfseries] {Chaleur};
% Enzyme dénaturée
\begin{scope}[xshift=7cm]
\draw[thick, poporange, fill=poporangeL, rounded corners=3pt, decorate, decoration={random steps, segment length=8pt, amplitude=2pt}]
    (0,0.4) .. controls (0.8,1.4) and (1.8,0.8) .. (2.6,1.1)
    .. controls (3.2,1.3) and (3.4,0.2) .. (2.8,-0.3)
    .. controls (2.2,-0.9) and (1.4,-0.2) .. (0.9,-0.8)
    .. controls (0.4,-1.3) and (-0.2,-0.6) .. (0,0.4) -- cycle;
\node[poporange, font=\small\bfseries] at (1.6,-2.0) {Enzyme dénaturée};
\draw[thin, popdark] (2.6,0.9) -- (3.8,1.6) node[right, popdark] {Site actif détruit};
% Substrat qui ne rentre pas
\draw[thick, popgreen, fill=popgreenL]
    (3.6,-0.6) -- (4.3,-0.6) -- (4.3,-0.2) -- (3.9,-0.2) -- (3.9,0.15) -- (3.6,0.15) -- cycle;
\draw[thick, poppink] (3.4,0.5) -- (3.75,0.15);
\draw[thick, poppink] (3.75,0.5) -- (3.4,0.15);
\node[poporange, align=center] at (1.6,-2.7) {Forme spatiale perdue :\\le substrat ne peut plus\\se fixer : pas de réaction};
\end{scope}
\end{tikzpicture}
\end{center}
\legende{Comparaison entre une enzyme native (forme repliée, site actif fonctionnel) et une enzyme dénaturée par la chaleur (chaîne dépliée, site actif déformé). Le substrat ne peut plus se fixer sur l'enzyme dénaturée : la catalyse est impossible.}
\end{popfigure}

\begin{exemplebox}[Exemple chiffré : lecture de valeurs sur la courbe]
Une amylase présente une activité de \textbf{20 \%} à 10 °C, de \textbf{100 \%} à 37 °C et de seulement \textbf{5 \%} à 70 °C. Interprétons :
\begin{itemize}
\item À 10 °C, l'agitation moléculaire est faible : peu de chocs enzyme-substrat, d'où une activité réduite à 20 \% — mais l'enzyme est intacte (effet réversible).
\item À 37 °C, l'activité est maximale (100 \%) : c'est la \textbf{température optimale}.
\item À 70 °C, l'activité tombe à 5 \% : la quasi-totalité des molécules d'enzyme sont \textbf{dénaturées} de façon irréversible.
\end{itemize}
\end{exemplebox}

\courssub{Méthode : exploiter une courbe activité = f(température)}

\begin{methode}[Exploiter une courbe d'activité enzymatique en fonction de la température]
Face à une courbe « activité = f(T) » au Probatoire, procède toujours en quatre étapes :
\begin{enumerate}
\item \textbf{Repérer le maximum} de la courbe : l'abscisse du sommet donne la \textbf{température optimale} de l'enzyme.
\item \textbf{Décrire la branche ascendante} (avant l'optimum) : l'activité augmente car l'élévation de température accroît l'agitation moléculaire et la fréquence des chocs enzyme-substrat. C'est la phase d'\textbf{activation}.
\item \textbf{Décrire la branche descendante} (après l'optimum) : l'activité chute brutalement car la chaleur détruit la structure spatiale de l'enzyme. C'est la phase de \textbf{dénaturation}, irréversible.
\item \textbf{Conclure en distinguant les deux effets} : le froid inactive de façon \emph{réversible} (enzyme ralentie mais intacte), la chaleur excessive inactive de façon \emph{irréversible} (enzyme dénaturée).
\end{enumerate}
\end{methode}

\courssub{Applications concrètes dans la vie quotidienne et en santé}

\begin{itemize}
\item \textbf{Conservation des aliments au froid.} Le réfrigérateur (4 °C) et le congélateur ($-18$ °C) ralentissent l'activité des enzymes des microbes et des tissus, sans les détruire : la décomposition est retardée. C'est pourquoi une viande décongelée doit être consommée rapidement : dès le réchauffement, les enzymes redeviennent actives !
\item \textbf{Cuisson des aliments.} La cuisson (au-delà de 60--70 °C) dénature les enzymes et les protéines des microbes : c'est une méthode de destruction irréversible des micro-organismes, essentielle pour l'hygiène alimentaire.
\item \textbf{Fièvre et hyperthermie.} La fièvre modérée (38--39 °C) accélère légèrement les réactions de défense de l'organisme. Mais au-delà de \textbf{42 °C} (hyperthermie grave, coup de chaleur), les enzymes commencent à se dénaturer : les réactions vitales ralentissent puis s'arrêtent, ce qui peut entraîner la mort. C'est pourquoi une fièvre très élevée chez un enfant est une urgence médicale dans nos centres de santé.
\end{itemize}

\begin{savaistu}[Des enzymes qui résistent à 80 °C : la révolution de la PCR]
Certaines bactéries vivent dans les sources chaudes (geysers, sources hydrothermales) à des températures proches de 80 °C. Leurs enzymes sont naturellement stables à la chaleur : leur structure spatiale est renforcée par des liaisons supplémentaires. L'une d'elles, la \emph{Taq polymérase} (extraite de la bactérie \emph{Thermus aquaticus}), est utilisée dans la technique de \cle{PCR} (amplification de l'ADN), qui exige des chauffages répétés à plus de 90 °C. Sans cette enzyme thermostable, le diagnostic moléculaire des infections (comme les tests PCR utilisés dans les laboratoires camerounais) serait impossible. La comparaison des courbes d'activité montre que la température optimale d'une enzyme dépend de l'organisme dont elle provient : environ 37 °C chez l'Homme, mais 70--80 °C chez ces bactéries.
\end{savaistu}

\begin{exoresolu}[Exploitation d'une courbe d'activité enzymatique]
\textbf{Énoncé.} On mesure l'activité de deux enzymes en fonction de la température : une enzyme humaine (enzyme A) et une enzyme d'une bactérie de source chaude (enzyme B). Résultats (en \% de l'activité maximale de chaque enzyme) :

\begin{center}
\begin{tabular}{|l|c|c|c|c|c|c|}
\hline
\rowcolor{popblueL} Température (°C) & 20 & 37 & 50 & 60 & 75 & 90 \\
\hline
Enzyme A & 60 & 100 & 45 & 10 & 0 & 0 \\
\hline
Enzyme B & 10 & 30 & 60 & 85 & 100 & 40 \\
\hline
\end{tabular}
\end{center}

\begin{enumerate}
\item Déterminer la température optimale de chaque enzyme.
\item Expliquer l'activité nulle de l'enzyme A à 75 °C.
\item L'enzyme A chauffée à 75 °C est ensuite refroidie à 37 °C. Retrouvera-t-elle son activité ? Justifier.
\end{enumerate}

\tcblower
\textbf{Solution détaillée.}

\textbf{1. Températures optimales.} La température optimale correspond au maximum d'activité (100 \%) :
\begin{itemize}
\item Enzyme A (humaine) : température optimale = \textbf{37 °C}, ce qui correspond à la température du corps humain.
\item Enzyme B (bactérie de source chaude) : température optimale = \textbf{75 °C}, adaptée au milieu très chaud où vit cette bactérie.
\end{itemize}

\textbf{2. Activité nulle de l'enzyme A à 75 °C.} À 75 °C, l'agitation thermique est si intense qu'elle rompt les liaisons fragiles qui maintiennent la structure tridimensionnelle de l'enzyme A. La protéine se déplie, son site actif est déformé, le substrat ne peut plus se fixer : c'est la \textbf{dénaturation thermique}. L'activité est donc nulle.

\textbf{3. Effet du refroidissement.} Non, l'enzyme A ne retrouvera \textbf{pas} son activité à 37 °C, car la dénaturation thermique est \textbf{irréversible} : la structure spatiale détruite ne se reforme pas spontanément. Il faut bien distinguer ce cas de celui d'une enzyme simplement refroidie (par exemple à 0 °C), qui est inactive de façon \emph{réversible} et retrouve son activité au réchauffement.
\end{exoresolu}

\begin{aretenir}[L'essentiel de la section]
\begin{itemize}
\item L'activité d'une enzyme dépend de la \cle{température} : elle croît jusqu'à un \cle{optimum} (environ 37 °C pour les enzymes humaines), puis chute brutalement.
\item \textbf{Branche ascendante} : la température augmente l'agitation moléculaire et la fréquence des chocs enzyme-substrat.
\item \textbf{Branche descendante} : la chaleur excessive détruit la structure spatiale de l'enzyme : c'est la \cle{dénaturation}, \textbf{irréversible}.
\item Le \textbf{froid} ne détruit pas les enzymes : il les ralentit de façon \textbf{réversible} (principe du réfrigérateur).
\item Applications : conservation au froid, cuisson, danger de l'hyperthermie au-delà de 42 °C.
\item La température optimale varie selon l'organisme : certaines bactéries de sources chaudes ont des enzymes actives jusqu'à 80 °C (utilisées en PCR).
\end{itemize}
\end{aretenir}

Maintenant que nous avons compris le rôle de la température, une question se pose naturellement : la température est-elle le seul facteur du milieu qui influence les enzymes ? Que se passe-t-il si l'on modifie l'acidité du milieu ou la quantité de substrat disponible ? C'est ce que nous étudierons dans la section suivante, consacrée à l'influence du pH et des concentrations.

\coursec{Influence du pH et des concentrations}

Dans la section précédente, nous avons vu que la température agit comme un véritable thermostat sur l'activité enzymatique : trop froide, l'enzyme est ralentie ; trop chaude, elle est dénaturée. Mais la température n'est pas le seul paramètre du milieu qui compte. Quiconque a déjà ressenti les brûlures d'estomac sait que notre tube digestif abrite des milieux très différents : l'estomac est un véritable bain d'acide, alors que l'intestin grêle est au contraire légèrement basique. Comment les enzymes digestives peuvent-elles fonctionner dans des milieux aussi opposés ? Et que se passe-t-il lorsque la quantité de nourriture — donc de substrat — varie d'un repas à l'autre ? C'est à ces deux questions que cette section est consacrée : l'influence du \cle{pH} et l'influence des \cle{concentrations} (en substrat et en enzyme) sur la vitesse des réactions enzymatiques.

\courssub{Influence du pH sur l'activité enzymatique}

\textbf{Observation et mise en situation.} Un comprimé antiacide (comme ceux vendus en pharmacie à Yaoundé ou Douala contre les brûlures d'estomac) neutralise l'acidité gastrique. Or les médecins déconseillent d'en abuser : un estomac trop peu acide digère mal les protéines. Pourquoi ? Parce que la \cle{pepsine}, l'enzyme qui découpe les protéines dans l'estomac, n'est active qu'en milieu très acide. Autre observation : l'amylase salivaire, qui commence la digestion de l'amidon du foutou ou du riz dès la bouche, cesse presque totalement d'agir une fois arrivée dans l'estomac. Le pH du milieu semble donc déterminer si une enzyme peut travailler ou non.

\textbf{Hypothèse.} Chaque enzyme possède une valeur de pH à laquelle elle est la plus efficace, et son activité diminue dès que le pH s'en écarte.

\begin{experience}[Activité d'une enzyme en fonction du pH]
\textbf{Matériel :} une solution d'enzyme (par exemple une amylase), une solution de substrat (amidon), des tubes à essai, des solutions tampons de pH 2, 4, 6, 7, 8, 10 et 12, un bain-marie à \SI{37}{\celsius}, un réactif permettant de doser le produit formé (ou la disparition du substrat, par exemple l'eau iodée pour l'amidon).\par
\textbf{Protocole :}
\begin{enumerate}
\item Répartir la même quantité de substrat dans 7 tubes.
\item Ajouter dans chaque tube un tampon de pH différent (2, 4, 6, 7, 8, 10, 12).
\item Ajouter la même quantité d'enzyme dans chaque tube et placer l'ensemble au bain-marie à \SI{37}{\celsius}.
\item Après un temps fixe, mesurer la quantité de produit formé (ou de substrat disparu) dans chaque tube.
\item Reporter les résultats sur un graphique : activité enzymatique (en \%) en ordonnée, pH en abscisse.
\end{enumerate}
\textbf{Observations :} l'activité est faible aux pH extrêmes, augmente progressivement, passe par un maximum à un pH donné, puis diminue à nouveau. La courbe obtenue a une allure \emph{en cloche}.\par
\textbf{Interprétation :} l'enzyme n'est pleinement efficace que dans un intervalle étroit de pH, centré sur une valeur caractéristique.
\end{experience}

\begin{definition}[pH optimal]
Le \cle{pH optimal} d'une enzyme est la valeur du pH à laquelle son activité est \textbf{maximale}. Il dépend du milieu dans lequel l'enzyme agit normalement dans l'organisme : la pepsine gastrique a un pH optimal voisin de 2 (milieu acide de l'estomac), l'amylase salivaire un pH optimal voisin de 7 (neutralité de la salive), et la trypsine intestinale un pH optimal voisin de 8 (milieu légèrement basique de l'intestin grêle).
\end{definition}

\begin{exemplebox}[Trois enzymes digestives, trois pH optima]
La digestion des aliments se déroule dans trois milieux de pH différents, et chaque enzyme digestive est adaptée à son milieu :
\begin{itemize}
\item la \textbf{pepsine} de l'estomac : pH optimal $\approx 2$. Son activité est de 100 \% à pH 2, tombe à environ 40 \% à pH 4 et devient quasi nulle à pH 7 ;
\item l'\textbf{amylase salivaire} de la bouche : pH optimal $\approx 7$ (salive neutre) ;
\item la \textbf{trypsine} de l'intestin grêle : pH optimal $\approx 8$ (milieu rendu basique par le suc pancréatique).
\end{itemize}
Ainsi, aucune de ces enzymes ne peut remplacer une autre : chacune travaille « à son étage » du tube digestif.
\end{exemplebox}

\begin{popfigure}
\begin{center}
\begin{tikzpicture}
\begin{axis}[
  width=0.95\linewidth, height=7cm,
  xlabel={pH du milieu}, ylabel={Activité enzymatique (\%)},
  xmin=0, xmax=13, ymin=0, ymax=115,
  xtick={0,2,4,6,7,8,10,12},
  ytick={0,20,40,60,80,100},
  grid=both, grid style={popline!40},
  legend style={at={(0.97,0.97)}, anchor=north east, font=\small, draw=popline},
  axis line style={popdark},
]
\addplot[poporange, very thick, smooth, samples=200, domain=0.5:6.5]
  {100*exp(-((x-2)^2)/(2*1.1^2))};
\addlegendentry{Pepsine (estomac)}
\addplot[popblue, very thick, smooth, samples=200, domain=3:11]
  {100*exp(-((x-7)^2)/(2*1.3^2))};
\addlegendentry{Amylase salivaire (bouche)}
\addplot[popgreen, very thick, smooth, samples=200, domain=4:12.5]
  {100*exp(-((x-8)^2)/(2*1.2^2))};
\addlegendentry{Trypsine (intestin)}
\draw[dashed, popmuted] (axis cs:2,0) -- (axis cs:2,100);
\draw[dashed, popmuted] (axis cs:7,0) -- (axis cs:7,100);
\draw[dashed, popmuted] (axis cs:8,0) -- (axis cs:8,100);
\end{axis}
\end{tikzpicture}
\end{center}
\legende{Activité de trois enzymes digestives en fonction du pH. Chaque courbe en cloche présente un maximum : le pH optimal (environ 2 pour la pepsine, 7 pour l'amylase salivaire, 8 pour la trypsine). Les pointillés verticaux repèrent les pH optima.}
\end{popfigure}

\textbf{Interprétation moléculaire.} Pourquoi le pH agit-il ainsi ? Rappelons que le pH mesure la concentration en ions \ce{H+} du milieu. Or le site actif d'une enzyme porte des \cle{charges électriques} (groupes ionisables des acides aminés, comme $-\ce{COO}^{-}$ ou $-\ce{NH3}^{+}$) qui participent à la fixation du substrat et à la catalyse. Lorsque le pH s'écarte de l'optimum :
\begin{itemize}
\item les ions \ce{H+} en excès (milieu acide) ou en défaut (milieu basique) \textbf{modifient les charges} du site actif : le substrat ne se fixe plus correctement, l'activité chute ;
\item à des pH très éloignés de l'optimum, les liaisons qui maintiennent la structure tridimensionnelle de la protéine (ponts ioniques, liaisons hydrogène) sont rompues : l'enzyme se \textbf{dénature}, c'est-à-dire qu'elle perd définitivement sa forme et donc sa fonction.
\end{itemize}

\begin{attention}[pH extrême = dénaturation, pas simple ralentissement]
Un pH extrême ne « ralentit » pas seulement l'enzyme : il peut la \textbf{dénaturer de façon irréversible}. Une pepsine placée à pH 12 ou une trypsine plongée dans l'acide gastrique concentré sont détruites : ramener ensuite le pH à la valeur optimale ne restaure pas l'activité. Ne confonds pas, dans tes rédactions, « activité réduite » (réversible, près de l'optimum) et « dénaturation » (irréversible, aux pH extrêmes).
\end{attention}

\begin{methode}[Exploiter un graphe d'activité en fonction du pH]
Pour comparer deux enzymes sur une courbe activité = f(pH) :
\begin{enumerate}
\item Repérer le \textbf{sommet} de chaque cloche : c'est le pH optimal de l'enzyme.
\item Noter l'\textbf{intervalle de pH} où l'activité dépasse 50 \% du maximum : c'est la zone de fonctionnement efficace.
\item Relier chaque pH optimal au \textbf{milieu physiologique} où l'enzyme agit : estomac acide (pH $\approx 2$), bouche neutre (pH $\approx 7$), intestin basique (pH $\approx 8$).
\item Conclure sur l'adaptation de l'enzyme à son milieu et sur ce qui se passerait si le pH du milieu changeait (ex. : antiacides, reflux gastrique).
\end{enumerate}
\end{methode}

\begin{savaistu}[Une digestion en étapes grâce au pH]
C'est précisément grâce aux pH différents de l'estomac et de l'intestin que la digestion se déroule en \textbf{étapes successives et coordonnées} : l'acidité gastrique active la pepsine et stérilise partiellement le bol alimentaire ; puis le suc pancréatique, riche en bicarbonates, neutralise cette acidité dans le duodénum, ce qui « éteint » la pepsine et « allume » la trypsine et les autres enzymes intestinales. Le pH agit comme un interrupteur qui allume et éteint les enzymes au bon moment. C'est aussi pourquoi les ulcères d'estomac, favorisés par la bactérie \emph{Helicobacter pylori}, sont si douloureux : la muqueuse n'est plus protégée de l'acide et de la pepsine, qui se mettent à digérer la paroi elle-même.
\end{savaistu}

\courssub{Influence de la concentration en substrat}

\textbf{Situation concrète.} Après un jeûne prolongé, un grand plat de ndolé ou de couscous de maïs apporte soudain une énorme quantité de substrats aux enzymes digestives. La vitesse de digestion augmente-t-elle indéfiniment avec la quantité de nourriture ? L'expérience montre que non.

\begin{experience}[Vitesse de réaction en fonction de la concentration en substrat]
\textbf{Protocole :} on prépare des tubes contenant tous la \textbf{même quantité d'enzyme} (concentration en enzyme constante) mais des concentrations croissantes de substrat $[S]$. On mesure la \textbf{vitesse initiale} $V$ de la réaction (quantité de produit formé par unité de temps, au tout début de la réaction) dans chaque tube.\par
\textbf{Observations :}
\begin{itemize}
\item aux faibles valeurs de $[S]$, la vitesse $V$ augmente presque proportionnellement à $[S]$ ;
\item puis l'augmentation ralentit ;
\item enfin, $V$ atteint un \textbf{plateau} : une vitesse maximale $V_{\max}$ qui ne dépend plus de $[S]$.
\end{itemize}
\textbf{Interprétation :} tant qu'il reste des sites actifs libres, ajouter du substrat accélère la réaction. Mais lorsque \textbf{tous les sites actifs sont occupés en permanence}, les enzymes travaillent à plein régime : ajouter encore du substrat ne sert à rien, il faut attendre qu'un site se libère. On dit que l'enzyme est \cle{saturée}.
\end{experience}

\begin{propriete}[Saturation enzymatique — admise]
À \textbf{concentration en enzyme constante}, la vitesse d'une réaction enzymatique croît avec la concentration en substrat jusqu'à atteindre une \textbf{vitesse maximale} $V_{\max}$, correspondant à la \cle{saturation} de tous les sites actifs. Cette propriété, admise en Première D, s'interprète par le modèle du site actif : chaque molécule d'enzyme ne peut transformer qu'un nombre limité de molécules de substrat par unité de temps.
\end{propriete}

\begin{popfigure}
\begin{center}
\begin{tikzpicture}
\begin{axis}[
  width=0.95\linewidth, height=7cm,
  xlabel={Concentration en substrat $[S]$ (mmol/L)},
  ylabel={Vitesse de réaction $V$ (unités arbitraires)},
  xmin=0, xmax=12, ymin=0, ymax=120,
  xtick={0,2,4,6,8,10,12},
  ytick={0,20,40,60,80,100},
  grid=both, grid style={popline!40},
  axis line style={popdark},
]
\addplot[popblue, very thick, smooth, samples=200, domain=0:12]
  {100*x/(2+x)};
\addplot[dashed, poporange, thick] coordinates {(0,100) (12,100)};
\node[poporange, anchor=south west, font=\small] at (axis cs:6.2,101) {$V_{\max}$ : plateau de saturation};
\draw[dashed, popmuted] (axis cs:2,0) -- (axis cs:2,50) -- (axis cs:0,50);
\node[popdark, font=\small, anchor=west] at (axis cs:4.5,35) {Zone où $V$ dépend de $[S]$};
\end{axis}
\end{tikzpicture}
\end{center}
\legende{Vitesse d'une réaction enzymatique en fonction de la concentration en substrat, à concentration en enzyme constante. La vitesse croît d'abord presque proportionnellement, puis plafonne à $V_{\max}$ lorsque tous les sites actifs sont saturés.}
\end{popfigure}

\textbf{Lecture de la courbe.} Deux zones sont à connaître :
\begin{itemize}
\item \textbf{zone initiale} (faible $[S]$) : beaucoup de sites actifs sont libres ; la vitesse est limitée par la quantité de substrat. Ajouter du substrat accélère la réaction ;
\item \textbf{plateau} (forte $[S]$) : tous les sites actifs sont occupés en permanence ; la vitesse est limitée par la \textbf{quantité d'enzyme}. Seule l'ajout d'enzyme supplémentaire pourrait augmenter $V_{\max}$.
\end{itemize}

\courssub{Influence de la concentration en enzyme}

\textbf{Observation.} Dans un centre de santé, un laborantin qui dose le glucose sait qu'avec deux fois plus de réactif enzymatique (glucose-oxydase), la réaction de dosage va deux fois plus vite — à condition que le glucose soit en excès. C'est une conséquence directe de la saturation vue ci-dessus.

\begin{aretenir}[Proportionnalité vitesse--enzyme]
Lorsque le substrat est en \textbf{excès} (enzymes saturées), la vitesse de la réaction est \textbf{proportionnelle à la concentration en enzyme} : doubler la quantité d'enzyme double la vitesse, tripler la quantité d'enzyme la triple, car on double ou triple le nombre de sites actifs disponibles. La valeur du plateau $V_{\max}$ dépend donc directement de la quantité d'enzyme présente.
\end{aretenir}

\textbf{Synthèse.} La vitesse d'une réaction enzymatique est limitée par le facteur le plus rare :
\begin{itemize}
\item si le substrat est rare, c'est lui qui limite : $V$ dépend de $[S]$ ;
\item si l'enzyme est rare (substrat en excès), c'est elle qui limite : $V$ dépend de $[E]$.
\end{itemize}
Cette logique de « facteur limitant » est essentielle pour comprendre comment la cellule règle son métabolisme : en ajustant la quantité d'enzymes qu'elle fabrique, elle ajuste la vitesse de ses réactions de synthèse et de dégradation.

\begin{exoresolu}[Exploitation de courbes : pH et saturation]
\textbf{Énoncé.} On mesure l'activité de la pepsine à différents pH : 100 \% à pH 2 ; 80 \% à pH 3 ; 40 \% à pH 4 ; 5 \% à pH 7. Par ailleurs, à pH 2 et à concentration en pepsine constante, on mesure la vitesse de digestion des protéines pour des concentrations croissantes en protéines : $V = 20$ ; $50$ ; $75$ ; $90$ ; $98$ ; $100$ ; $100$ unités pour $[S] = 0,5$ ; $1$ ; $2$ ; $3$ ; $5$ ; $8$ ; $12$ mmol/L.\par
1. Déterminer le pH optimal de la pepsine et le relier à son milieu d'action.\par
2. Que vaut approximativement $V_{\max}$ ? Pourquoi la vitesse plafonne-t-elle ?\par
3. Comment pourrait-on dépasser cette vitesse maximale ?
\tcblower
\textbf{Solution détaillée.}\par
1. L'activité est maximale (100 \%) à \textbf{pH 2} : c'est le pH optimal de la pepsine. Cette valeur très acide correspond au milieu gastrique, dont l'acidité est due à l'acide chlorhydrique sécrété par l'estomac. La pepsine est donc adaptée à son milieu physiologique. À pH 7 (neutralité), son activité est quasi nulle (5 \%) : elle ne peut pas fonctionner dans l'intestin.\par
2. La vitesse plafonne à $V_{\max} \approx 100$ unités : au-delà de $[S] = 8$ mmol/L, augmenter la concentration en protéines n'accélère plus la réaction. Ce plateau traduit la \textbf{saturation des sites actifs} : toutes les molécules de pepsine sont occupées en permanence, et la vitesse est alors limitée par la quantité d'enzyme, non par le substrat.\par
3. Pour dépasser $V_{\max}$, il faut \textbf{augmenter la concentration en enzyme} : avec davantage de pepsine (et le substrat toujours en excès), le nombre de sites actifs augmente et le nouveau plateau est plus élevé, la vitesse étant proportionnelle à $[E]$.
\end{exoresolu}

\begin{attention}[Pièges fréquents au Probatoire]
\begin{itemize}
\item Ne dis pas que « le pH optimal est toujours 7 » : il \textbf{dépend de l'enzyme} et de son milieu (2 pour la pepsine, 8 pour la trypsine).
\item Sur une courbe $V = f([S])$, le plateau ne signifie pas que l'enzyme est « épuisée » ou « détruite » : elle est \textbf{saturée}, c'est-à-dire occupée au maximum de ses capacités, mais elle continue de fonctionner indéfiniment.
\item N'oublie pas la condition d'application : la proportionnalité entre $V$ et $[E]$ n'est valable que si le \textbf{substrat est en excès}.
\item Dans l'exploitation d'une courbe en cloche, cite toujours la \textbf{valeur numérique} du pH optimal lue sur le graphe, puis relie-la au milieu physiologique.
\end{itemize}
\end{attention}

En résumé, le pH et les concentrations sont, avec la température, les grands paramètres qui conditionnent l'activité enzymatique. Chaque enzyme possède un pH optimal adapté à son milieu, et sa vitesse de travail est limitée tantôt par le substrat, tantôt par sa propre quantité. Ces propriétés expliquent, comme nous le verrons dans la section suivante, comment la cellule orchestre finement les réactions de synthèse et de dégradation qui assurent le renouvellement permanent de ses molécules.

\coursec{Rôle des enzymes dans le renouvellement moléculaire}

Dans la section précédente, nous avons vu que la température, le pH et les concentrations modulent finement l'activité de chaque enzyme. Mais à quoi servent concrètement ces catalyseurs dans la vie quotidienne de la cellule ? Pourquoi l'organisme fabrique-t-il des milliers d'enzymes différentes ? C'est ce que nous allons découvrir maintenant : les enzymes sont les ouvrières permanentes du \cle{renouvellement moléculaire}, ce grand chantier de construction et de démolition qui ne s'arrête jamais, ni le jour, ni la nuit, tant que nous sommes vivants.

\courssub{La cellule : un chantier permanent de construction et de démolition}

\textbf{Une observation de la vie courante.} Observe un élève de Première D pendant l'année scolaire : il grandit, ses muscles se développent après le sport, une coupure au doigt cicatrise en quelques jours, ses cheveux et ses ongles poussent sans cesse. Pourtant, il mange chaque jour du manioc, du ndolé, du poisson... Comment ces aliments deviennent-ils de la peau, du muscle, du sang ? Et d'où vient l'énergie qui lui permet de courir, de réfléchir, de maintenir sa température à 37~°C sous le soleil de Yaoundé ?

\textbf{Hypothèse.} À l'intérieur de chaque cellule, des réactions chimiques doivent en permanence \emph{détruire} certaines molécules (pour en tirer de l'énergie et des matériaux) et en \emph{construire} d'autres (pour fabriquer les constituants de l'organisme).

\textbf{Données.} Les biochimistes ont mesuré qu'un adulte renouvelle chaque jour environ 200 à 300~g de protéines : des protéines « usées » sont dégradées en acides aminés, et de nouvelles protéines sont synthétisées à partir de ces acides aminés ou de ceux apportés par l'alimentation. De même, les membranes cellulaires, les globules rouges (durée de vie : environ 120 jours) sont constamment détruits et refabriqués.

\textbf{Conclusion.} La matière vivante n'est pas figée : elle est le siège d'un flux permanent de réactions de construction et de destruction. L'ensemble de ces réactions constitue le \cle{métabolisme}, qui se divise en deux grands ensembles complémentaires.

\begin{definition}[Anabolisme et catabolisme]
L'\cle{anabolisme} regroupe l'ensemble des réactions de \textbf{synthèse moléculaire} : de petites molécules simples (acides aminés, glucose, acides gras) sont assemblées pour fabriquer de grosses molécules complexes (protéines, glycogène, lipides membranaires). L'anabolisme \textbf{consomme de l'énergie}.
\par\smallskip
Le \cle{catabolisme} regroupe l'ensemble des réactions de \textbf{dégradation moléculaire} : de grosses molécules (glucides, protéines, lipides) sont fragmentées en petites molécules. Le catabolisme \textbf{libère de l'énergie}.
\par\smallskip
Anabolisme et catabolisme forment ensemble le \cle{métabolisme}, et chacune de leurs réactions est catalysée par une enzyme.
\end{definition}

\begin{popfigure}
\begin{center}
\begin{tikzpicture}[font=\small, >=Stealth]
% Branche anabolisme (gauche)
\node[draw=popgreen, fill=popgreenL, rounded corners, thick, minimum width=3.6cm, minimum height=1.1cm, align=center] (ana) at (-3.8,2.2) {\textbf{ANABOLISME}\\ (synthèse)};
\node[draw=popdark, fill=popcream, rounded corners, align=center, minimum width=3.2cm] (petites) at (-3.8,0) {Pool des petites molécules\\ acides aminés, glucose,\\ acides gras, ATP/ADP};
\node[draw=popdark, fill=popcream, rounded corners, align=center, minimum width=3.2cm] (grosses) at (-3.8,4.4) {Grosses molécules\\ protéines, glycogène,\\ lipides membranaires};
\draw[->, very thick, popgreen] (petites) -- (ana);
\draw[->, very thick, popgreen] (ana) -- (grosses);
\node[draw=poporange, fill=poporangeL, rounded corners, align=center, font=\footnotesize] (energiein) at (-7.2,2.2) {Énergie\\ \textbf{consommée (ATP)}};
\draw[->, thick, poporange] (energiein) -- (ana);
% Branche catabolisme (droite)
\node[draw=poppink, fill=poppinkL, rounded corners, thick, minimum width=3.6cm, minimum height=1.1cm, align=center] (cata) at (3.8,2.2) {\textbf{CATABOLISME}\\ (dégradation)};
\node[draw=popdark, fill=popcream, rounded corners, align=center, minimum width=3.2cm] (grosses2) at (3.8,4.4) {Grosses molécules\\ glucides, protéines,\\ lipides de réserve};
\node[draw=popdark, fill=popcream, rounded corners, align=center, minimum width=3.2cm] (petites2) at (3.8,0) {Pool des petites molécules\\ \ce{CO2}, \ce{H2O}, urée,\\ acides aminés, ATP/ADP};
\draw[->, very thick, poppink] (grosses2) -- (cata);
\draw[->, very thick, poppink] (cata) -- (petites2);
\node[draw=popgold, fill=popgoldL, rounded corners, align=center, font=\footnotesize] (energieout) at (7.2,2.2) {Énergie\\ \textbf{libérée (ATP,}\\\textbf{ chaleur)}};
\draw[->, thick, popgold] (cata) -- (energieout);
% Lien central : pool commun
\draw[<->, dashed, thick, popmuted] (petites) -- (petites2) node[midway, below, font=\footnotesize\itshape, text=popmuted, align=center] {Pool commun des\\ petites molécules\\ et de l'ATP/ADP};
\end{tikzpicture}
\end{center}
\legende{Les deux branches complémentaires du métabolisme. À gauche, l'anabolisme assemble de petites molécules en grosses molécules en consommant de l'énergie (ATP). À droite, le catabolisme dégrade les grosses molécules, libère de l'énergie (ATP, chaleur) et restitue les petites molécules au pool commun. Toutes les flèches pleines représentent des réactions catalysées par des enzymes spécifiques.}
\end{popfigure}

\courssub{Les enzymes de dégradation : démolir pour récupérer}

\textbf{Situation concrète.} Au petit-déjeuner, tu manges du pain (riche en amidon) et des {\oe}ufs (riches en protéines). Une heure plus tard, ton sang transporte du glucose et des acides aminés. L'amidon et les protéines, molécules trop grosses pour traverser la paroi de l'intestin, ont été découpés en petits fragments. Qui a fait ce découpage ? Des enzymes de dégradation.

\begin{exemplebox}[Les hydrolyses digestives : un catabolisme extracellulaire]
La digestion est un ensemble d'\cle{hydrolyses} : des enzymes brisent les liaisons des grosses molécules alimentaires en ajoutant de l'eau. Chaque enzyme est spécifique d'un type de molécule :
\begin{itemize}
\item l'\cle{amylase} (salivaire puis pancréatique) hydrolyse l'\textbf{amidon} en maltose puis en glucose ;
\item les \cle{protéases} (pepsine de l'estomac, trypsine du pancréas) hydrolysent les \textbf{protéines} en peptides puis en acides aminés ;
\item les \cle{lipases} hydrolysent les \textbf{lipides} en acides gras et glycérol.
\end{itemize}
Les petites molécules obtenues traversent la paroi intestinale, passent dans le sang et sont distribuées à toutes les cellules de l'organisme.
\end{exemplebox}

\textbf{Et à l'intérieur des cellules ?} Le catabolisme ne s'arrête pas à la digestion. Chaque cellule dégrade le glucose pour en extraire l'énergie, lors de la \cle{respiration cellulaire}, dont l'équation bilan est :
\[ \ce{C6H12O6 + 6O2 -> 6CO2 + 6H2O} + \text{énergie (ATP)} \]
Cette réaction, qui semble simple écrite ainsi, est en réalité une succession d'une vingtaine de réactions (glycolyse, cycle de Krebs, chaîne respiratoire), chacune catalysée par une enzyme spécifique. L'énergie libérée est stockée dans une molécule « pile » rechargeable, l'\cle{ATP} (adénosine triphosphate), qui alimente ensuite tous les travaux cellulaires : contraction musculaire, transport actif, synthèses.

\begin{aretenir}[Le double but du catabolisme]
Les dégradations enzymatiques, digestives ou intracellulaires, ont deux fonctions complémentaires :
\begin{enumerate}
\item \textbf{fournir de l'énergie} (stockée dans l'ATP) en oxydant les molécules énergétiques comme le glucose ;
\item \textbf{fournir de petites molécules} (glucose, acides aminés, acides gras) qui serviront de « briques » pour les synthèses de l'anabolisme.
\end{enumerate}
\end{aretenir}

\courssub{Les enzymes de synthèse : construire l'édifice vivant}

\textbf{Situation concrète.} Un enfant atteint de kwashiorkor (malnutrition protéique, encore observée dans certaines zones rurales du Cameroun) présente un ventre ballonné, des muscles fondus, des cheveux décolorés. Pourquoi ? Parce que son alimentation ne lui fournit pas assez d'acides aminés : ses cellules ne peuvent plus fabriquer suffisamment de protéines, alors que la synthèse des protéines est indispensable à la construction et au renouvellement des tissus.

La fabrication des protéines illustre parfaitement l'anabolisme : à partir des 20 acides aminés disponibles, la cellule assemble, dans un ordre précis dicté par ses gènes, des chaînes de centaines d'acides aminés. Chaque étape de cet assemblage (activation des acides aminés, formation de la liaison peptidique au ribosome) fait intervenir des enzymes ou des complexes enzymatiques. De même :
\begin{itemize}
\item le \textbf{glycogène} (réserve de glucides du foie et des muscles) est synthétisé à partir du glucose par la \textbf{glycogène synthase} (enzyme de synthèse), et dégradé en glucose-1-phosphate par la \textbf{phosphorylase} (enzyme de dégradation) quand l'organisme en a besoin ;
\item les \textbf{lipides des membranes} (phospholipides, cholestérol) sont fabriqués à partir d'acides gras et de glycérol-3-phosphate, ce qui permet la croissance des cellules et la réparation des membranes usées.
\end{itemize}

\begin{attention}[L'anabolisme consomme, le catabolisme libère]
Ne jamais inverser : les réactions de \textbf{synthèse} (anabolisme) \textbf{consomment} de l'énergie — il faut de l'ATP (et souvent de l'UTP ou GTP) pour assembler des molécules, comme il faut du ciment et de la main-d'{\oe}uvre pour monter un mur. Les réactions de \textbf{dégradation} (catabolisme) \textbf{libèrent} de l'énergie. Les deux sont couplés : l'énergie libérée par le catabolisme (production d'ATP) alimente l'anabolisme (consommation d'ATP).
\end{attention}

\courssub{Les chaînes enzymatiques : les voies métaboliques}

\textbf{Observation.} Comment la cellule transforme-t-elle la phénylalanine (un acide aminé) en mélanine (le pigment qui colore la peau) ? Pas en une seule réaction, mais par une succession d'étapes. Les biochimistes ont découvert que pratiquement toutes les transformations cellulaires fonctionnent ainsi, comme une \emph{chaîne de montage} (ou de démontage) dans une usine.

\begin{definition}[Voie métabolique]
Une \cle{voie métabolique} (ou chaîne enzymatique) est une succession ordonnée de réactions chimiques dans laquelle le produit d'une réaction devient le substrat de la réaction suivante :
\[ A \xrightarrow{E_1} B \xrightarrow{E_2} C \xrightarrow{E_3} D \]
$A$ est le \textbf{substrat initial}, $D$ le \textbf{produit final}, $B$ et $C$ les \textbf{intermédiaires métaboliques}, et chaque étape est catalysée par une enzyme spécifique ($E_1$, $E_2$, $E_3$).
\end{definition}

\begin{popfigure}
\begin{center}
\begin{tikzpicture}[font=\small, >=Stealth, node distance=2.8cm]
\node[draw=popblue, fill=popblueL, circle, thick, minimum size=1.1cm] (A) {A};
\node[draw=popblue, fill=popblueL, circle, thick, minimum size=1.1cm, right=of A] (B) {B};
\node[draw=popblue, fill=popblueL, circle, thick, minimum size=1.1cm, right=of B] (C) {C};
\node[draw=popgreen, fill=popgreenL, circle, thick, minimum size=1.1cm, right=of C] (D) {D};
\draw[->, very thick, popdark] (A) -- (B) node[midway, above, font=\small\bfseries, text=poporange] {$E_1$};
\draw[->, very thick, popdark] (B) -- (C) node[midway, above, font=\small\bfseries, text=poporange] {$E_2$};
\draw[->, very thick, popdark] (C) -- (D) node[midway, above, font=\small\bfseries, text=poporange] {$E_3$};
\node[below=0.3cm of A, font=\footnotesize\itshape, text=popmuted, align=center] {substrat\\ initial};
\node[below=0.3cm of B, font=\footnotesize\itshape, text=popmuted, align=center] {intermédiaire};
\node[below=0.3cm of C, font=\footnotesize\itshape, text=popmuted, align=center] {intermédiaire};
\node[below=0.3cm of D, font=\footnotesize\itshape, text=popmuted, align=center] {produit\\ final};
\node[above=0.7cm of B, font=\footnotesize\itshape, text=poporange, align=center] {chaque étape est catalysée\\ par une enzyme spécifique};
\end{tikzpicture}
\end{center}
\legende{Schéma d'une voie métabolique linéaire. Le substrat initial A est transformé en produit final D par une succession de réactions ; chaque flèche correspond à une réaction catalysée par une enzyme spécifique ($E_1$, $E_2$, $E_3$). Le produit d'une étape est le substrat de l'étape suivante.}
\end{popfigure}

\begin{propriete}[Spécificité et enchaînement des étapes]
Chaque étape d'une voie métabolique est catalysée par \textbf{une enzyme spécifique}. Par conséquent, le \textbf{blocage d'une seule étape bloque toute la chaîne} : l'intermédiaire situé en amont s'accumule (parfois de façon toxique) et le produit final n'est plus fabriqué. C'est exactement ce qui se passe dans les maladies métaboliques héréditaires (section suivante).
\end{propriete}

\begin{methode}[Analyser une voie métabolique]
Face à un schéma de voie métabolique au Probatoire, procède toujours dans l'ordre :
\begin{enumerate}
\item \textbf{Identifier le substrat initial} (la molécule de départ) et le \textbf{produit final} (la molécule fabriquée ou le déchet éliminé).
\item \textbf{Repérer l'enzyme de chaque étape} (nom au-dessus ou sous chaque flèche) et vérifier que chaque flèche correspond à une seule transformation.
\item \textbf{Prévoir les conséquences d'un blocage} : si l'enzyme $E_2$ est absente ou inactive, l'intermédiaire $B$ s'accumule et les produits $C$ et $D$ ne sont plus formés.
\item \textbf{Classer la voie} : s'agit-il d'une voie de dégradation (catabolisme, libération d'énergie) ou de synthèse (anabolisme, consommation d'énergie) ?
\end{enumerate}
\end{methode}

\begin{exoresolu}[La voie de dégradation de la phénylalanine : la phénylcétonurie]
Chez l'Homme, l'acide aminé phénylalanine (apporté par les protéines alimentaires) est dégradé selon la voie : phénylalanine $\xrightarrow{E_1\ (\text{phénylalanine hydroxylase})}$ tyrosine $\xrightarrow{E_2}$ ... $\xrightarrow{E_n}$ \ce{CO2} + \ce{H2O}. Chez certains enfants, l'enzyme $E_1$ est absente ou inactive par mutation du gène \emph{PAH} (chromosome 12).
\begin{enumerate}
\item Quel intermédiaire s'accumule dans l'organisme de ces enfants ?
\item Quelle substance manque à leur organisme ?
\item Propose une mesure diététique pour limiter les dégâts neurologiques.
\end{enumerate}
\tcblower
\textbf{Solution détaillée.}
\begin{enumerate}
\item L'enzyme $E_1$ (phénylalanine hydroxylase) catalyse la transformation phénylalanine $\to$ tyrosine. Si $E_1$ est absente, le substrat situé \emph{en amont} du blocage, la \textbf{phénylalanine}, s'accumule dans le sang et le cerveau. En excès, elle est transformée en dérivés toxiques (phénylpyruvate, phénylacétate) qui perturbent le développement du système nerveux : c'est la \textbf{phénylcétonurie} (PKU), dépistée à la naissance au Cameroun par le test de Guthrie.
\item Le produit situé \emph{en aval} du blocage, la \textbf{tyrosine}, n'est plus fabriquée par la voie endogène : elle devient un acide aminé \textbf{indispensable} qu'il faut apporter par l'alimentation. La tyrosine est par ailleurs précurseur de la mélanine (pigmentation) et des catécholamines (neurotransmetteurs).
\item Un \textbf{régime strictement pauvre en phénylalanine} (lait maternel ou laits maternisés spéciaux sans phénylalanine, apport contrôlé en protéines naturelles, supplémentation en tyrosine) mis en place \textbf{dès les premières semaines de vie} limite l'accumulation du substrat toxique et permet un développement intellectuel quasi normal. C'est l'application directe de la méthode : réduire l'entrée de la voie quand une enzyme de la chaîne est défaillante.
\end{enumerate}
\end{exoresolu}

\courssub{Le renouvellement moléculaire : une question de vitesse}

\textbf{Le problème posé.} Les réactions du métabolisme — hydrolyser une protéine, oxyder le glucose, assembler un acide aminé à une chaîne en construction — sont chimiquement possibles sans enzyme (thermodynamiquement favorables). Mais les calculs de cinétique chimique montrent qu'à 37~°C et sans catalyseur, la plupart prendraient des \emph{années}, voire des siècles. Or une cellule musculaire qui se contracte a besoin d'énergie en une fraction de seconde, et une plaie doit être réparée en quelques jours.

\textbf{La réponse des enzymes.} En abaissant l'énergie d'activation, une enzyme multiplie la vitesse d'une réaction par un facteur de $10^{6}$ à $10^{12}$ (parfois jusqu'à $10^{17}$ pour les plus efficaces comme la carbonique anhydrase). C'est cette accélération prodigieuse qui rend possible le \cle{renouvellement moléculaire permanent} :
\begin{itemize}
\item les \textbf{protéines} de l'organisme sont continuellement dégradées (protéasome, lysosomes) et resynthétisées : une enzyme du foie peut être entièrement remplacée en quelques heures (demi-vie courte) ;
\item les \textbf{membranes cellulaires} sont renouvelées en permanence : leurs lipides et protéines usés sont détruits et remplacés (turnover membranaire) ;
\item les \textbf{globules rouges}, détruits au rythme de plusieurs millions par seconde dans la rate (hémolyse sénescente), sont remplacés au même rythme par la moelle osseuse (érythropoïèse).
\end{itemize}

\begin{attention}[Une enzyme ne fournit pas d'énergie]
Erreur fréquente au Probatoire : écrire qu'une enzyme « fournit de l'énergie » ou « donne de la force » à la réaction. C'est faux. Une enzyme est un \textbf{catalyseur} : elle \textbf{abaisse l'énergie d'activation} (elle accélère une réaction déjà thermodynamiquement possible) et elle est retrouvée \textbf{inchangée} en fin de réaction. L'énergie, elle, provient des \textbf{molécules elles-mêmes} : elle est libérée par la dégradation des molécules énergétiques (glucose, lipides) et stockée temporairement dans l'ATP.
\end{attention}

\begin{savaistu}[Des protéines qui ne vivent que quelques heures]
Les protéines de ton organisme sont renouvelées en permanence, mais à des rythmes très différents. Certaines enzymes régulatrices du foie (ex. : HMG-CoA réductase, enzymes de la glycolyse) ont une durée de vie de \textbf{quelques heures seulement} ; l'hémoglobine de tes globules rouges vit environ 120 jours ; le collagène de tes os et de ta peau, plusieurs années. Au total, même sans grandir, tu dégrades et resynthétises chaque jour l'équivalent de \textbf{200 à 300~g de protéines} — un chantier invisible rendu possible uniquement par la rapidité conférée par les enzymes. Sans elles, les réactions du métabolisme seraient des millions de fois trop lentes pour maintenir la vie : c'est pourquoi on dit que les enzymes sont une condition \emph{sine qua non} de la vie elle-même.
\end{savaistu}

\begin{exercice}[Bilan énergétique d'une journée de jeûne et de réalimentation]
Un élève jeûne depuis le matin ; pourtant, sa température reste à 37~°C et son c{\oe}ur bat.
\begin{enumerate}
\item Quelle branche du métabolisme fournit l'énergie nécessaire au maintien de ces fonctions vitales ?
\item Quelles molécules de réserve sont dégradées en priorité, et par quel type d'enzymes (donnez les noms précis) ?
\item Le soir, après un repas riche en amidon (manioc, maïs), que devient le glucose en excès absorbé par le foie, et cette transformation appartient-elle à l'anabolisme ou au catabolisme ? Justifie ta réponse en termes de sens de la réaction et de bilan énergétique.
\end{enumerate}
\end{exercice}

\begin{corrige}[Exercice : Bilan énergétique]
\begin{enumerate}
\item C'est le \textbf{catabolisme} qui libère l'énergie nécessaire au maintien de la température corporelle (thermogénèse) et aux battements cardiaques (travail mécanique).
\item En priorité, le \textbf{glycogène} hépatique (et musculaire pour le muscle lui-même) est dégradé en glucose-1-phosphate (puis glucose-6-phosphate) par la \textbf{glycogène phosphorylase} (enzyme de la glycogénolyse, une phosphorylase, \textbf{pas une hydrolase}). Le glucose-6-phosphate entre dans la glycolyse puis la respiration cellulaire. En cas de jeûne prolongé, les \textbf{lipides} de réserve (triglycérides du tissu adipeux) sont hydrolysés par des \textbf{lipases} (lipoprotéine lipase, lipase hormono-sensible) en acides gras et glycérol, qui sont oxydés.
\item Le glucose en excès est transformé en \textbf{glycogène} (glycogenèse) dans le foie et les muscles sous l'action de la \textbf{glycogène synthase}. Cette transformation est une \textbf{synthèse} (assemblage de molécules de glucose en un polymère) : elle appartient à l'\textbf{anabolisme}. Elle \textbf{consomme de l'énergie} (hydrolyse d'UTP en UDP pour l'activation du glucose, et consommation indirecte d'ATP).
\end{enumerate}
\end{corrige}

\begin{aretenir}[L'essentiel de la section : Enzymes et renouvellement moléculaire]
\begin{itemize}
\item Le métabolisme = \textbf{anabolisme} (synthèses, consomme de l'énergie/ATP) + \textbf{catabolisme} (dégradations, libère de l'énergie/ATP).
\item Les enzymes de dégradation (amylase, protéases, lipases, enzymes de la glycolyse et respiration, phosphorylase) fournissent énergie (ATP) et petites molécules (briques) ; les enzymes de synthèse (glycogène synthase, ribosome/peptidyl-transférase, synthases des lipides) fabriquent les macromolécules structurales et de réserve.
\item Une voie métabolique est une chaîne de réactions dont chaque étape est catalysée par une enzyme spécifique ; le blocage d'une étape (mutation, inhibiteur) entraîne accumulation du substrat amont et déficit du produit aval (pathologie métabolique).
\item Le renouvellement permanent des protéines (turnover) et des membranes n'est possible que grâce à l'accélération prodigieuse ($10^{6}$ à $10^{12}$ fois) conférée par les enzymes.
\item Une enzyme ne fournit jamais d'énergie : elle accélère la réaction (cinétique) ; l'énergie (thermodynamique) vient des molécules dégradées et transite par l'ATP.
\end{itemize}
\end{aretenir}

\coursec{Déficits enzymatiques et pathologies métaboliques}

Dans la section précédente, nous avons vu que les enzymes assurent en permanence le \cle{renouvellement moléculaire} de nos cellules : elles dégradent les molécules usées et en synthétisent de nouvelles, au sein de véritables chaînes de montage appelées \cle{voies métaboliques}. Mais que se passe-t-il lorsqu'un maillon de cette chaîne vient à manquer ? C'est toute la chaîne qui s'arrête. Cette situation, loin d'être théorique, est à l'origine de nombreuses maladies humaines, certaines fréquentes au Cameroun et en Afrique. C'est ce que nous allons découvrir maintenant.

\courssub{Principe général : une enzyme défaillante bloque toute une voie métabolique}

\textbf{Une situation concrète pour comprendre.} Imagine une chaîne de production de jus de bâton de manioc dans un village : un ouvrier épluche, un autre râpe, un troisième presse. Si l'ouvrier qui râpe est absent, le manioc épluché s'accumule inutilement sur la table, et aucun jus ne sort en bout de chaîne. Une voie métabolique fonctionne exactement ainsi : chaque enzyme transforme un \cle{substrat} en un produit, qui devient le substrat de l'enzyme suivante.

\begin{definition}[Maladie métabolique]
Une \cle{maladie métabolique} (ou maladie héréditaire du métabolisme) est une pathologie due à l'\textbf{absence} ou au \textbf{dysfonctionnement} d'une enzyme d'une voie métabolique. L'enzyme défaillante ne peut plus transformer son substrat : la réaction est bloquée.
\end{definition}

\begin{propriete}[Conséquences du blocage enzymatique]
Le blocage d'une enzyme dans une voie métabolique entraîne deux conséquences directes, liées à la structure en chaîne de la voie :
\begin{itemize}
\item l'\textbf{accumulation du substrat en amont} du blocage : cette substance, parfois toxique à forte dose, s'accumule dans les cellules et le sang ;
\item la \textbf{carence du produit en aval} du blocage : la substance normalement fabriquée vient à manquer, privant l'organisme d'une molécule indispensable.
\end{itemize}
Les symptômes de la maladie résultent de l'une de ces deux anomalies, ou des deux à la fois.
\end{propriete}

\begin{popfigure}
\begin{center}
\begin{tikzpicture}[font=\small, >=Stealth, node distance=0.8cm and 1.2cm]
% Voie normale (haut)
\node[font=\bfseries\small, popgreen] (titreN) at (0,3.0) {Voie normale};
\node[draw, rounded corners, fill=popgreenL, minimum width=1.8cm, minimum height=0.7cm] (S) at (0,2.0) {Substrat $S$};
\node[draw, circle, fill=popblueL, minimum size=0.8cm] (E1) at (3.0,2.0) {$E_1$};
\node[draw, rounded corners, fill=popgreenL, minimum width=1.8cm, minimum height=0.7cm] (I) at (6.0,2.0) {Intermédiaire $I$};
\node[draw, circle, fill=popblueL, minimum size=0.8cm] (E2) at (9.0,2.0) {$E_2$};
\node[draw, rounded corners, fill=popgreenL, minimum width=1.8cm, minimum height=0.7cm] (P) at (12.0,2.0) {Produit $P$};
\draw[->, thick, popgreen] (S) -- (E1);
\draw[->, thick, popgreen] (E1) -- (I);
\draw[->, thick, popgreen] (I) -- (E2);
\draw[->, thick, popgreen] (E2) -- (P);
% Voie bloquée (bas)
\node[font=\bfseries\small, poppink] (titreB) at (0,-1.2) {Voie bloquée ($E_2$ défaillante)};
\node[draw, rounded corners, fill=poporangeL, minimum width=2.5cm, minimum height=0.7cm] (S2) at (1.5,-2.4) {Accumulation de $S$ et $I$};
\node[draw, circle, fill=popblueL, minimum size=0.8cm] (E1b) at (4.5,-2.4) {$E_1$};
\node[draw, circle, fill=poppinkL, minimum size=0.8cm, cross out, draw=poppink, thick] (E2b) at (7.5,-2.4) {$E_2$};
\node[draw, rounded corners, fill=popmuted!20, minimum width=2.2cm, minimum height=0.7cm, dashed] (P2) at (11.5,-2.4) {Absence de $P$};
\draw[->, thick, popgreen] (S2) -- (E1b);
\draw[->, thick, popgreen] (E1b) -- (E2b);
\draw[->, thick, poppink, dashed] (E2b) -- (P2);
% Flèches d'accumulation
\draw[<->, thick, poporange, bend left=20] (S.south) to node[left, poporange, font=\tiny] {Accumulation} (S2.north);
\draw[<->, thick, poporange, bend left=20] (I.south) to node[right, poporange, font=\tiny] {Accumulation} (S2.north east);
\end{tikzpicture}
\end{center}
\legende{Schéma comparé d'une voie métabolique normale (en haut) et d'une voie bloquée par une enzyme $E_2$ défaillante (en bas, barrée). Le substrat $S$ et l'intermédiaire $I$ s'accumulent en amont du blocage (flèches orange) ; le produit $P$ n'est plus fabriqué en aval (flèche pointillée rose).}
\end{popfigure}

\begin{methode}[Relier un symptôme à un déficit enzymatique]
Face à un exercice du Probatoire présentant une pathologie métabolique, procède en quatre étapes :
\begin{enumerate}
\item \textbf{Identifier la voie bloquée} : quelle enzyme est absente ou inactive ?
\item \textbf{Déterminer le substrat accumulé} : c'est souvent lui qui est toxique (il explique les symptômes d'intoxication).
\item \textbf{Déterminer le produit manquant} : sa carence explique les symptômes de déficience.
\item \textbf{Conclure} en reliant chaque symptôme à l'accumulation ou à la carence, puis proposer la prise en charge logique (supprimer le substrat de l'alimentation, ou apporter le produit manquant).
\end{enumerate}
\end{methode}

\courssub{Exemples de déficits enzymatiques}

\courssub{L'intolérance au lactose : un déficit en lactase}

\textbf{Observation.} Dans de nombreuses familles camerounaises, on remarque que les adultes qui boivent du lait frais souffrent de ballonnements, de douleurs abdominales et de diarrhées, alors que les jeunes enfants le digèrent sans problème. Pourquoi ?

Le lait contient un sucre, le \cle{lactose}, qui doit être hydrolysé en glucose et galactose par une enzyme intestinale : la \cle{lactase} (une $\beta$-galactosidase). Chez la plupart des adultes africains (comme chez la majorité des humains adultes dans le monde), l'expression du gène de la lactase diminue fortement après le sevrage : c'est la \cle{non-persistance de la lactase}, un déficit enzymatique physiologique. Le lactose non digéré \textbf{s'accumule} dans la lumière intestinale, où il exerce un effet osmotique (attire l'eau $\rightarrow$ diarrhée) et est fermenté par le microbiote colique, produisant des gaz (ballonnements, douleurs).

\begin{attention}[Intolérance n'est pas allergie]
L'\textbf{intolérance au lactose n'est pas une allergie au lait}. C'est un simple déficit en lactase : il n'y a \textbf{aucune réaction immunitaire}. L'allergie aux protéines de lait de vache (APLV), elle, est une réaction immunitaire dirigée contre les \emph{protéines} du lait (caséines, $\beta$-lactoglobuline), potentiellement grave (choc anaphylactique). Ne confonds jamais ces deux notions dans une copie.
\end{attention}

\courssub{La phénylcétonurie (PCU) : une accumulation toxique évitable}

La \cle{phénylalanine} (Phe) est un acide aminé essentiel apporté par les protéines alimentaires. Normalement, l'enzyme hépatique \cle{phénylalanine-hydroxylase} (PAH), avec son cofacteur la tétrahydrobioptérine, la transforme en \cle{tyrosine} (Tyr). Chez les enfants atteints de \cle{phénylcétonurie} (PCU), cette enzyme est absente ou très peu active : la phénylalanine s'accumule massivement dans le sang (hyperphénylalaninémie) et atteint le cerveau en développement, qu'elle endommage de façon irréversible par toxicité directe et perturbation de la synthèse des neurotransmetteurs, provoquant un retard mental profond. L'excès de phénylalanine est aussi détourné vers une voie mineure produisant des phénylcétones (acide phénylpyruvique) excrétées dans les urines (d'où le nom de la maladie).

\begin{exemplebox}[Un handicap évitable]
Chez un nourrisson atteint de PCU \textbf{non dépisté}, la phénylalanine accumulée détruit progressivement les cellules nerveuses : à quelques années, l'enfant présente un retard mental irréversible, des troubles moteurs et de l'épilepsie. En revanche, si la maladie est détectée dès la naissance (avant 3 semaines), un \textbf{régime pauvre en phénylalanine} (lait maternisé spécial sans phénylalanine, apport contrôlé en protéines naturelles, supplémentation en tyrosine) empêche toute accumulation toxique : l'enfant se développe alors \textbf{normalement}. C'est un exemple remarquable où la connaissance du mécanisme enzymatique permet de prévenir totalement la maladie.
\end{exemplebox}

\begin{savaistu}[Le test de Guthrie : une victoire de la médecine préventive]
Mis au point en 1961 par le médecin américain Robert Guthrie, le \cle{test de Guthrie} consiste à prélever une goutte de sang sur papier buvard au talon du nouveau-né (généralement entre le 3\textsuperscript{e} et le 5\textsuperscript{e} jour de vie), pour doser la phénylalanine par une méthode microbiologique (ou aujourd'hui par spectrométrie de masse en tandem). Généralisé dans de nombreux pays (au Cameroun, le dépistage néonatal s'étend progressivement dans les grandes maternités), ce dépistage systématique a sauvé des \textbf{milliers d'enfants} du handicap mental. Il illustre parfaitement comment une découverte en biochimie enzymatique peut transformer la santé publique.
\end{savaistu}

\courssub{L'albinisme : un déficit en tyrosinase}

La \cle{mélanine} est le pigment brun-noir qui colore la peau, les cheveux et les yeux, et qui nous protège des rayons ultraviolets (UV) du soleil. Elle est fabriquée par les mélanocytes à partir de la tyrosine, grâce à l'enzyme clé \cle{tyrosinase} (qui catalyse l'hydroxylation de la tyrosine en DOPA puis son oxydation en dopaquinone). Chez les personnes atteintes d'albinisme oculocutané de type 1 (le plus fréquent), la tyrosinase est déficiente ou absente : la voie est bloquée, le \textbf{produit manque} (absence de mélanine), d'où une peau et des cheveux très clairs (blancs), une grande sensibilité au soleil (coups de soleil répétés, risque majeur de cancers cutanés précoces : carcinomes, mélanomes) et des troubles de la vision (photophobie, nystagmus, acuité visuelle réduite) dus à l'absence de pigment dans la rétine et l'iris. Ici, contrairement à la PCU, ce n'est pas l'accumulation du substrat (tyrosine) qui est en cause, mais bien la \textbf{carence du produit final}.

\courssub{Le déficit en G6PD (favisme) : une fragilité des globules rouges}

La \cle{glucose-6-phosphate-déshydrogénase} (G6PD) est une enzyme clé de la voie des pentoses phosphate dans les globules rouges. Elle permet la régénération du \cle{NADPH}, indispensable pour maintenir le \cle{glutathion} sous sa forme réduite (GSH), qui protège l'hémoglobine et la membrane érythrocytaire contre l'oxydation. Son déficit, transmis sur le chromosome X (mode \textbf{récessif lié à l'X}), est très répandu en Afrique subsaharienne (fréquence allélique jusqu'à 20-25\% dans certaines zones) car il confère une protection partielle contre le paludisme grave (\emph{Plasmodium falciparum}), ce qui explique sa sélection positive. Lorsque la personne (souvent un homme hémizygote ou une femme homozygote) est exposée à des substances oxydantes (contenues dans les \textbf{fèves} \emph{Vicia faba}, d'où le nom de \cle{favisme}, ou certains médicaments : primaquine, sulfamides, dapsone, quinine à haute dose), ses globules rouges ne peuvent pas neutraliser le stress oxydant : l'hémoglobine précipite (corps de Heinz), la membrane se rigidifie et les globules rouges éclatent massivement dans la rate et la circulation : c'est une \textbf{anémie hémolytique aiguë} (fatigue brutale, pâleur, jaunisse par hyperbilirubinémie, urines foncées « couleur Coca-Cola » par hémoglobinurie). La prise en charge consiste en l'\textbf{éviction stricte} des aliments et médicaments déclencheurs, et parfois une transfusion sanguine en cas d'anémie sévère.

\begin{center}
\begin{tabularx}{\linewidth}{|l|X|X|X|}
\hline
\rowcolor{popblueL} \textbf{Pathologie} & \textbf{Enzyme déficiente} & \textbf{Mécanisme principal \& Conséquences} & \textbf{Prise en charge} \\
\hline
Intolérance au lactose & Lactase ($\beta$-galactosidase) & Accumulation de lactose : effet osmotique + fermentation $\rightarrow$ diarrhée, ballonnements & Régime pauvre en lactose (lait fermenté, fromages affinés), lactase en supplément (comprimés) \\
\hline
Phénylcétonurie (PCU) & Phénylalanine-hydroxylase (PAH) & Accumulation de Phe (neurotoxique) + carence en Tyr $\rightarrow$ retard mental, troubles neurologiques & Dépistage néonatal (test de Guthrie), régime pauvre en Phe (lait spécial, tyrosine supplémentée) à vie \\
\hline
Albinisme (type 1) & Tyrosinase & Absence de mélanine (carence produit) $\rightarrow$ peau/cheveux blancs, photosensibilité extrême, troubles visuels & Protection solaire rigoureuse (vêtements, crème indice 50+), suivi dermatologique et ophtalmologique régulier \\
\hline
Déficit en G6PD (Favisme) & Glucose-6-P-déshydrogénase (G6PD) & Perte protection anti-oxydante $\rightarrow$ hémolyse aiguë déclenchée par fèves/médicaments $\rightarrow$ anémie, ictère, urines foncées & Éviction fèves et médicaments oxydants (liste à remettre au patient), transfusion si anémie sévère \\
\hline
\end{tabularx}
\end{center}

\courssub{Origine génétique des déficits enzymatiques}

Pourquoi une enzyme est-elle absente ou inactive ? Rappelle-toi : une enzyme est une \textbf{protéine}, et toute protéine est fabriquée selon les instructions d'un \cle{gène}. La plupart des déficits enzymatiques sont donc \textbf{héréditaires} : une \cle{mutation} (modification de la séquence d'ADN) du gène codant l'enzyme produit une enzyme anormale (mauvais pliage, site actif altéré, instabilité thermique), inactive, ou empêche totalement sa fabrication (allèle nul). L'enfant reçoit l'allèle muté de ses parents.

\begin{itemize}
\item \textbf{Transmission autosomique récessive} (PCU, Albinisme type 1, la plupart des maladies métaboliques) : il faut recevoir l'allèle muté des \textbf{deux} parents pour être malade. Les parents sont des \textbf{porteurs sains} (hétérozygotes) : ils ont une activité enzymatique résiduelle (souvent 50\%) suffisante pour le métabolisme normal. Deux porteurs ont 25\% de risque d'avoir un enfant malade à chaque grossesse.
\item \textbf{Transmission récessive liée à l'X} (Déficit en G6PD) : le gène est sur le chromosome X. Les \textbf{hommes} (XY) n'ont qu'un seul X : s'ils reçoivent l'allèle muté de leur mère, ils sont \textbf{malades} (hémizygotes). Les \textbf{femmes} (XX) doivent recevoir l'allèle muté des \textbf{deux} parents pour être malades (homozygotes) ; si elles n'en ont qu'un, elles sont porteuses (souvent asymptomatiques ou forme mineure grâce à l'inactivation aléatoire de l'X, mais peuvent avoir des crises hémolytiques).
\end{itemize}
C'est pourquoi deux parents en bonne santé (porteurs) peuvent avoir un enfant atteint de PCU ou d'albinisme, et pourquoi le favisme touche predominantly les hommes.

\courssub{Dépistage, prévention et inhibiteurs enzymatiques}

La compréhension des mécanismes enzymatiques offre trois stratégies de lutte contre ces maladies :
\begin{itemize}
\item le \textbf{dépistage précoce} : test de Guthrie à la naissance pour la PCU (et maintenant d'autres maladies : hypothyroïdie, mucoviscidose, drépanocytose au Cameroun), recherche du déficit en G6PD avant de prescrire certains antipaludéens (primaquine) ou sulfamides ;
\item les \textbf{régimes alimentaires adaptés} : supprimer ou limiter le substrat qui s'accumule (phénylalanine, lactose, fèves) ;
\item la \textbf{supplémentation} : apporter l'enzyme manquante (lactase en comprimés avant le repas) ou le produit manquant (tyrosine dans la PCU).
\end{itemize}

\begin{definition}[Inhibiteur enzymatique]
Un \cle{inhibiteur enzymatique} est une substance qui se fixe sur une enzyme (souvent sur son \textbf{site actif}) et bloque ou ralentit son activité. L'inhibition est dite \textbf{compétitive} si l'inhibiteur entre en compétition avec le substrat pour le site actif (structure ressemblant au substrat) ; elle est \textbf{non compétitive} si l'inhibiteur se fixe sur un site distinct (site allostérique), modifiant la conformation de l'enzyme.
\end{definition}

Cette notion est capitale en pharmacologie : de nombreux \textbf{médicaments} sont des inhibiteurs enzymatiques. Par exemple, l'aspirine inhibe irréversiblement les cyclooxygénases (COX), enzymes de la synthèse des prostaglandines de l'inflammation ; les inhibiteurs de la protéase du VIH bloquent la maturation virale ; les statines inhibent l'HMG-CoA réductase (cholestérol). À l'inverse, de nombreux \textbf{poisons} agissent aussi en bloquant des enzymes vitales : c'est le cas des insecticides organophosphorés (parathion, malathion) qui inhibent l'acétylcholinestérase (synapses cholinergiques), provoquant une crise cholinergique mortelle ; ou du cyanure qui inhibe le cytochrome c oxydase (complexe IV de la chaîne respiratoire), bloquant la respiration cellulaire.

\begin{exoresolu}[Exploiter un document : la phénylcétonurie]
\textbf{Énoncé.} Chez un nouveau-né, le dosage sanguin réalisé trois jours après la naissance donne un taux de phénylalanine de \SI{1.8}{mmol/L}, alors que la normale est inférieure à \SI{0.12}{mmol/L}. Le taux de tyrosine est, lui, anormalement bas. 1) Quelle enzyme est vraisemblablement déficiente ? 2) Explique les deux anomalies observées. 3) Propose une mesure préventive.
\tcblower
\textbf{Solution détaillée.}
1) L'enzyme déficiente est la \textbf{phénylalanine-hydroxylase} (PAH), qui transforme normalement la phénylalanine en tyrosine en présence de tétrahydrobioptérine : il s'agit d'une phénylcétonurie (PCU) classique.

2) Le blocage de la voie métabolique a deux conséquences directes conformes au principe général : le \textbf{substrat} (phénylalanine) s'\textbf{accumule} en amont du blocage, d'où son taux sanguin 15 fois supérieur à la normale (\SI{1.8}{mmol/L} vs $<\SI{0.12}{mmol/L}$) ; le \textbf{produit} (tyrosine) n'est plus fabriqué en aval, d'où son taux anormalement bas (carence en tyrosine, qui devient un acide aminé essentiel chez ces patients). La phénylalanine accumulée est neurotoxique pour le cerveau en développement (myélinisation, synthèse neurotransmetteurs).

3) Il faut mettre en place dès la naissance (avant 3 semaines de vie) un \textbf{régime pauvre en phénylalanine} (lait maternisé spécial sans phénylalanine, alimentation contrôlée en protéines naturelles, supplémentation en tyrosine) afin d'empêcher l'accumulation du substrat toxique : le développement cérébral de l'enfant sera alors normal. Le régime doit être poursuivi au moins jusqu'à l'adolescence, voire à vie.
\end{exoresolu}

\begin{aretenir}[L'essentiel de la section]
\begin{itemize}
\item Une maladie métabolique résulte de l'absence ou du dysfonctionnement d'une enzyme : la voie est bloquée, le \textbf{substrat s'accumule} (souvent toxique) et le \textbf{produit manque}.
\item Quatre exemples types à connaître pour le Probatoire :
  \begin{itemize}
  \item Intolérance au lactose : déficit en \textbf{lactase} $\rightarrow$ accumulation lactose (digestif).
  \item Phénylcétonurie (PCU) : déficit en \textbf{phénylalanine-hydroxylase} $\rightarrow$ accumulation Phe (neurotoxique) + carence Tyr.
  \item Albinisme : déficit en \textbf{tyrosinase} $\rightarrow$ carence mélanine (protection UV, vision).
  \item Favisme (G6PD) : déficit en \textbf{G6PD} (lié à l'X) $\rightarrow$ perte protection anti-oxydante $\rightarrow$ hémolyse aiguë déclenchée (fèves, médicaments).
  \end{itemize}
\item Ces déficits sont le plus souvent \textbf{héréditaires} (mutation du gène codant l'enzyme) : transmission autosomique récessive (PCU, albinisme) ou récessive liée à l'X (G6PD).
\item La prévention repose sur le dépistage néonatal (test de Guthrie pour PCU), les régimes adaptés (exclusion substrat) et la supplémentation (enzyme ou produit).
\item Les inhibiteurs enzymatiques (compétitifs ou non) bloquent l'activité enzymatique : principe d'action de nombreux médicaments (aspirine, antirétroviraux, statines) et de poisons (cyanure, organophosphorés).
\end{itemize}
\end{aretenir}

\begin{exercice}[Pour s'entraîner]
Une personne d'origine africaine présente, quelques heures après avoir mangé un plat de fèves, une fatigue intense, une jaunisse et des urines foncées couleur « Coca-Cola ». L'analyse sanguine montre une anémie brutale avec réticulocytose et la présence de corps de Heinz dans les globules rouges. 1) Quel déficit enzymatique suspectes-tu ? 2) Explique le mécanisme biophysique reliant l'ingestion de fèves à l'hémolyse. 3) Quels conseils précis donnerais-tu à cette personne pour l'avenir ?
\end{exercice}

\begin{corrige}[Exercice : le favisme]
1) On suspecte un \textbf{déficit en glucose-6-phosphate-déshydrogénase (G6PD)}, appelé favisme. La survenue brutale après ingestion de fèves, le sexe masculin (fréquent), l'origine africaine (zone palustre), l'anémie hémolytique avec corps de Heinz (dénaturation de l'hémoglobine par oxydation) sont pathognomoniques.

2) \textbf{Mécanisme :} La G6PD est l'enzyme limitante de la voie des pentoses phosphate dans l'érythrocyte. Elle produit le \textbf{NADPH} nécessaire à la réduction du glutathion oxydé (GSSG) en glutathion réduit (GSH) via la glutathion réductase. Le GSH est le principal piégeur des espèces réactives de l'oxygène (ERO). En l'absence de G6PD, le stock de NADPH s'effondre, le GSH chute, et les ERO (produites notamment par les composés oxydants de la fève : vicine, convicine, isouramil) oxydent les groupements sulfhydryles de l'hémoglobine (formation de corps de Heinz) et les protéines de la membrane (rigidification). Les globules rouges deviennent rigides, sont piégés et détruits dans la rate (hémolyse extravasculaire) et lyse intravasculaire (hémoglobinurie $\rightarrow$ urines foncées, bilirubine libre $\rightarrow$ ictère).

3) \textbf{Conseils :}
\begin{itemize}
\item \textbf{Éviction stricte et définitive des fèves} (fraîches, sèches, farine) et de l'inhalation de leur pollen.
\item \textbf{Informer tout personnel soignant} de son déficit G6PD avant toute prescription : éviter les médicaments oxydants formels (primaquine, dapsone, sulfamides, phénazopyridine, quinine haute dose, rasburicase) et utiliser avec prudence ceux à risque discuté (aspirine haute dose, vitamine K, chloroquine -- bien que souvent tolérée aux doses antipaludéennes standard).
\item Porter une \textbf{carte d'alerte médicale} (type carte de groupe sanguin) mentionnant le déficit.
\item En cas de crise hémolytique : consulter en urgence (transfusion possible si Hb $< 7$ g/dL ou signes de décompensation), hydratation alcaline pour protéger les reins de l'hémoglobine.
\item Conseil génétique : ses frères et sœurs (surtout les garçons) ont un risque d'être atteints ; ses filles seront obligatoirement porteuses ; ses fils ne le seront pas (père transmet son Y).
\end{itemize}
\end{corrige}

\newpage

\coursec{Activité d'intégration}

\courssub{Objectifs de l'activité}

À l'issue de cette activité d'intégration, tu seras capable de :

\begin{itemize}
\item \cle{exploiter} une courbe d'activité enzymatique en fonction de la température pour déterminer graphiquement une \cle{température optimale} ;
\item \cle{exploiter} un tableau de données pour déterminer un \cle{pH optimal} ;
\item \cle{interpréter} un échec expérimental en invoquant la \cle{dénaturation} irréversible des enzymes par la chaleur ;
\item \cle{mobiliser} la notion de \cle{spécificité enzymatique} (enzyme--substrat) pour expliquer le rôle de la lactase ;
\item \cle{relier} un déficit enzymatique (en lactase) à une pathologie métabolique (l'intolérance au lactose) et proposer une solution concrète ;
\item \cle{concevoir} un protocole argumenté de fabrication optimale d'un aliment fermenté.
\end{itemize}

\begin{aretenir}[Compétence visée]
Intégrer les savoirs et savoir-faire de la leçon (catalyse enzymatique, spécificité, facteurs de l'activité enzymatique, déficits et pathologies) pour résoudre un problème authentique de conservation et de transformation alimentaire.
\end{aretenir}

\courssub{Situation}

\textbf{Le mystère du lait caillé : enquête enzymatique au marché de Mokolo.}

Au marché de Mokolo, à Yaoundé, Madame Aïcha vend chaque matin du \emph{kossam} (lait caillé) très apprécié des habitants du quartier. Pour le fabriquer, elle ajoute au lait de vache frais un peu de \cle{ferment lactique} (une petite quantité de kossam de la veille, riche en bactéries lactiques dont les enzymes transforment le lactose en acide lactique, ce qui fait cailler le lait). Depuis quelques semaines, sa production est devenue capricieuse :

\begin{itemize}
\item lorsqu'elle laisse fermenter le lait à \SI{37}{\celsius}, le kossam est onctueux et réussi en 6 heures ;
\item un jour, pressée, elle a fait \textbf{bouillir} le lait (\SI{100}{\celsius}) puis a ajouté le ferment \emph{immédiatement}, sans laisser refroidir : le lait n'a \textbf{jamais caillé} ;
\item un après-midi de forte chaleur, la cuve est restée au soleil à \SI{45}{\celsius} : la fermentation a été beaucoup trop rapide et le kossam est devenu acide et grumeleux, invendable.
\end{itemize}

Par ailleurs, un client fidèle, M. Etoundi, se plaint de ballonnements, de crampes et de diarrhées après avoir consommé du lait frais. Son médecin au Centre Médical d'Arrondissement lui a diagnostiqué une \cle{intolérance au lactose} : son intestin grêle ne produit presque plus de \cle{lactase}, l'enzyme qui hydrolyse le lactose en glucose et en galactose.

Madame Aïcha, intriguée, sollicite les élèves de Première D du lycée voisin. Elle leur fournit les documents suivants.

\textbf{Document 1 : activité des enzymes du ferment lactique en fonction de la température} (activité mesurée par la vitesse de production d'acide lactique, en unités arbitraires, u.a.).

\begin{popfigure}
\begin{tikzpicture}
\begin{axis}[
width=0.85\linewidth, height=6.5cm,
xlabel={Température (\si{\celsius})},
ylabel={Activité enzymatique (u.a.)},
xmin=0, xmax=110, ymin=0, ymax=105,
xtick={0,10,20,30,37,45,60,80,100},
ytick={0,20,40,60,80,100},
grid=both, grid style={popline!40},
legend style={at={(0.03,0.97)}, anchor=north west, font=\small}
]
\addplot[popcurve, domain=5:100, samples=200]
{100*exp(-((x-37)/22)^2) * (1/(1+exp((x-70)/3)))};
\addlegendentry{Enzymes du ferment lactique}
\addplot[poporange, very thick, dashed, domain=5:100, samples=200]
{95*exp(-((x-37)/25)^2) * (1/(1+exp((x-55)/3)))};
\addlegendentry{Lactase (enzyme digestive humaine)}
\draw[popgreen, thick, dashed] (axis cs:37,0) -- (axis cs:37,100);
\node[popgreen, font=\small, anchor=west] at (axis cs:38,95) {optimum : \SI{37}{\celsius}};
\end{axis}
\end{tikzpicture}
\legende{Courbes d'activité enzymatique en fonction de la température. L'activité est nulle au-delà de \SI{60}{\celsius} environ (dénaturation thermique irréversible).}
\end{popfigure}

\textbf{Document 2 : activité des enzymes du ferment en fonction du pH} (mesurée à \SI{37}{\celsius}).

\begin{center}
\begin{tabularx}{\linewidth}{|l|*{8}{>{\centering\arraybackslash}X|}}
\hline
\rowcolor{popblueL}
\textbf{pH} & 3 & 4 & 4,5 & 5 & 5,5 & 6 & 7 & 8 \\
\hline
\textbf{Activité (u.a.)} & 5 & 35 & 70 & 95 & 80 & 45 & 15 & 5 \\
\hline
\end{tabularx}
\legende{Activité des enzymes lactiques en fonction du pH du milieu.}
\end{center}

\textbf{Document 3 : teneur en lactose de quelques produits laitiers} (pour 100 g).

\begin{center}
\begin{tabularx}{\linewidth}{|l|X|}
\hline
\rowcolor{popblueL}
\textbf{Produit} & \textbf{Lactose (g/100 g)} \\
\hline
Lait frais de vache & 4,8 \\
\hline
Kossam (lait caillé fermenté 6 h à \SI{37}{\celsius}) & 3,0 \\
\hline
Kossam très fermenté (24 h) & 1,5 \\
\hline
\end{tabularx}
\legende{La fermentation lactique consomme une partie du lactose : plus la fermentation est longue, moins le produit contient de lactose.}
\end{center}

\courssub{Consignes}

Par groupes de 4 à 6 élèves, vous produirez une \textbf{affiche} comportant les courbes annotées, vos calculs et vos conclusions argumentées. Répondez aux questions suivantes.

\begin{enumerate}
\item À l'aide du document 1, déterminer graphiquement la \cle{température optimale} des enzymes du ferment lactique. Justifier alors pourquoi la production de Madame Aïcha réussit à \SI{37}{\celsius}.
\item Expliquer, en vous appuyant sur le document 1 et sur vos connaissances, pourquoi le lait n'a pas caillé lorsque le ferment a été ajouté juste après ébullition. Nommer précisément le phénomène subi par les enzymes et préciser s'il est réversible.
\item Le document 1 montre qu'à \SI{45}{\celsius} l'activité est encore élevée. Expliquer pourquoi la fermentation y est trop rapide et pourquoi le kossam devient acide et grumeleux. Quel phénomène menace les enzymes si la température continue d'augmenter ?
\item À l'aide du document 2, déterminer le \cle{pH optimal} des enzymes du ferment. Expliquer pourquoi l'acidification progressive du lait finit par ralentir elle-même la fermentation (on parle d'auto-limitation).
\item Rappeler ce qu'est la \cle{spécificité enzymatique} et l'illustrer avec le couple lactase--lactose. Pourquoi la lactase ne peut-elle pas digérer l'amidon ?
\item Expliquer l'origine des troubles de M. Etoundi en reliant le \cle{déficit enzymatique} à la \cle{pathologie}. Que devient le lactose non digéré dans le gros intestin ?
\item En utilisant le document 3, proposer à M. Etoundi un conseil argumenté lui permettant de continuer à consommer un produit laitier.
\item Rédiger le \textbf{protocole optimal de fabrication du kossam} (température, hygiène, durée, conservation), en justifiant chaque étape par un savoir de la leçon.
\end{enumerate}

\courssub{Ressources à mobiliser}

\begin{aretenir}[Ce qu'il faut avoir en tête]
\begin{itemize}
\item Une \cle{enzyme} est un \textbf{catalyseur biologique} de nature protéique : elle accélère une réaction sans être consommée, en abaissant l'énergie d'activation.
\item Le \cle{substrat} se fixe sur le \cle{site actif} de l'enzyme ; cette reconnaissance est \textbf{spécifique} (modèle clé--serrure, ou ajustement induit).
\item L'activité enzymatique dépend de la \textbf{température} : elle croît jusqu'à un \cle{optimum} (souvent voisin de \SI{37}{\celsius} chez l'Homme), puis s'effondre par \cle{dénaturation} (destruction irréversible de la structure spatiale de la protéine) au-delà d'environ \SI{60}{\celsius}.
\item Chaque enzyme possède un \cle{pH optimal} ; un écart de pH modifie le site actif et diminue l'activité.
\item Un \cle{déficit enzymatique} (ex. : lactase) entraîne l'accumulation du substrat non transformé et une \textbf{pathologie métabolique} (ex. : intolérance au lactose).
\item Lecture graphique : l'optimum correspond au \textbf{maximum} de la courbe ; on projette le sommet sur l'axe des abscisses.
\end{itemize}
\end{aretenir}

\courssub{Démarche et corrigé commenté}

\begin{exoresolu}[Corrigé commenté de l'activité]
\textbf{Questions 1 à 4 : exploitation des documents.}

\textbf{Question 2 :} déterminer la température optimale des enzymes du ferment et justifier le succès à \SI{37}{\celsius}.

\textbf{Question 4 :} expliquer l'échec après ébullition.

\textbf{Question 6 :} déterminer le pH optimal et expliquer l'auto-limitation.

\textbf{Questions 5 et 6 (savoirs) :} spécificité enzymatique et lien déficit--pathologie.

\textbf{Questions 7 et 8 :} conseil au client et protocole optimal.

\tcblower
\textbf{Question 1.} Sur le document 1, la courbe des enzymes du ferment présente un \textbf{maximum} dont l'abscisse est \SI{37}{\celsius} (activité voisine de 100 u.a.). La température optimale est donc \SI{37}{\celsius}. À cette température, l'agitation thermique est suffisante pour des chocs fréquents enzyme--substrat, sans dénaturer les protéines : la vitesse de production d'acide lactique est maximale, d'où un caillage rapide et régulier. C'est exactement la condition empirique utilisée par Madame Aïcha.

\textbf{Question 2.} Le document 1 montre qu'au-delà d'environ \SI{60}{\celsius}, l'activité devient \textbf{nulle}. Or le ferment a été ajouté dans un lait proche de \SI{100}{\celsius}. À cette température, les vibrations thermiques rompent les liaisons faibles qui maintiennent la structure tridimensionnelle des enzymes : c'est la \cle{dénaturation thermique}. Le site actif est déformé, il ne reconnaît plus le lactose. Cette dénaturation est \textbf{irréversible} : même en refroidissant ensuite, les enzymes ne retrouvent pas leur forme active, et le lait ne peut plus cailler.

\textbf{Question 3.} À \SI{45}{\celsius}, la courbe indique une activité encore élevée (environ 85--90 u.a.) : la fermentation démarre très vite. La production massive d'acide lactique fait chuter brutalement le pH ; or le document 2 montre que l'activité s'effondre aux pH acides, et surtout l'acidité excessive fait précipiter les protéines du lait de façon désordonnée : le kossam devient acide et grumeleux. De plus, à \SI{45}{\celsius} on s'approche de la zone de dénaturation : si la température monte encore, les enzymes seront détruites.

\textbf{Question 4.} Le document 2 présente un maximum d'activité (95 u.a.) à \textbf{pH = 5} : c'est le pH optimal. Au cours de la fermentation, l'acide lactique produit abaisse progressivement le pH du lait ; en passant sous pH 5 puis vers pH 4, l'activité tombe à 35 u.a. : la fermentation se ralentit d'elle-même. Cette \textbf{auto-limitation} est favorable : elle empêche une acidification totale et contribue à la conservation du kossam.

\textbf{Question 5.} La \cle{spécificité enzymatique} signifie qu'une enzyme ne catalyse qu'une réaction précise, sur un substrat précis, car la forme du site actif est complémentaire de celle du substrat (modèle clé--serrure). La \cle{lactase} hydrolyse uniquement le \cle{lactose} en glucose et galactose ; elle ne peut pas hydrolyser l'amidon, dont la géométrie moléculaire ne correspond pas à son site actif (l'amidon est dégradé par les \emph{amylases}).

\textbf{Question 6.} M. Etoundi présente un \cle{déficit en lactase} : le lactose du lait frais n'est pas hydrolysé dans l'intestin grêle. Il arrive intact dans le gros intestin, où les bactéries le fermentent en produisant des gaz (ballonnements, crampes) et des acides qui appellent l'eau dans le tube digestif (diarrhées). Le déficit enzymatique est donc bien la \textbf{cause directe} de la pathologie.

\textbf{Question 7.} Le document 3 montre que le kossam fermenté 24 h ne contient plus que 1,5 g de lactose pour 100 g, contre 4,8 g pour le lait frais, car les enzymes du ferment ont consommé une grande partie du lactose. Conseil : M. Etoundi peut consommer du \textbf{kossam bien fermenté} (ou du lait pré-traité à la lactase), en quantités modérées, plutôt que du lait frais.

\textbf{Question 8. Protocole optimal.}
\begin{enumerate}
\item Faire bouillir le lait pour l'hygiène (destruction des microbes contaminants), \textbf{puis le laisser refroidir à environ \SI{37}{\celsius}} avant d'ajouter le ferment (pour ne pas dénaturer les enzymes).
\item Ajouter le ferment, mélanger, maintenir la cuve à \SI{37}{\celsius} (à l'ombre, dans un récipient couvert) pendant environ 6 h.
\item Surveiller le caillage ; dès que le lait est pris, réfrigérer : le froid (\SI{4}{\celsius}) \textbf{ralentit} l'activité enzymatique sans dénaturer les enzymes, ce qui stoppe l'acidification et conserve le kossam.
\end{enumerate}
Chaque étape repose sur un savoir : température optimale, dénaturation irréversible, effet réversible du froid, auto-limitation par le pH.
\end{exoresolu}

\courssub{Bilan de l'activité}

Cette enquête au marché de Mokolo a permis de mettre en évidence, sur un problème concret de la vie quotidienne camerounaise, l'ensemble des acquis de la leçon :

\begin{itemize}
\item l'activité enzymatique possède une \cle{température optimale} (\SI{37}{\celsius} ici) et un \cle{pH optimal} (pH 5 ici), lisibles sur des documents expérimentaux ;
\item la \cle{dénaturation} par la chaleur est \textbf{irréversible}, alors que le froid ne fait que \textbf{ralentir} l'activité de façon \textbf{réversible} — distinction capitale, souvent source d'erreur au Probatoire ;
\item la \cle{spécificité} enzyme--substrat explique à la fois le fonctionnement du ferment et le déficit en lactase ;
\item un \cle{déficit enzymatique} se traduit par une pathologie métabolique, et la connaissance des enzymes permet de proposer des solutions concrètes (aliment fermenté appauvri en lactose, protocole de fabrication maîtrisé).
\end{itemize}

\begin{attention}[Erreurs fréquentes à éviter]
\begin{itemize}
\item Dire que la chaleur « tue » les enzymes : une enzyme n'est pas un être vivant, on dit qu'elle est \textbf{dénaturée}.
\item Croire que le froid dénature les enzymes : il ne fait que \textbf{ralentir} leur activité (c'est le principe du réfrigérateur).
\item Confondre optimum de température et optimum de pH : ce sont deux paramètres indépendants, chacun lu sur son propre document.
\item Affirmer qu'une enzyme est « consommée » par la réaction : un catalyseur est \textbf{retrouvé intact} en fin de réaction.
\end{itemize}
\end{attention}

\newpage

\coursec{Méthodes et exercices résolus}

\coursec{Leçon ND04 : Sensibilisation sur l'influence des enzymes sur les réactions chimiques indispensables au renouvellement moléculaire}

\courssub{Notion d'enzyme, de substrat et de catalyse enzymatique}

\begin{definition}[Enzyme]
Une \cle{enzyme} est une macromolécule biologique, le plus souvent une \cle{protéine} (parfois un ARN catalytique ou \cle{ribozyme}), qui agit comme un \cle{catalyseur} des réactions chimiques du métabolisme cellulaire. Elle accélère la vitesse d'une réaction sans être consommée ni modifiée de façon permanente à l'issue de celle-ci.
\end{definition}

\begin{definition}[Substrat et produit]
Le \cle{substrat} est la molécule (ou les molécules) sur laquelle l'enzyme agit. Au cours de la réaction, le substrat est transformé en \cle{produit(s)}. L'enzyme se libère intacte et peut catalyser un nouveau cycle.
\end{definition}

\begin{propriete}[Propriétés fondamentales des enzymes]
\begin{enumerate}
\item \textbf{Nature protéique (généralement) :} structure primaire, secondaire, tertiaire (et parfois quaternaire) ; la conformation spatiale 3D est indispensable à l'activité.
\item \textbf{Catalyse :} l'enzyme \cle{abaisse l'énergie d'activation} ($E_a$) de la réaction, augmentant ainsi la vitesse de réaction (souvent d'un facteur $10^6$ à $10^{12}$) sans modifier l'équilibre thermodynamique ($\Delta G$ inchangé).
\item \textbf{Spécificité :} une enzyme donnée ne catalyse qu'une seule réaction (ou un groupe de réactions très voisines) sur un substrat précis (ou un groupe de substrats structuralement proches).
\item \textbf{Réversibilité :} la plupart des réactions enzymatiques sont réversibles ; le sens net dépend des concentrations en substrats et produits (loi de l'action de masse).
\item \textbf{Sensibilité aux conditions physico-chimiques :} l'activité dépend fortement de la température, du pH, de la force ionique et de la présence de \cle{cofactors} (ions métalliques \ce{Mg^{2+}}, \ce{Zn^{2+}}, \ce{Fe^{2+/3+}}) ou de \cle{coenzymes} (dérivés de vitamines : NAD$^+$, FAD, CoA).
\end{enumerate}
\end{propriete}

\begin{experience}[Mise en évidence qualitative de l'activité catalytique : la catalase]
\textbf{Matériel :} eau oxygénée (\ce{H2O2} \SI{10}{\%} volumes), foie frais de bœuf (source de catalase), foie bouilli 10 min (témoin dénaturé), tubes à essai, bâtonnets de bois incandescents.
\textbf{Protocole :}
\begin{enumerate}
\item Tube A : \SI{5}{mL} \ce{H2O2} + \SI{1}{cm^3} foie frais.
\item Tube B : \SI{5}{mL} \ce{H2O2} + \SI{1}{cm^3} foie bouilli.
\item Tube C (témoin) : \SI{5}{mL} \ce{H2O2} seul.
\item On approche un bâtonnet incandescent de l'embouchure de chaque tube.
\end{enumerate}
\textbf{Observations :}
\begin{itemize}
\item Tube A : moussage intense (dégagement gazeux) ; le bâtonnet se rallume violemment $\rightarrow$ \ce{O2}.
\item Tube B : pas de moussage, pas de rallumage.
\item Tube C : pas de réaction visible.
\end{itemize}
\textbf{Interprétation :} La catalase du foie frais catalyse la décomposition de l'eau oxygénée : $\ce{2 H2O2 ->[catalase] 2 H2O + O2}$. L'ébullition a dénaturé l'enzyme (destruction de la structure 3D), abolissant l'activité. L'eau oxygénée est stable spontanément à température ambiante.
\legende{Expérience classique de mise en évidence de la catalase (foie frais vs foie bouilli).}
\end{experience}

\begin{exemplebox}[Modèles d'interaction Enzyme--Substrat]
\begin{itemize}
\item \textbf{Modèle « Clé--Serrure » (Emil Fischer, 1894) :} le site actif possède une forme rigide, géométriquement complémentaire à celle du substrat. Le substrat s'y insère comme une clé dans une serrure.
\item \textbf{Modèle « Ajustement Induit » (Daniel Koshland, 1958) :} le site actif est flexible ; la fixation du substrat induit un changement de conformation de l'enzyme qui « embrasse » le substrat et place les groupes catalytiques en position optimale. Ce modèle rend compte de la catalyse et de la régulation allostérique.
\end{itemize}
Dans les deux cas, se forme un \cle{complexe enzyme--substrat} (ES) transitoire.
\legende{Schémas comparatifs : (a) Clé--Serrure : enzyme rigide. (b) Ajustement Induit : enzyme flexible se déformant au contact du substrat.}
\end{exemplebox}

\begin{aretenir}[Vocabulaire clé de la catalyse enzymatique]
\begin{itemize}
\item \textbf{Site actif :} région creuse de l'enzyme (quelques acides aminés) où se lie le substrat et où a lieu la catalyse.
\item \textbf{Complexe enzyme--substrat (ES) :} association transitoire enzyme + substrat.
\item \textbf{Energie d'activation ($E_a$) :} barrière énergétique à franchir pour que la réaction démarre. L'enzyme la diminue en stabilisant l'état de transition.
\item \textbf{Vitesse de réaction ($v$) :} quantité de produit formé (ou substrat disparu) par unité de temps. Elle est proportionnelle à la concentration en complexe ES.
\end{itemize}
\end{aretenir}

\courssub{Spécificité enzymatique et site actif}

\begin{definition}[Spécificité enzymatique]
La \cle{spécificité} d'une enzyme se manifeste à deux niveaux :
\begin{itemize}
\item \textbf{Spécificité de substrat (ou de liaison) :} l'enzyme ne reconnaît qu'un substrat précis (ex. : uréase $\rightarrow$ urée) ou un groupe de substrats très voisins (ex. : hexokinase $\rightarrow$ glucose, fructose, mannose).
\item \textbf{Spécificité d'action (ou de réaction) :} l'enzyme ne catalyse qu'un seul type de transformation chimique (hydrolyse, oxydo-réduction, transfert de groupe, etc.).
\end{itemize}
\end{definition}

\begin{propriete}[Base moléculaire de la spécificité]
La spécificité résulte de la \textbf{complémentarité stéréochimique et chimique} entre le site actif et le substrat :
\begin{itemize}
\item Forme géométrique (emboîtement).
\item Répartition des charges (interactions ioniques, liaisons hydrogène, forces de Van der Waals, interactions hydrophobes).
\item Présence de résidus catalytiques (acides/bases, nucléophiles) précisément positionnés.
\end{itemize}
Toute modification de la structure de l'enzyme (mutation, dénaturation, fixation d'un inhibiteur) altère cette complémentarité et donc l'activité.
\end{propriete}

\courssub{Exploiter des courbes d'activité enzymatique en fonction de la température}

\begin{methode}[Lire et interpréter une courbe activité = f(température)]
\begin{enumerate}
\item \textbf{Identifier les axes} : l'axe des abscisses porte la température (en \si{\celsius}), l'axe des ordonnées porte l'activité enzymatique (vitesse de réaction, quantité de produit formé par unité de temps, ou absorbance). Toujours citer les unités.
\item \textbf{Décrire la courbe en deux phases} : d'abord une \cle{augmentation} de l'activité avec la température (l'agitation moléculaire accroît les chocs efficaces enzyme--substrat), puis une \cle{diminution} brutale au-delà d'un maximum.
\item \textbf{Repérer la température optimale} : c'est l'abscisse du maximum de la courbe (souvent voisine de \SI{37}{\celsius} pour les enzymes humaines/mammifères). La nommer précisément : « température optimale ».
\item \textbf{Interpréter les deux phases} :
\begin{itemize}
\item Phase montante : augmentation de l'énergie cinétique moyenne des molécules $\rightarrow$ augmentation de la fréquence des chocs et de la proportion de molécules ayant une énergie $\ge E_a$.
\item Phase descendante : \cle{dénaturation} de l'enzyme (rupture des liaisons faibles : liaisons H, ponts ioniques, interactions hydrophobes $\rightarrow$ perte de la structure tertiaire/quaternaire $\rightarrow$ déformation du site actif). Phénomène \emph{irréversible} à haute température.
\end{itemize}
\item \textbf{Conclure} en une phrase reliant structure et fonction : « L'activité enzymatique est maximale à la température optimale ; au-delà, l'enzyme se dénature et perd son activité de façon irréversible. »
\end{enumerate}
\textbf{Réflexes Probatoire} : ne jamais écrire « l'enzyme meurt » (une enzyme n'est pas vivante) mais « l'enzyme est dénaturée » ; distinguer \emph{ralentissement} (basse température, réversible) et \emph{dénaturation} (haute température, irréversible).
\end{methode}

\begin{exoresolu}[Activité de la catalase en fonction de la température]
On mesure le volume de dioxygène dégagé (en mL) en 2 minutes lors de la décomposition de l'eau oxygénée par la catalase d'un extrait de foie de bœuf, à différentes températures :

\begin{center}
\begin{tabularx}{\linewidth}{|l|X|X|X|X|X|X|}\hline
\rowcolor{popblueL}
Température (\si{\celsius}) & 0 & 10 & 20 & 37 & 50 & 70 \\ \hline
Volume de \ce{O2} (mL) & 1 & 4 & 9 & 15 & 8 & 0,5 \\ \hline
\end{tabularx}
\end{center}

\textbf{1.} Trace l'allure de la courbe de l'activité enzymatique en fonction de la température.\\
\textbf{2.} Détermine la température optimale de la catalase.\\
\textbf{3.} Interprète les résultats obtenus à \SI{0}{\celsius} et à \SI{70}{\celsius}.
\tcblower
\textbf{1. Tracé de la courbe.} On porte la température en abscisse et le volume de \ce{O2} (proportionnel à l'activité enzymatique) en ordonnée. La courbe monte de \SI{0}{\celsius} à \SI{37}{\celsius}, passe par un maximum, puis redescend fortement jusqu'à presque zéro à \SI{70}{\celsius} : c'est une courbe en cloche asymétrique.

\begin{popfigure}
\begin{tikzpicture}
\begin{axis}[width=0.85\linewidth,height=6cm,
xlabel={Température (\si{\celsius})},ylabel={Volume de \ce{O2} (mL)},
xmin=0,xmax=75,ymin=0,ymax=17,
xtick={0,10,20,37,50,70},grid=both,grid style={popline!40}]
\addplot[popcurve,mark=*,mark size=1.5pt] coordinates {(0,1) (10,4) (20,9) (37,15) (50,8) (70,0.5)};
\addplot[dashed,poporange] coordinates {(37,0) (37,15)};
\end{axis}
\end{tikzpicture}
\legende{Activité de la catalase en fonction de la température. Le maximum (température optimale) est atteint vers \SI{37}{\celsius}.}
\end{popfigure}

\textbf{2. Température optimale.} Le volume de \ce{O2} est maximal (\SI{15}{mL}) à $T = \SI{37}{\celsius}$. La température optimale de la catalase est donc \SI{37}{\celsius}, ce qui correspond à la température corporelle des mammifères : l'enzyme est adaptée à son milieu de fonctionnement.

\textbf{3. Interprétation.}
\begin{itemize}
\item À \SI{0}{\celsius}, l'activité est très faible (\SI{1}{mL}) car l'agitation thermique est réduite : les chocs entre l'enzyme et son substrat (l'eau oxygénée) sont rares et peu énergétiques. L'enzyme n'est \textbf{pas détruite} : si l'on réchauffe le milieu, l'activité reprend. C'est un \cle{ralentissement réversible}.
\item À \SI{70}{\celsius}, l'activité est quasi nulle (\SI{0,5}{mL}) car la chaleur a rompu les liaisons faibles qui maintiennent la structure spatiale (repliement) de la protéine enzymatique : le site actif est déformé et ne reconnaît plus le substrat. C'est une \cle{dénaturation irréversible}.
\end{itemize}
\textbf{Conclusion} : l'activité de la catalase est maximale à sa température optimale (\SI{37}{\celsius}) ; le froid la ralentit de façon réversible, la chaleur excessive la dénature de façon irréversible.
\end{exoresolu}

\courssub{Exploiter des courbes d'activité enzymatique en fonction du pH}

\begin{methode}[Lire et interpréter une courbe activité = f(pH)]
\begin{enumerate}
\item \textbf{Identifier les axes} : pH en abscisse (échelle de 0 à 14), activité enzymatique en ordonnée.
\item \textbf{Décrire la courbe} : elle a généralement une forme de \cle{cloche} : l'activité croît jusqu'à un maximum puis décroît.
\item \textbf{Repérer le pH optimal} : abscisse du sommet de la courbe. Exemples à connaître : pepsine gastrique vers pH 2 (milieu acide de l'estomac), amylase salivaire vers pH 7, trypsine pancréatique vers pH 8 (milieu basique de l'intestin grêle).
\item \textbf{Interpréter} : loin du pH optimal, la concentration en ions \ce{H+} modifie les charges électriques portées par les chaînes latérales des acides aminés du site actif (protonation/déprotonation), ce qui déforme le site actif ou empêche la catalyse ; à pH extrême, il y a dénaturation globale.
\item \textbf{Comparer éventuellement deux courbes} (deux enzymes) et relier chaque pH optimal au pH du milieu où l'enzyme agit normalement dans l'organisme.
\end{enumerate}
\textbf{Réflexe Probatoire} : au Probatoire, relie toujours le pH optimal au \emph{milieu physiologique} de l'enzyme (estomac, salive, intestin) : c'est la justification attendue.
\end{methode}

\begin{exoresolu}[Comparaison de la pepsine et de la trypsine]
On mesure l'activité de deux enzymes digestives, la pepsine (estomac) et la trypsine (intestin grêle), en fonction du pH du milieu. Les résultats (activité en unités arbitraires, u.a.) sont :

\begin{center}
\begin{tabularx}{\linewidth}{|l|X|X|X|X|X|X|X|}\hline
\rowcolor{popblueL}
pH & 1 & 2 & 4 & 6 & 7 & 8 & 10 \\ \hline
Pepsine (u.a.) & 8 & 12 & 6 & 2 & 1 & 0 & 0 \\ \hline
Trypsine (u.a.) & 0 & 0 & 2 & 6 & 9 & 12 & 6 \\ \hline
\end{tabularx}
\end{center}

\textbf{1.} Détermine le pH optimal de chaque enzyme.\\
\textbf{2.} Justifie ces valeurs par rapport au milieu de fonctionnement de chaque enzyme.\\
\textbf{3.} Que se passe-t-il pour la pepsine lorsque le contenu de l'estomac passe dans l'intestin grêle (pH voisin de 8) ?
\tcblower
\textbf{1. pH optimaux.} L'activité de la pepsine est maximale (\SI{12}{u.a.}) à pH 2 : son pH optimal est \textbf{2}. L'activité de la trypsine est maximale (\SI{12}{u.a.}) à pH 8 : son pH optimal est \textbf{8}.

\textbf{2. Justification.} Dans l'estomac, les cellules pariétales sécrètent de l'acide chlorhydrique (\ce{HCl}) : le pH gastrique est voisin de 2. La pepsine, qui digère les protéines dans l'estomac, est donc adaptée à un milieu très acide. Dans l'intestin grêle, le suc pancréatique riche en bicarbonates (\ce{HCO3-}) neutralise l'acidité : le pH est voisin de 8, ce qui correspond au pH optimal de la trypsine, enzyme du suc pancréatique. Chaque enzyme présente un pH optimal correspondant au pH de son milieu naturel d'action.

\textbf{3. Devenir de la pepsine dans l'intestin.} À pH 8, l'activité de la pepsine est nulle (0 u.a.) : l'excès d'ions \ce{OH-} (ou la faiblesse en \ce{H+}) modifie les charges des acides aminés du site actif, qui se déforme ; la pepsine est inactivée, voire dénaturée. La digestion des protéines est alors poursuivie par la trypsine, active à ce pH.

\textbf{Conclusion} : chaque enzyme possède un pH optimal ; la succession estomac (pH 2) puis intestin (pH 8) permet l'action coordonnée de la pepsine puis de la trypsine dans la digestion des protéines.
\end{exoresolu}

\courssub{Interpréter une expérience mettant en évidence la spécificité d'une enzyme}

\begin{methode}[Analyser une expérience de spécificité enzymatique]
\begin{enumerate}
\item \textbf{Dégager le problème} : une enzyme donnée agit-elle sur n'importe quel substrat ?
\item \textbf{Identifier les variables} : ce qui change entre les tubes (la nature du substrat, ou la nature de l'enzyme) et ce qui est maintenu constant (température, pH, volumes, durée). Citer le \cle{tube témoin}.
\item \textbf{Décrire les résultats tube par tube} avec les observations exactes (coloration, dégagement gazeux, présence/absence de produit révélée par un test : liqueur de Fehling, eau iodée, etc.).
\item \textbf{Comparer} : la réaction n'a lieu que dans le(s) tube(s) où l'enzyme rencontre « son » substrat.
\item \textbf{Interpréter par le modèle clé--serrure / ajustement induit} : le substrat se fixe dans le \cle{site actif} de l'enzyme, dont la forme est complémentaire de la sienne ; un autre substrat ne peut pas s'y fixer.
\item \textbf{Conclure} : « l'enzyme est spécifique de son substrat » (spécificité de substrat) et « elle catalyse un seul type de réaction » (spécificité d'action).
\end{enumerate}
\textbf{Réflexe} : cite toujours le complexe \cle{enzyme--substrat} dans l'interprétation, et précise que la spécificité repose sur la forme spatiale du site actif.
\end{methode}

\begin{exoresolu}[L'amylase salivaire est-elle spécifique ?]
Un élève de Première D du lycée de Nkongsamba réalise l'expérience suivante à \SI{37}{\celsius} :

\begin{center}
\begin{tabularx}{\linewidth}{|l|X|X|X|}\hline
\rowcolor{popblueL}
 & Tube 1 & Tube 2 & Tube 3 \\ \hline
Substrat & Amidon & \textbf{Saccharose} & Amidon \\ \hline
Ajout & Salive (amylase) & Salive (amylase) & Eau distillée \\ \hline
Test après 15 min & Eau iodée : jaune ; Fehling : rouge brique & Fehling : reste bleu & Eau iodée : bleu-noir \\ \hline
\end{tabularx}
\end{center}

\textbf{1.} Quel est le rôle du tube 3 ?\\
\textbf{2.} Interprète les résultats des tubes 1 et 2.\\
\textbf{3.} Conclus sur la spécificité de l'amylase salivaire.
\tcblower
\textbf{1. Rôle du tube 3.} Le tube 3 est le \textbf{tube témoin} : il contient le substrat (amidon) sans enzyme (l'eau distillée remplace la salive). Il permet de vérifier que l'amidon ne se transforme pas spontanément dans les conditions de l'expérience. Le test à l'eau iodée bleu-noir confirme que l'amidon est encore présent : sans enzyme, aucune réaction n'a lieu.

\textbf{2. Interprétation.}
\begin{itemize}
\item \textbf{Tube 1} : l'eau iodée donne une coloration jaune, donc l'amidon a \emph{disparu} (hydrolysé) ; la liqueur de Fehling donne un précipité rouge brique, donc un sucre réducteur (le maltose) est \emph{apparu}. L'amylase salivaire a donc hydrolysé l'amidon en maltose : la réaction a eu lieu. Le complexe \textbf{amylase--amidon} s'est formé.
\item \textbf{Tube 2} : le substrat est le saccharose (sucre non réducteur, non hydrolysable par l'amylase). La liqueur de Fehling reste bleue : aucun sucre réducteur n'est formé. L'amylase n'a exercé aucune action sur le saccharose. Le saccharose ne peut pas former de complexe avec le site actif de l'amylase.
\end{itemize}
La même enzyme, dans les mêmes conditions de température et de durée, agit sur l'amidon mais pas sur le saccharose.

\textbf{3. Conclusion.} L'amylase salivaire n'agit que sur l'amidon (et polysaccharides apparentés) : elle est \textbf{spécifique de son substrat}. Cette spécificité s'explique par la complémentarité de forme et de charges entre le site actif de l'amylase et l'amidon (modèle clé--serrure / ajustement induit) ; le saccharose ne peut pas s'y fixer. L'enzyme catalyse en outre une seule réaction : l'hydrolyse des liaisons $\alpha$-1,4-glucosidiques (spécificité d'action).
\end{exoresolu}

\courssub{Étudier l'influence de la concentration (enzyme ou substrat)}

\begin{methode}[Exploiter des données de cinétique enzymatique]
\begin{enumerate}
\item \textbf{Distinguer les deux situations} : on fait varier soit la concentration en \cle{substrat} (enzyme constante), soit la concentration en \cle{enzyme} (substrat en excès, dit « saturant »).
\item \textbf{Substrat variable (enzyme constante) :} la vitesse de réaction augmente d'abord proportionnellement à la concentration en substrat (premier ordre), puis se \cle{stabilise} (palier, ordre zéro) : tous les sites actifs sont occupés en permanence, on dit que l'enzyme est \cle{saturé}. La vitesse maximale est notée $V_{\max}$.
\item \textbf{Enzyme variable (substrat en excès saturant) :} la vitesse est proportionnelle à la concentration en enzyme : doubler l'enzyme double la vitesse (droite passant par l'origine).
\item \textbf{Exploiter numériquement} : calcule des rapports de vitesses, vérifie la proportionnalité, ou repère le palier à partir duquel l'ajout de substrat n'augmente plus la vitesse.
\item \textbf{Conclure} en précisant le facteur limitant dans chaque cas (sites actifs disponibles ou quantité d'enzyme).
\end{enumerate}
\textbf{Réflexe} : devant un palier, la seule explication attendue est la \emph{saturation des sites actifs} ; devant une droite passant par l'origine, la \emph{proportionnalité}.
\end{methode}

\begin{exoresolu}[Influence de la concentration en substrat sur la vitesse de réaction]
On mesure la vitesse initiale $V$ d'une réaction enzymatique (en µmol de produit formé par minute) pour différentes concentrations en substrat $[S]$ (en mmol/L), la concentration en enzyme étant constante :

\begin{center}
\begin{tabularx}{\linewidth}{|l|X|X|X|X|X|X|}\hline
\rowcolor{popblueL}
$[S]$ (mmol/L) & 0 & 1 & 2 & 4 & 8 & 16 \\ \hline
$V$ (µmol/min) & 0 & 10 & 18 & 30 & 38 & 40 \\ \hline
\end{tabularx}
\end{center}

\textbf{1.} Décris l'évolution de $V$ en fonction de $[S]$.\\
\textbf{2.} Calcule le rapport $\dfrac{V}{[S]}$ pour $[S] = 1$ puis pour $[S] = 16$. La proportionnalité est-elle conservée ?\\
\textbf{3.} Interprète le palier observé.
\tcblower
\textbf{1. Description.} Pour $[S]$ passant de 0 à 8 mmol/L, $V$ augmente fortement (de 0 à 38 µmol/min) ; entre 8 et 16 mmol/L, $V$ n'augmente plus que de 2 µmol/min et tend vers une valeur limite voisine de 40 µmol/min : la courbe présente une phase de croissance puis un \textbf{palier}.

\textbf{2. Test de proportionnalité.}
\[ \frac{V}{[S]} = \frac{10}{1} = 10 \quad \text{pour } [S] = 1 \text{ mmol/L} \]
\[ \frac{V}{[S]} = \frac{40}{16} = 2{,}5 \quad \text{pour } [S] = 16 \text{ mmol/L} \]
Les rapports sont différents ($10 \neq 2{,}5$) : la vitesse \textbf{n'est pas proportionnelle} à la concentration en substrat sur tout le domaine. La proportionnalité n'est approchée qu'aux faibles concentrations.

\textbf{3. Interprétation du palier.} Aux faibles concentrations, de nombreux sites actifs sont libres : chaque molécule de substrat ajoutée trouve une enzyme disponible, donc la vitesse augmente. Aux fortes concentrations, toutes les molécules d'enzyme sont occupées en permanence par le substrat : les sites actifs sont \textbf{saturés}. Ajouter du substrat ne peut plus accélérer la réaction, car le facteur limitant devient la quantité d'enzyme. La vitesse atteint sa valeur maximale $V_{\max} \approx 40$ µmol/min.

\textbf{Conclusion} : à concentration en enzyme constante, la vitesse de la réaction croît avec la concentration en substrat jusqu'à un maximum correspondant à la saturation des sites actifs de l'enzyme.
\end{exoresolu}

\courssub{Relier un dysfonctionnement enzymatique à une pathologie métabolique}

\begin{methode}[Raisonner sur un déficit enzymatique]
\begin{enumerate}
\item \textbf{Identifier la réaction bloquée} : quelle enzyme est déficiente ou absente, et quelle transformation (substrat $\rightarrow$ produit) ne se fait plus ?
\item \textbf{Déduire les deux conséquences logiques} : \cle{accumulation} du substrat en amont du blocage (éventuellement toxique) et \cle{carence} du produit en aval.
\item \textbf{Relier aux signes cliniques} : chaque symptôme doit être rattaché à l'accumulation ou à la carence.
\item \textbf{Proposer une prise en charge cohérente} : régime limitant le substrat accumulé, supplémentation en produit manquant, ou apport de l'enzyme (thérapie enzymatique).
\item \textbf{Rédiger une chaîne causale complète} : déficit enzymatique $\rightarrow$ perturbation métabolique (accumulation/carence) $\rightarrow$ symptôme.
\end{enumerate}
\textbf{Réflexe Probatoire} : la note est gagnée si la \emph{chaîne causale complète} est écrite, pas seulement le nom de la maladie. Exemples de référence : intolérance au lactose (déficit en lactase), phénylcétonurie (déficit en phénylalanine-hydroxylase), albinisme (déficit en tyrosinase), maladie de Pompe (déficit en $\alpha$-glucosidase), mucoviscidose (déficit fonctionnel de la protéine CFTR, canal chlorure).
\end{methode}

\begin{exoresolu}[L'intolérance au lactose]
Après avoir bu du lait frais, plusieurs élèves d'un internat de Bafoussam présentent des douleurs abdominales, des ballonnements et des diarrhées. Le médecin scolaire explique que leur muqueuse intestinale produit peu de \textbf{lactase}, enzyme qui hydrolyse le lactose du lait en glucose et en galactose. Des analyses montrent la présence de lactose intact et d'acides organiques dans les selles.

\textbf{1.} Écris l'équation de la réaction normalement catalysée par la lactase.\\
\textbf{2.} Explique l'origine des symptômes observés.\\
\textbf{3.} Propose une mesure de prévention.
\tcblower
\textbf{1. Équation de la réaction.} La lactase catalyse l'hydrolyse du lactose (disaccharide du lait) :
\[ \text{lactose} + \ce{H2O} \xrightarrow{\text{lactase}} \text{glucose} + \text{galactose} \]
Ces deux monosaccharides sont ensuite absorbés par la muqueuse de l'intestin grêle (transport actif \ce{Na+}-dépendant).

\textbf{2. Origine des symptômes.} Chez ces élèves, le déficit en lactase empêche l'hydrolyse du lactose. On établit la chaîne causale :
\begin{itemize}
\item le \textbf{lactose non digéré s'accumule} dans la lumière intestinale (d'où sa présence dans les selles) ; il retient l'eau par osmose (effet osmotique), ce qui provoque les \textbf{diarrhées} aqueuses ;
\item les bactéries de la flore colique \textbf{fermentent} ce lactose (glycolyse anaérobie), produisant des gaz (\ce{H2}, \ce{CO2}, \ce{CH4} $\rightarrow$ \textbf{ballonnements} et douleurs abdominales) et des acides organiques (lactate, acétate, butyrate $\rightarrow$ pH fécal acide, irritation muqueuse, accélération du transit).
\end{itemize}
Il y a en outre une \textbf{carence relative en glucose et galactose} absorbés, mais ce sont surtout l'accumulation du substrat et sa fermentation bactérienne qui expliquent les signes cliniques.

\textbf{3. Prévention.} Éviter ou réduire la consommation de lait frais ; utiliser des laits « sans lactose » (pré-hydrolysés) ou des produits fermentés (yaourts, lait caillé, fromages affinés) dont une partie du lactose est déjà hydrolysée par les ferments lactiques ; éventuellement prendre des gélules de lactase (enzyme exogène) avant le repas.

\textbf{Conclusion} : l'intolérance au lactose illustre comment le déficit d'une seule enzyme bloque une étape du métabolisme et entraîne des troubles par accumulation du substrat non transformé et fermentation bactérienne.
\end{exoresolu}

\courssub{Relier les enzymes au renouvellement moléculaire de la cellule}

\begin{methode}[Expliquer le rôle des enzymes dans les réactions de synthèse et de dégradation]
\begin{enumerate}
\item \textbf{Rappeler le principe} : sans enzyme, les réactions du métabolisme seraient trop lentes (cinétique négligeable à \SI{37}{\celsius}) pour être compatibles avec la vie ; l'enzyme \cle{abaisse l'énergie d'activation} et accélère la réaction sans être consommée (rôle catalytique).
\item \textbf{Distinguer les deux grands types de réactions métaboliques} :
\begin{itemize}
\item réactions de \cle{dégradation} (\cle{catabolisme}) : grosses molécules $\rightarrow$ petites molécules, généralement exergoniques (libération d'énergie, ex. hydrolyses digestives, glycolyse, $\beta$-oxydation, respiration cellulaire) ;
\item réactions de \cle{synthèse} (\cle{anabolisme}) : petites molécules $\rightarrow$ grosses molécules, généralement endergoniques (nécessitant un apport d'énergie, souvent couplées à l'hydrolyse de l'ATP, ex. synthèse des protéines, de l'ADN, du glycogène, des lipides).
\end{itemize}
\item \textbf{Associer chaque réaction à son enzyme} : une enzyme catalyse une réaction précise ; le renouvellement moléculaire permanent de la cellule résulte de l'action coordonnée de nombreuses enzymes (voies métaboliques).
\item \textbf{Vérifier le bilan d'une équation} : réactifs, produits, et nom de l'enzyme au-dessus de la flèche.
\end{enumerate}
\textbf{Réflexe} : une enzyme n'est ni consommée ni modifiée à la fin de la réaction : elle peut servir à nouveau (turnover catalytique). Le dire, c'est définir un catalyseur.
\end{methode}

\begin{exoresolu}[Synthèse et dégradation du glycogène dans le foie]
Chez un élève bien nourri, le foie stocke du glycogène après le repas ; pendant un effort sportif prolongé (par exemple un match de football au quartier), ce glycogène est dégradé pour fournir du glucose au sang. Ces deux transformations sont catalysées par des enzymes distinctes : la glycogène-synthase (synthèse) et la glycogène-phosphorylase (dégradation).

\textbf{1.} Pourquoi ces deux réactions nécessitent-elles des enzymes ?\\
\textbf{2.} Classe chaque réaction dans le catabolisme ou l'anabolisme.\\
\textbf{3.} En quoi cet exemple illustre-t-il le renouvellement moléculaire ?
\tcblower
\textbf{1. Nécessité des enzymes.} À la température du corps (\SI{37}{\celsius}) et au pH cellulaire, les réactions de synthèse et de dégradation du glycogène sont spontanément extrêmement lentes (haute énergie d'activation). Les enzymes, catalyseurs biologiques, abaissent cette énergie d'activation et multiplient la vitesse de ces réactions (facteur $10^{10}$ et plus), les rendant compatibles avec les besoins rapides de l'organisme (par exemple la libération de glucose en quelques secondes pendant l'effort). Elles ne sont pas consommées et peuvent catalyser des millions de cycles (nombre de turnover élevé).

\textbf{2. Classification.}
\begin{itemize}
\item La \textbf{synthèse du glycogène} à partir du glucose (catalysée par la glycogène-synthase, utilisant l'UDP-glucose et l'ATP) est une réaction d'\textbf{anabolisme} : construction d'une grosse molécule (polymère) à partir de petites molécules (monomères), avec consommation d'énergie.
\item La \textbf{dégradation du glycogène} en glucose-1-phosphate (catalysée par la glycogène-phosphorylase, utilisant le phosphate inorganique) est une réaction de \textbf{catabolisme} : fragmentation d'une grosse molécule, mettant à disposition du glucose-6-phosphate dont l'oxydation (glycolyse + respiration) libérera de l'énergie utilisable (ATP).
\end{itemize}

\textbf{3. Renouvellement moléculaire.} Les molécules de l'organisme ne sont pas figées : le glycogène du foie est continuellement synthétisé (état post-prandial) puis dégradé (jeûne, effort) selon les besoins ; de même, les protéines (turnover protéique), les lipides et les acides nucléiques des cellules sont sans cesse détruits (catabolisme) et reconstruits (anabolisme). Ce renouvellement permanent, qui maintient la structure, l'homéostasie et le fonctionnement de la cellule, n'est possible que grâce à l'action spécifique, rapide et finement régulée (hormones, allostérie) de nombreuses enzymes.

\textbf{Conclusion} : les enzymes assurent la vitesse et la spécificité des réactions de synthèse (anabolisme) et de dégradation (catabolisme) qui réalisent le renouvellement moléculaire indispensable à la vie de la cellule.
\end{exoresolu}

\begin{attention}[Les 5 erreurs qui coûtent des points au Probatoire]
\begin{enumerate}
\item \textbf{Écrire « l'enzyme meurt » ou « est tuée ».} Une enzyme est une molécule, pas un être vivant. À haute température, elle est \emph{dénaturée} (structure spatiale détruite) ; à basse température, elle est seulement \emph{ralentie}, de façon réversible. Confondre ces deux situations est l'erreur la plus fréquente.
\item \textbf{Oublier le tube témoin dans l'analyse d'une expérience.} Toute interprétation expérimentale exige de citer le tube témoin et de préciser les conditions maintenues constantes (température, pH, durée, volumes). Sans cela, la conclusion n'est pas validée.
\item \textbf{Confondre spécificité de substrat et conditions optimales.} Dire « l'amylase n'agit pas à pH 2 parce qu'elle est spécifique de l'amidon » mélange deux notions : la spécificité concerne la nature du substrat (modèle clé--serrure), pas la température ou le pH, qui sont des \emph{facteurs} de l'activité.
\item \textbf{Interpréter un palier par « l'enzyme est épuisée ».} Une enzyme n'est jamais consommée : le palier de la courbe $V = f([S])$ traduit la \emph{saturation des sites actifs}, c'est-à-dire que toutes les molécules d'enzyme sont occupées en permanence.
\item \textbf{Relier une pathologie à un déficit enzymatique sans chaîne causale.} Écrire seulement « le déficit en lactase provoque des diarrhées » est insuffisant : il faut enchaîner déficit enzymatique $\rightarrow$ réaction bloquée $\rightarrow$ accumulation du substrat (et/ou carence du produit) $\rightarrow$ symptôme. Chaque étape rapporte des points.
\end{enumerate}
\end{attention}

\newpage

\coursec{Exercices}

\coursec{Sensibilisation sur l'influence des enzymes sur les réactions chimiques indispensables au renouvellement moléculaire}

\courssub{Je vérifie mes connaissances}

\begin{exercice}[QCM : Enzyme, substrat et catalyse]
Parmi les affirmations suivantes, une seule est correcte. Laquelle ?
\begin{enumerate}
    \item Une enzyme est un glucide qui accélère une réaction chimique sans être modifiée.
    \item Le substrat se fixe sur le site allostérique de l'enzyme pour former le complexe enzyme-substrat.
    \item Une enzyme abaisse l'énergie d'activation d'une réaction chimique et accélère l'atteinte de l'équilibre.
    \item La dénaturation d'une enzyme par la température est toujours réversible.
\end{enumerate}
\end{exercice}

\begin{exercice}[Vrai ou Faux justifié]
Indiquez si les propositions suivantes sont vraies (V) ou fausses (F) et justifiez brièvement votre réponse.
\begin{enumerate}
    \item La spécificité d'une enzyme est due à la complémentarité de forme et de charges entre son site actif et son substrat.
    \item L'augmentation indéfinie de la concentration en substrat entraîne une augmentation indéfinie de la vitesse de réaction enzymatique.
    \item La pepsine, enzyme digestive de l'estomac, a un pH optimal acide (autour de 2).
    \item Un inhibiteur compétitif se fixe de manière irréversible sur le site actif de l'enzyme.
\end{enumerate}
\end{exercice}

\begin{exercice}[Définitions et notions clés]
Donnez une définition précise et rigoureuse des termes suivants :
\begin{enumerate}
    \item Enzyme
    \item Site actif
    \item Spécificité enzymatique (spécificité de substrat et de réaction)
    \item Dénaturation enzymatique
    \item Vitesse maximale ($V_{\text{max}}$) et constante de Michaelis ($K_m$)
\end{enumerate}
\end{exercice}

\begin{exercice}[Calcul direct : Rendement catalytique]
Une enzyme catalyse la conversion de 120 µmol de substrat par minute dans un volume de 3 mL. La concentration de l'enzyme est de $2 \times 10^{-8}$ mol·L$^{-1}$.
Calculer le nombre de tours catalytiques ($k_{\text{cat}}$ ou turnover number) de cette enzyme en s$^{-1}$.
\emph{Rappel : $k_{\text{cat}} = \frac{V_{\text{max}}}{[E]_{\text{total}}}$.}
\end{exercice}

\courssub{Je m'entraîne}

\begin{exercice}[Exploitation de courbe : Activité d'une amylase en fonction de la température]
On étudie l'activité d'une amylase salivaire humaine à différentes températures. Le graphique ci-dessous représente la vitesse initiale de réaction ($v_0$) en fonction de la température.

\begin{popfigure}
\begin{tikzpicture}[scale=0.7]
\begin{axis}[
    width=\linewidth, height=8cm,
    xlabel={Température ($^\circ$C)}, ylabel={Vitesse initiale $v_0$ (UA)},
    xmin=0, xmax=80, ymin=0, ymax=110,
    xtick={0,10,20,30,37,40,50,60,70,80}, ytick={0,20,40,60,80,100},
    grid=major, grid style={popline},
    axis lines=middle, enlargelimits=0.05,
    clip=false
]
\addplot[popcurve, thick, domain=0:37, samples=100] {100*exp(-(x-37)^2/50)};
\addplot[popcurve, thick, domain=37:80, samples=100] {100*exp(-(x-37)^2/50) - 2*(x-37)^1.5};
\addplot[popcurve, thick, domain=65:80, samples=50] {5};
\node[popdark, anchor=south] at (axis cs:37,100) {$\bullet$};
\node[popdark, anchor=north, font=\small] at (axis cs:37,105) {Optimum};
\end{axis}
\end{tikzpicture}
\legende{Activité de l'amylase salivaire humaine en fonction de la température. Optimum à 37 °C.}
\end{popfigure}

\begin{enumerate}
    \item Déterminer graphiquement la température optimale de cette amylase.
    \item Expliquer l'allure de la courbe entre 0 °C et la température optimale (phénomène cinétique).
    \item Expliquer l'allure de la courbe au-delà de la température optimale (phénomène structural).
    \item Prédire et justifier ce qui se passe si on porte l'enzyme à 80 °C pendant 10 minutes puis qu'on la replace à 37 °C.
\end{enumerate}
\end{exercice}

\begin{exercice}[Interprétation d'expérience : Spécificité et thermosensibilité de l'amylase]
On réalise l'expérience suivante dans trois tubes à essai maintenus à 37 °C :
\begin{itemize}
    \item \textbf{Tube 1 :} 2 mL de solution d'amidon (1 \%) + 1 mL de salive fraîche.
    \item \textbf{Tube 2 :} 2 mL de solution de cellulose (1 \%) + 1 mL de salive fraîche.
    \item \textbf{Tube 3 :} 2 mL de solution d'amidon (1 \%) + 1 mL de salive bouillie 5 min puis refroidie.
\end{itemize}
Après 15 minutes d'incubation, on effectue deux tests de révélation sur chaque tube :
\begin{itemize}
    \item Test à l'eau iodée (révèle l'amidon : couleur bleu-noir si amidon présent).
    \item Test à la liqueur de Fehling (révèle les sucres réducteurs : précipité rouge brique si sucres réducteurs présents).
\end{itemize}
Les résultats sont consignés dans le tableau :

\begin{center}
\begin{tabularx}{\linewidth}{|c|X|X|}\hline
\rowcolor{popblueL}
\textbf{Tube} & \textbf{Eau iodée} & \textbf{Liqueur de Fehling} \\ \hline
1 & Jaune (pas de coloration bleu-noir) & Précipité rouge brique \\ \hline
2 & Bleu-noir & Bleu (pas de précipité) \\ \hline
3 & Bleu-noir & Bleu (pas de précipité) \\ \hline
\end{tabularx}
\end{center}

\begin{enumerate}
    \item Interpréter les résultats du Tube 1. Quelle réaction a eu lieu ? Écrire l'équation bilan simplifiée (avec \ce{->}).
    \item Interpréter les résultats du Tube 2. Que conclure sur la spécificité de l'amylase ?
    \item Interpréter les résultats du Tube 3. Que conclure sur la nature chimique de l'enzyme et sa sensibilité à la température ?
    \item Résumer les deux propriétés fondamentales des enzymes mises en évidence par cette expérience.
\end{enumerate}
\end{exercice}

\begin{exercice}[Comparaison de courbes d'activité en fonction du pH]
Le document ci-dessous présente les courbes d'activité relative (en \%) de trois enzymes digestives (A, B, C) en fonction du pH.

\begin{popfigure}
\begin{tikzpicture}[scale=0.7]
\begin{axis}[
    width=\linewidth, height=7cm,
    xlabel={pH}, ylabel={Activité relative (\%)},
    xmin=0, xmax=10, ymin=0, ymax=105,
    xtick={0,1,2,3,4,5,6,7,8,9,10}, ytick={0,25,50,75,100},
    grid=major, grid style={popline},
    axis lines=middle, enlargelimits=0.05,
    legend pos=north east, legend cell align=left
]
\addplot[popblue, thick, domain=0:10, samples=200] {100*exp(-(x-2)^2/1.5)};
\addlegendentry{Enzyme A}
\addplot[popgreen, thick, domain=0:10, samples=200] {100*exp(-(x-6.8)^2/2)};
\addlegendentry{Enzyme B}
\addplot[poporange, thick, domain=0:10, samples=200] {100*exp(-(x-8)^2/1.8)};
\addlegendentry{Enzyme C}
\end{axis}
\end{tikzpicture}
\legende{Activité relative de trois enzymes digestives en fonction du pH.}
\end{popfigure}

Sachant que :
\begin{itemize}
    \item La pepsine agit dans l'estomac (pH $\approx$ 1,5--3).
    \item L'amylase salivaire agit dans la bouche (pH $\approx$ 6,5--7,5).
    \item La trypsine agit dans l'intestin grêle (pH $\approx$ 7,5--8,5).
\end{itemize}

\begin{enumerate}
    \item Identifier chacune des enzymes A, B et C (pepsine, amylase salivaire, trypsine) en justifiant par leur pH optimal et leur lieu d'action.
    \item Expliquer au niveau moléculaire pourquoi l'activité chute brutalement de part et d'autre du pH optimal (rôle des groupements ionisables du site actif).
    \item Que se passerait-il si on introduisait de la pepsine dans le duodénum (pH $\approx$ 8) ? Justifier.
\end{enumerate}
\end{exercice}

\begin{exercice}[Construction et interprétation : Vitesse en fonction de la concentration en substrat]
On mesure la vitesse initiale ($v_0$) d'une réaction enzymatique pour différentes concentrations en substrat $[S]$. Résultats :

\begin{center}
\begin{tabular}{|c|c|c|c|c|c|c|c|}\hline
\rowcolor{popblueL}
$[S]$ (mmol·L$^{-1}$) & 0,5 & 1,0 & 2,0 & 4,0 & 8,0 & 16,0 & 32,0 \\ \hline
$v_0$ (µmol·min$^{-1}$·mL$^{-1}$) & 12 & 20 & 33 & 50 & 67 & 80 & 89 \\ \hline
\end{tabular}
\end{center}

\begin{enumerate}
    \item Tracer la courbe $v_0 = f([S])$ sur papier millimétré (ou décrire la forme attendue).
    \item Déterminer graphiquement (ou estimer) $V_{\text{max}}$ et $K_m$ (concentration en substrat pour $v_0 = V_{\text{max}}/2$).
    \item Expliquer la forme hyperbolique de la courbe en utilisant le modèle du complexe enzyme-substrat (saturation des sites actifs).
    \item Calculer la vitesse initiale attendue pour $[S] = K_m$ et pour $[S] = 10 \times K_m$. Commenter.
\end{enumerate}
\end{exercice}

\courssub{Je me prépare au Probatoire}

\begin{exercice}[Cas clinique : Intolérance au lactose du nourrisson]
\textbf{Contexte :} Un nourrisson de 3 mois, allaité exclusivement au sein, présente depuis la naissance des diarrhées aqueuses, des ballonnements, des crampes abdominales et des vomissements après chaque tétée. Il a une courbe de poids qui fléchit. Le pédiatre suspecte une intolérance au lactose. Il demande un test d'expiration à l'hydrogène (H$_2$) après ingestion de lactose et un dosage de l'activité lactasique sur biopsie intestinale.

\textbf{Résultats :}
\begin{itemize}
    \item Test H$_2$ : Pic d'hydrogène expiré à 45 ppm (normal $<$ 10 ppm) 90 min après l'ingestion.
    \item Activité lactasique (biopsie) : 0,2 U/g de protéines (normal : 10--30 U/g).
    \item Activité sucrase-isomaltase : Normale.
\end{itemize}

\textbf{Rappels :} Le lactose (glucide du lait) est un disaccharide (Glucose + Galactose). La lactase ($\beta$-galactosidase) est une enzyme de la bordure en brosse des entérocytes de l'intestin grêle. L'hydrogène expiré provient de la fermentation bactérienne colique des glucides non absorbés.

\begin{enumerate}
    \item Rappeler l'équation de l'hydrolyse du lactose catalysée par la lactase (utiliser \ce{->}).
    \item Expliquer le mécanisme physiopathologique reliant le déficit en lactase aux symptômes digestifs (diarrhée, ballonnements, H$_2$ expiré). Utiliser les termes : substrat, accumulation, pression osmotique, fermentation, flore colique.
    \item Pourquoi l'activité de la sucrase-isomaltase est-elle normale ? Que cela prouve-t-il sur l'origine du déficit (génétique vs acquis généralisé) ?
    \item Proposer deux mesures diététiques ou thérapeutiques adaptées pour ce nourrisson. Justifier chacune.
    \item Différencier « intolérance au lactose » et « allergie aux protéines de lait de vache » sur le plan mécanistique.
\end{enumerate}
\end{exercice}

\begin{exercice}[Étude expérimentale : Inhibition enzymatique et thérapeutique]
\textbf{Contexte :} L'enzyme de conversion de l'angiotensine (ECA) catalyse la transformation de l'angiotensine I (inactive) en angiotensine II (puissant vasoconstricteur). L'inhibition de l'ECA est une stratégie thérapeutique majeure dans le traitement de l'hypertension artérielle (fréquente au Cameroun). On étudie l'effet d'un inhibiteur (I) sur la cinétique de l'ECA.

On mesure la vitesse initiale $v_0$ pour différentes concentrations en substrat $[S]$, en l'absence (témoin) et en présence de deux concentrations de l'inhibiteur $[I_1]$ et $[I_2]$ ($[I_2] > [I_1]$). Les résultats sont représentés dans le graphique de Lineweaver-Burk ($1/v_0 = f(1/[S])$) ci-dessous.

\begin{popfigure}
\begin{tikzpicture}[scale=0.7]
\begin{axis}[
    width=\linewidth, height=8cm,
    xlabel={$1/[S]$ (L·mmol$^{-1}$)}, ylabel={$1/v_0$ (min·mL·µmol$^{-1}$)},
    xmin=0, xmax=4, ymin=0, ymax=0.6,
    xtick={0,1,2,3,4}, ytick={0,0.1,0.2,0.3,0.4,0.5,0.6},
    grid=major, grid style={popline},
    axis lines=middle, enlargelimits=0.05,
    legend pos=north west, legend cell align=left
]
% Temoin
\addplot[popblue, thick, domain=0.2:4, samples=50] {0.1 + 0.1*x};
\addlegendentry{Témoin (sans I)}
% I1
\addplot[popgreen, thick, domain=0.2:4, samples=50] {0.1 + 0.2*x};
\addlegendentry{$[I_1]$}
% I2
\addplot[poporange, thick, domain=0.2:4, samples=50] {0.1 + 0.35*x};
\addlegendentry{$[I_2]$}
\end{axis}
\end{tikzpicture}
\legende{Représentation de Lineweaver-Burk de l'activité ECA en présence d'inhibiteur. Les droites se coupent sur l'axe des ordonnées.}
\end{popfigure}

\begin{enumerate}
    \item En analysant le point d'intersection des trois droites, déterminer le type d'inhibition (compétitive, non compétitive, anti-compétitive). Justifier votre réponse en précisant ce qui varie ($V_{\text{max}}$, $K_m$ apparent) et ce qui reste constant.
    \item Expliquer le mécanisme moléculaire de ce type d'inhibition (site de fixation de l'inhibiteur par rapport au site actif, compétition ou non avec le substrat).
    \item Calculer $V_{\text{max}}$ (en µmol·min$^{-1}$·mL$^{-1}$) et $K_m$ (en mmol·L$^{-1}$) pour le témoin à partir de l'ordonnée à l'origine et de l'abscisse à l'origine de la droite témoin.
    \item Déduire l'effet de l'augmentation de la concentration en inhibiteur sur la pente de la droite. Relier cela à la constante d'inhibition $K_i$.
    \item Pourquoi un inhibiteur compétitif de l'ECA est-il cliniquement efficace pour baisser la tension artérielle, même si la concentration en substrat (angiotensine I) peut varier physiologiquement ?
\end{enumerate}
\end{exercice}

\begin{exercice}[Synthèse : Enzymes, environnement et santé publique au Cameroun]
\textbf{Document 1 :} La pollution par les métaux lourds (plomb Pb$^{2+}$, mercure Hg$^{2+}$, cadmium Cd$^{2+}$) issue de l'exploitation minière artisanale (ex : région de l'Est, Bétaré-Oya) ou des batteries usagées contaminent les cours d'eau. Ces ions métalliques sont de puissants inhibiteurs enzymatiques non compétitifs irréversibles. Ils se fixent sur les groupements -SH (thiols) des cystéines essentielles au site actif ou à la structure tertiaire de nombreuses enzymes (ex : $\delta$-ALA déshydratase pour le plomb, enzymes du cycle de Krebs).

\textbf{Document 2 :} La fièvre jaune, endémique dans certaines zones forestières du Sud-Cameroun, est due à un flavivirus. Le virus code pour une protéase virale (NS2B-NS3) essentielle à la maturation des protéines virales. Des recherches visent à concevoir des inhibiteurs spécifiques de cette protéase.

\textbf{Document 3 :} Dans l'industrie agroalimentaire camerounaise (brasseries, huileries, jus de fruits), on utilise des enzymes exogènes : amylases (hydrolyse amidon), pectinases (clarification jus), protéases (attendrissement viande, détergents).

\begin{enumerate}
    \item \textbf{(Doc 1)} Expliquer au niveau moléculaire pourquoi la fixation d'un ion Hg$^{2+}$ sur un groupement -SH d'une enzyme entraîne une perte d'activité irréversible. Préciser le type d'inhibition et son incidence sur $V_{\text{max}}$ et $K_m$.
    \item \textbf{(Doc 1)} Le plomb inhibe la $\delta$-ALA déshydratase, enzyme clé de la synthèse de l'hème (constituant de l'hémoglobine). Expliquer comment ce blocage enzymatique conduit à une anémie microcytaire hypochrome. Citer un signe clinique d'alerte chez l'enfant exposé.
    \item \textbf{(Doc 2)} Pourquoi la recherche d'inhibiteurs spécifiques de la protéase virale NS2B-NS3 est-elle une stratégie antivirale prometteuse ? Discuter de la notion de « fenêtre thérapeutique » et de toxicité pour l'hôte.
    \item \textbf{(Doc 3)} Pour la clarification des jus de fruits (ex : ananas, mangue), on utilise des pectinases. Expliquer le substrat de cette enzyme (structure de la pectine) et le mécanisme d'action (hydrolyse) qui conduit à la baisse de viscosité et à la clarification.
    \item \textbf{Transversal} : Citer deux mesures de prévention primaire (niveau population) et deux mesures de prévention secondaire (niveau individuel) pour limiter l'impact sanitaire des inhibiteurs enzymatiques environnementaux (métaux lourds, pesticides organophosphorés) au Cameroun.
\end{enumerate}
\end{exercice}



\newpage

\coursec{Corrigés des exercices}

\coursec{Corrigés des exercices d'application — Leçon ND04}

\begin{corrige}[Exercice 1 — QCM : Enzyme, substrat et catalyse]
\textbf{Bonne réponse : 3.} Une enzyme abaisse l'énergie d'activation d'une réaction chimique et accélère l'atteinte de l'équilibre.

\textbf{Analyse détaillée de chaque proposition :}
\begin{enumerate}
    \item \textbf{Fausse.} Une enzyme est une \cle{protéine} (ou un ARN catalytique dans le cas des ribozymes), jamais un glucide. C'est la confusion classique entre le catalyseur (l'enzyme) et le substrat (qui peut être un glucide, comme l'amidon).
    \item \textbf{Fausse.} Le substrat se fixe sur le \cle{site actif} de l'enzyme pour former le complexe enzyme-substrat (ES). Le site allostérique est un site \emph{distinct} du site actif, réservé aux effecteurs (activateurs ou inhibiteurs) qui modulent l'activité enzymatique.
    \item \textbf{Exacte.} L'enzyme stabilise l'état de transition et abaisse l'\cle{énergie d'activation} $E_a$. D'après la loi d'Arrhenius ($k = A\,e^{-E_a/RT}$), la constante de vitesse $k$ augmente fortement : la réaction atteint son équilibre beaucoup plus vite. \textbf{Point essentiel :} l'enzyme ne modifie ni l'équilibre final, ni le bilan énergétique global de la réaction ; elle accélère uniquement la cinétique.
    \item \textbf{Fausse.} La dénaturation thermique est le plus souvent \emph{irréversible} : les liaisons faibles (ponts hydrogène, interactions hydrophobes) sont rompues et la protéine coagule (exemple : le blanc d'œuf cuit ne redevient jamais liquide en refroidissant).
\end{enumerate}
\end{corrige}

\begin{attention}[Point de vigilance]
Au Probatoire, la formulation « accélère l'atteinte de l'équilibre \emph{sans le déplacer} » est la justification attendue. Dire simplement « l'enzyme accélère la réaction » sans mentionner l'énergie d'activation est insuffisant.
\end{attention}

\begin{corrige}[Exercice 2 — Vrai ou Faux justifié]
\begin{enumerate}
    \item \textbf{VRAI.} La spécificité repose sur la \cle{complémentarité} géométrique (forme) et chimique (charges, groupements polaires) entre le site actif et le substrat : c'est le modèle « clé-serrure » d'Emil Fischer, affiné par le modèle de l'« ajustement induit » de Koshland (le site actif se moule légèrement sur le substrat lors de la fixation).
    \item \textbf{FAUX.} La vitesse initiale suit une cinétique de type Michaelis-Menten : $v_0 = \dfrac{V_{\text{max}}\,[S]}{K_m + [S]}$. Quand $[S]$ devient très grande, tous les sites actifs sont saturés en permanence et la vitesse plafonne à $V_{\text{max}}$ : la courbe est une hyperbole avec asymptote horizontale, et non une droite indéfiniment croissante.
    \item \textbf{VRAI.} La pepsine est sécrétée dans l'estomac où le suc gastrique (acide chlorhydrique) maintient un pH de 1,5 à 2. Son pH optimal est donc très acide ($\approx$ 2) : c'est une adaptation de l'enzyme à son milieu de fonctionnement.
    \item \textbf{FAUX.} Un inhibiteur compétitif se fixe sur le site actif de manière \emph{réversible} (liaisons non covalentes) et entre en compétition avec le substrat : un excès de substrat peut le déloger. La fixation irréversible caractérise les inhibiteurs irréversibles (liaison covalente), comme l'aspirine sur la cyclo-oxygénase ou les organophosphorés sur l'acétylcholinestérase.
\end{enumerate}
\end{corrige}

\begin{attention}[Point de vigilance]
Ne pas confondre « compétitif » (lieu de fixation : le site actif) et « irréversible » (nature de la liaison : covalente). Ce sont deux critères de classification indépendants : un inhibiteur compétitif classique est réversible.
\end{attention}

\begin{corrige}[Exercice 3 — Définitions et notions clés]
\begin{enumerate}
    \item \textbf{Enzyme :} catalyseur biologique de nature protéique (ou parfois ribonucléique) qui accélère considérablement une réaction chimique spécifique en abaissant son énergie d'activation, sans être consommée ni modifiée de façon permanente au cours de la réaction, et sans modifier l'équilibre final.
    \item \textbf{Site actif :} région tridimensionnelle de l'enzyme, en creux ou en fente, formée par quelques acides aminés rapprochés grâce au repliement de la chaîne polypeptidique, sur laquelle se fixe le substrat et où s'accomplit la catalyse.
    \item \textbf{Spécificité enzymatique :} propriété d'une enzyme de ne reconnaître qu'un substrat déterminé (\cle{spécificité de substrat} — exemple : l'uréase n'hydrolyse que l'urée) et de ne catalyser qu'un seul type de réaction chimique (\cle{spécificité de réaction} — exemple : une lipase catalyse une hydrolyse, jamais une oxydation).
    \item \textbf{Dénaturation enzymatique :} destruction de la structure spatiale native de l'enzyme (structures secondaire, tertiaire et quaternaire) par rupture des liaisons faibles, sous l'effet de la chaleur, d'un pH extrême ou de certains agents chimiques. Elle entraîne la déformation du site actif et la perte, le plus souvent irréversible, de l'activité catalytique.
    \item \textbf{Vitesse maximale $V_{\text{max}}$ :} vitesse de la réaction lorsque tous les sites actifs de l'enzyme sont saturés en substrat ($[S] \gg K_m$) ; elle dépend de la concentration en enzyme. \textbf{Constante de Michaelis $K_m$ :} concentration en substrat pour laquelle $v_0 = V_{\text{max}}/2$ ; elle traduit l'affinité de l'enzyme pour son substrat : plus $K_m$ est faible, plus l'affinité est grande.
\end{enumerate}
\end{corrige}

\begin{attention}[Point de vigilance]
$K_m$ est une concentration (unité : mol·L$^{-1}$ ou mmol·L$^{-1}$), pas une vitesse. Au Probatoire, on exige souvent la phrase : « $K_m$ traduit l'affinité : $K_m$ faible = affinité forte ».
\end{attention}

\begin{corrige}[Exercice 4 — Calcul direct : rendement catalytique ($k_{\text{cat}}$)]
\textbf{Étape 1 : exprimer $V_{\text{max}}$ en mol·L$^{-1}$·s$^{-1}$.}

La vitesse mesurée est de \SI{120}{\micro\mol\per\minute} dans \SI{3}{\mL} :
\begin{align*}
V_{\text{max}} &= \frac{\SI{120}{\micro\mol\per\minute}}{\SI{3}{\mL}} = \SI{40}{\micro\mol\per\minute\per\mL} \\
&= \SI{40e3}{\micro\mol\per\minute\per\liter} = \SI{4e-2}{\mol\per\liter\per\minute}
\end{align*}
Conversion en secondes :
\[ V_{\text{max}} = \frac{\SI{4e-2}{\mol\per\liter\per\minute}}{60} \approx \SI{6,67e-4}{\mol\per\liter\per\second} \]

\textbf{Étape 2 : appliquer la formule du turnover.}
\begin{align*}
k_{\text{cat}} &= \frac{V_{\text{max}}}{[E]_{\text{total}}} = \frac{\SI{6,67e-4}{\mol\per\liter\per\second}}{\SI{2e-8}{\mol\per\liter}} \\
k_{\text{cat}} &\approx \textbf{3,3} \times \textbf{10}^{\textbf{4}}\ \textbf{s}^{\textbf{-1}}
\end{align*}

\textbf{Vérification (chemin alternatif) :} quantité d'enzyme totale = $\SI{2e-8}{\mol\per\liter} \times \SI{3e-3}{\liter} = \SI{6e-11}{\mol}$. Nombre de tours par minute = $\frac{\SI{120e-6}{\mol\per\minute}}{\SI{6e-11}{\mol}} = \SI{2e6}{\per\minute}$. $k_{\text{cat}} = \dfrac{\SI{2e6}{\per\minute}}{60} \approx \SI{3,3e4}{\per\second}$. Les deux calculs concordent.

\textbf{Interprétation :} chaque molécule d'enzyme transforme environ \num{33000} molécules de substrat par seconde : cela illustre l'extraordinaire efficacité de la catalyse enzymatique (l'enzyme n'est pas consommée et « tourne » en permanence).
\end{corrige}

\begin{attention}[Point de vigilance]
Deux pièges classiques : (1) oublier de diviser par le volume (\SI{3}{\mL}) pour obtenir une vitesse \emph{par litre} ; (2) oublier de convertir les minutes en secondes. Toujours vérifier l'homogénéité des unités : $k_{\text{cat}}$ s'exprime en s$^{-1}$ (mol de substrat / mol d'enzyme / seconde).
\end{attention}

\begin{corrige}[Exercice 5 — Exploitation de courbe : activité d'une amylase en fonction de la température]
\begin{enumerate}
    \item \textbf{Lecture graphique :} la vitesse initiale est maximale au sommet de la courbe, soit à une \textbf{température optimale de \SI{37}{\celsius}} ($v_0 \approx 100$ UA). Cette valeur correspond à la température corporelle humaine, milieu naturel de l'amylase salivaire.
    \item \textbf{Entre \SI{0}{\celsius} et \SI{37}{\celsius} (montée de la courbe) — phénomène cinétique :} l'élévation de température augmente l'agitation thermique des molécules, donc la fréquence et l'énergie des chocs entre l'enzyme et le substrat. Le nombre de complexes enzyme-substrat formés par unité de temps augmente, et la vitesse de réaction croît (effet général de la température sur toute réaction chimique, décrit par la loi d'Arrhenius).
    \item \textbf{Au-delà de \SI{37}{\celsius} (descente brutale) — phénomène structural :} l'agitation thermique devient excessive et rompt les liaisons faibles (ponts hydrogène, interactions hydrophobes et ioniques) qui stabilisent la structure tertiaire de l'enzyme. Le site actif se déforme : c'est la \cle{dénaturation thermique}. L'activité s'effondre et devient quasi nulle vers \SI{60-65}{\celsius}.
    \item \textbf{Prédiction :} après 10 minutes à \SI{80}{\celsius}, l'enzyme replacée à \SI{37}{\celsius} reste \textbf{totalement inactive}. \textbf{Justification :} la dénaturation thermique est \emph{irréversible} : le retour à la température optimale ne permet pas à la protéine de retrouver spontanément son repliement natif (analogie : le blanc d'œuf coagulé ne redevient pas liquide en refroidissant). Cette expérience prouve que la chute d'activité à haute température n'est pas un simple ralentissement cinétique mais une destruction structurale.
\end{enumerate}
\end{corrige}

\begin{attention}[Point de vigilance]
Erreur fréquente au Probatoire : expliquer la descente de la courbe par « les molécules vont trop vite ». La justification attendue est la \textbf{dénaturation de la protéine} (rupture des liaisons faibles, perte du site actif). Bien distinguer les deux versants : versant ascendant = effet cinétique ; versant descendant = effet structural.
\end{attention}

\begin{corrige}[Exercice 6 — Interprétation d'expérience : spécificité et thermosensibilité de l'amylase]
\begin{enumerate}
    \item \textbf{Tube 1 :} l'eau iodée reste jaune $\Rightarrow$ l'amidon a \emph{disparu} ; la liqueur de Fehling donne un précipité rouge brique $\Rightarrow$ des \emph{sucres réducteurs} sont apparus. L'amylase salivaire a donc hydrolysé l'amidon en sucres simples (maltose puis glucose). Équation bilan simplifiée :
    \[ \ce{(C6H10O5)_{n} + n H2O ->[amylase] n C6H12O6} \]
    (amidon + eau $\rightarrow$ glucose, le maltose étant l'intermédiaire principal).
    \item \textbf{Tube 2 :} la liqueur de Fehling reste bleue $\Rightarrow$ aucun sucre réducteur n'a été produit : la cellulose n'a \emph{pas} été hydrolysée par l'amylase. L'eau iodée, spécifique de l'amidon, reste jaune/orange (pas de coloration bleu-noir) en présence de cellulose. \textbf{Conclusion :} l'amylase est \cle{spécifique} de son substrat : elle hydrolyse les liaisons $\alpha$-1,4 de l'amidon mais est incapable d'hydrolyser les liaisons $\beta$-1,4 de la cellulose, bien que les deux soient des polymères de glucose. C'est la \textbf{spécificité de substrat}.
    \item \textbf{Tube 3 :} amidon toujours présent (eau iodée bleu-noir) et aucun sucre réducteur (Fehling bleu) : la salive bouillie a perdu toute activité. \textbf{Conclusion :} l'ébullition a \cle{dénaturé} l'amylase de façon irréversible, ce qui prouve la \textbf{nature protéique} de l'enzyme et sa \textbf{thermosensibilité}.
    \item \textbf{Deux propriétés fondamentales mises en évidence :}
    \begin{itemize}
        \item la \textbf{spécificité enzymatique} (comparaison tubes 1 et 2 : amidon hydrolysé, cellulose non hydrolysée) ;
        \item la \textbf{thermosensibilité} liée à la nature protéique des enzymes (comparaison tubes 1 et 3 : salive fraîche active, salive bouillie inactive).
    \end{itemize}
\end{enumerate}
\end{corrige}

\begin{attention}[Point de vigilance]
La démarche d'interprétation exigée suit toujours le schéma : observation (résultat des tests) $\rightarrow$ déduction (présence/absence des molécules) $\rightarrow$ conclusion (propriété de l'enzyme). Le tube 3 est le \textbf{tube témoin} : il prouve que c'est bien l'enzyme (et non un composé chimique quelconque de la salive) qui réalise l'hydrolyse. Attention : l'eau iodée ne colore pas la cellulose (reste jaune), le bleu-noir n'apparaît que pour l'amidon.
\end{attention}

\begin{corrige}[Exercice 7 — Comparaison de courbes d'activité en fonction du pH]
\begin{enumerate}
    \item \textbf{Identification :}
    \begin{itemize}
        \item \textbf{Enzyme A = pepsine} : pH optimal $\approx$ 2, adapté au milieu très acide de l'estomac (pH 1,5--3).
        \item \textbf{Enzyme B = amylase salivaire} : pH optimal $\approx$ 6,8, proche de la neutralité, correspondant au milieu buccal (pH 6,5--7,5).
        \item \textbf{Enzyme C = trypsine} : pH optimal $\approx$ 8, adapté au milieu légèrement alcalin de l'intestin grêle (pH 7,5--8,5), où le suc pancréatique neutralise l'acidité gastrique.
    \end{itemize}
    Le pH optimal de chaque enzyme correspond au pH de son milieu naturel de fonctionnement : c'est une remarquable adaptation fonctionnelle.
    \item \textbf{Explication moléculaire :} le site actif comporte des acides aminés à chaînes latérales \emph{ionisables} (acide aspartique, acide glutamique, histidine, lysine...). Leur état de protonation (charge positive, négative ou neutre) dépend du pH et conditionne à la fois la fixation du substrat (complémentarité de charges) et le mécanisme catalytique (transferts de protons). Tout écart par rapport au pH optimal modifie l'ionisation de ces groupements : la complémentarité enzyme-substrat est perdue et l'activité chute. Aux pH extrêmes, les liaisons ioniques et ponts hydrogène de toute la protéine sont rompus : c'est la \cle{dénaturation}, généralement irréversible.
    \item \textbf{Pepsine placée à pH 8 :} elle serait rapidement et \textbf{irréversiblement inactivée}. Le pH alcalin modifie l'état de charge de ses groupements ionisables et rompt ses liaisons faibles (dénaturation alcaline) : son site actif se déforme et elle ne peut plus hydrolyser les protéines. C'est pourquoi la digestion protéique dans le duodénum est assurée par d'autres enzymes (trypsine, chymotrypsine) dont le pH optimal est alcalin.
\end{enumerate}
\end{corrige}

\begin{attention}[Point de vigilance]
Ne pas écrire « le pH tue l'enzyme » : une enzyme n'est pas un être vivant. La formulation correcte est « le pH dénature l'enzyme » ou « modifie l'ionisation du site actif ». Préciser aussi que l'effet du pH est réversible pour de petits écarts, irréversible aux pH extrêmes.
\end{attention}

\begin{corrige}[{Exercice 8 — Construction et interprétation : $v_0 = f({[S]})$}]
\begin{enumerate}
    \item \textbf{Allure de la courbe :} en reportant $v_0$ en ordonnée et $[S]$ en abscisse, on obtient une courbe partant de l'origine, d'abord fortement croissante, puis s'infléchissant et tendant vers un \textbf{plateau} : c'est une \cle{hyperbole} de Michaelis-Menten, avec une asymptote horizontale correspondant à $V_{\text{max}}$.
    \item \textbf{Détermination des paramètres :} la vitesse plafonne vers \SI{89-95}{\micro\mol\per\minute\per\mL} aux fortes concentrations ; on estime $V_{\text{max}} \approx \textbf{\SI{100}{\micro\mol\per\minute\per\mL}}$. La demi-vitesse maximale vaut $V_{\text{max}}/2 = \SI{50}{\micro\mol\per\minute\per\mL}$, valeur atteinte pour $[S] = \SI{4,0}{\mmol\per\liter}$ (donnée du tableau). Donc :
    \[ K_m \approx \textbf{\SI{4,0}{\mmol\per\liter}} \]
    \item \textbf{Explication par le modèle du complexe enzyme-substrat :}
    \begin{itemize}
        \item À \emph{faible} $[S]$ : les sites actifs sont majoritairement libres ; chaque molécule de substrat supplémentaire trouve facilement une enzyme disponible, donc $v_0$ est quasi proportionnelle à $[S]$ (cinétique d'ordre 1).
        \item À $[S]$ \emph{intermédiaire} : de plus en plus de sites sont occupés ; la probabilité de trouver un site libre diminue, la courbe s'infléchit.
        \item À \emph{forte} $[S]$ ($[S] \gg K_m$) : tous les sites actifs sont saturés en permanence ; l'enzyme travaille à plein régime et $v_0 = V_{\text{max}}$, indépendante de $[S]$ (cinétique d'ordre 0). C'est le phénomène de \cle{saturation}.
    \end{itemize}
    \item \textbf{Calculs :} d'après l'équation de Michaelis-Menten $v_0 = \dfrac{V_{\text{max}}\,[S]}{K_m + [S]}$ :
    \begin{align*}
    [S] = K_m : \quad v_0 &= \frac{V_{\text{max}} \times K_m}{K_m + K_m} = \frac{V_{\text{max}}}{2} = \textbf{\SI{50}{\micro\mol\per\minute\per\mL}} \\
    [S] = 10\,K_m : \quad v_0 &= \frac{V_{\text{max}} \times 10\,K_m}{K_m + 10\,K_m} = \frac{10}{11}\,V_{\text{max}} \approx 0,909 \times 100 \approx \textbf{\SI{91}{\micro\mol\per\minute\per\mL}}
    \end{align*}
    \textbf{Commentaire :} à $[S] = 10\,K_m$, l'enzyme fonctionne déjà à plus de 90 \% de sa capacité maximale ; multiplier encore la concentration en substrat n'augmenterait quasiment plus la vitesse. C'est la conséquence directe de la saturation des sites actifs.
\end{enumerate}
\end{corrige}

\begin{attention}[Point de vigilance]
$K_m$ se lit toujours sur l'\textbf{axe des abscisses} (c'est une concentration) à la hauteur de $V_{\text{max}}/2$ — ne jamais le lire sur l'axe des vitesses. Vérifier la cohérence avec le tableau : ici $v_0 = 50$ correspond exactement à $[S] = \SI{4,0}{\mmol\per\liter}$, ce qui confirme l'estimation.
\end{attention}

\begin{corrige}[Exercice 9 — Cas clinique : intolérance au lactose du nourrisson]
\begin{enumerate}
    \item \textbf{Équation d'hydrolyse du lactose :}
    \[ \ce{C12H22O11 + H2O ->[lactase] C6H12O6 + C6H12O6} \]
    (lactose + eau $\rightarrow$ glucose + galactose).
    \item \textbf{Mécanisme physiopathologique :} le déficit en lactase (\SI{0,2}{U.g^{-1}} au lieu de \SIrange{10}{30}{U.g^{-1}}) empêche l'hydrolyse du \cle{substrat} (le lactose du lait) dans l'intestin grêle. Le lactose non digéré s'\textbf{accumule} dans la lumière intestinale : sa présence augmente la \textbf{pression osmotique}, ce qui attire l'eau du sang vers l'intestin $\Rightarrow$ diarrhée aqueuse. Arrivé intact dans le côlon, le lactose est dégradé par \textbf{fermentation} par la \textbf{flore colique} bactérienne, qui produit des gaz (H$_2$, CO$_2$, CH$_4$) et des acides organiques $\Rightarrow$ ballonnements, crampes, et hydrogène absorbé dans le sang puis \textbf{expiré} (d'où le test H$_2$ très positif à \SI{45}{ppm}). La malabsorption chronique explique la cassure de la courbe de poids.
    \item \textbf{Activité sucrase-isomaltase normale :} cela montre que la bordure en brosse des entérocytes est globalement intacte et que le déficit est \textbf{sélectif}, limité à la lactase. Il s'agit donc d'un déficit d'origine \textbf{génétique} (déficit congénital en lactase) et non d'une cause acquise généralisée (gastro-entérite, malnutrition, atrophie villositaire), qui aurait abaissé l'activité de \emph{toutes} les disaccharidases.
    \item \textbf{Mesures adaptées :}
    \begin{itemize}
        \item \textbf{Substitution du lait par un lait sans lactose} (lait hydrolysé industriellement) ou une préparation adaptée : on supprime le substrat que l'enzyme déficiente ne peut pas traiter, ce qui fait disparaître les symptômes tout en couvrant les besoins du nourrisson.
        \item \textbf{Supplémentation en lactase exogène} (gouttes de lactase ajoutées au lait avant la tétée) : l'enzyme apportée réalise l'hydrolyse du lactose à la place de la lactase intestinale déficiente — c'est un traitement de substitution enzymatique.
    \end{itemize}
    \item \textbf{Différence mécanistique :} l'\textbf{intolérance au lactose} est un trouble \emph{enzymatique et métabolique} (déficit en lactase $\Rightarrow$ malabsorption) ; les symptômes sont purement digestifs et proportionnels à la quantité ingérée, sans danger vital. L'\textbf{allergie aux protéines de lait de vache} est une réaction \emph{immunitaire} (souvent IgE-médiée) dirigée contre des protéines (caséine, $\beta$-lactoglobuline) : elle peut provoquer des signes cutanés (urticaire, eczéma), respiratoires et digestifs, déclenchés par des quantités infimes, avec un risque de choc anaphylactique. Aucun déficit enzymatique n'est en cause.
\end{enumerate}
\end{corrige}

\begin{attention}[Point de vigilance]
La confusion intolérance/allergie est un piège classique du Probatoire : retenir le couple « déficit enzymatique $\leftrightarrow$ réponse immunitaire ». Autre point attendu : relier explicitement chaque symptôme à une étape du mécanisme (osmose pour la diarrhée, fermentation bactérienne pour les gaz et le H$_2$).
\end{attention}

\begin{corrige}[Exercice 10 — Étude expérimentale : inhibition de l'ECA et thérapeutique]
\begin{enumerate}
    \item \textbf{Type d'inhibition : compétitive.} \textbf{Justification :} dans la représentation de Lineweaver-Burk, les trois droites se coupent sur l'\textbf{axe des ordonnées} (même ordonnée à l'origine $1/V_{\text{max}} = 0,1$). Donc :
    \begin{itemize}
        \item $V_{\text{max}}$ \textbf{reste constante} (\SI{10}{\micro\mol\per\minute\per\mL}) : à concentration de substrat saturante, le substrat déloge totalement l'inhibiteur ;
        \item les pentes augmentent avec $[I]$ (0,1 $\to$ 0,2 $\to$ 0,35), donc le $K_m$ \textbf{apparent augmente} : l'affinité apparente de l'enzyme pour le substrat diminue.
    \end{itemize}
    Ce double critère ($V_{\text{max}}$ inchangée, $K_m$ apparent augmenté, intersection sur l'axe des $y$) est la signature cinétique de l'inhibition compétitive.
    \item \textbf{Mécanisme moléculaire :} l'inhibiteur possède une structure chimique analogue à celle du substrat (l'angiotensine I) et se fixe de façon \emph{réversible} (liaisons non covalentes) sur le \textbf{site actif} de l'ECA. Substrat et inhibiteur sont en \textbf{compétition} pour le même site : $E + S \rightleftharpoons ES \rightarrow E + P$ et $E + I \rightleftharpoons EI$ (complexe inactif). Une forte concentration de substrat déplace l'équilibre vers ES et l'emporte sur l'inhibiteur : c'est pourquoi $V_{\text{max}}$ peut être atteinte malgré la présence de l'inhibiteur.
    \item \textbf{Calcul des paramètres du témoin :} la droite témoin a pour équation $1/v_0 = 0,1 + 0,1 \times (1/[S])$.
    \begin{itemize}
        \item Ordonnée à l'origine : $\dfrac{1}{V_{\text{max}}} = 0,1$ min·mL·µmol$^{-1}$ $\Rightarrow$ $V_{\text{max}} = \dfrac{1}{0,1} = \textbf{\SI{10}{\micro\mol\per\minute\per\mL}}$.
        \item Abscisse à l'origine : $0 = 0,1 + 0,1 \times \dfrac{1}{[S]}$ $\Rightarrow$ $\dfrac{1}{[S]} = -1$ L·mmol$^{-1}$. Or cette abscisse vaut $-\dfrac{1}{K_m}$, donc $K_m = \textbf{\SI{1}{\mmol\per\liter}}$.
    \end{itemize}
    \item \textbf{Effet de $[I]$ sur la pente :} en inhibition compétitive, la pente vaut $\dfrac{K_m}{V_{\text{max}}}\left(1 + \dfrac{[I]}{K_i}\right)$. Quand $[I]$ augmente ($[I_2] > [I_1]$), la pente augmente proportionnellement (0,1 $\to$ 0,2 $\to$ 0,35) alors que l'ordonnée à l'origine ne bouge pas. La \textbf{constante d'inhibition $K_i$} (constante de dissociation du complexe EI) se déduit du rapport des pentes : $\dfrac{\text{pente}_I}{\text{pente}_0} = 1 + \dfrac{[I]}{K_i}$. Plus $K_i$ est faible, plus l'inhibiteur est puissant (forte affinité pour l'enzyme).
    \item \textbf{Efficacité clinique :} même si la concentration en angiotensine I varie physiologiquement, l'inhibiteur compétitif réduit à tout moment la \emph{vitesse de formation} de l'angiotensine II, car à concentration d'inhibiteur fixée, la fraction d'enzyme bloquée sous forme EI reste importante ($v_i = \dfrac{V_{\text{max}}[S]}{K_m(1+[I]/K_i) + [S]} < v_0$ pour toute valeur de $[S]$). En maintenant par le dosage médicamenteux $[I] \gg K_i$, on obtient une inhibition forte et constante : la baisse durable d'angiotensine II entraîne une vasodilatation et une réduction de la rétention sodée (moins d'aldostérone), donc une \textbf{baisse de la tension artérielle}. C'est le principe des IEC (inhibiteurs de l'enzyme de conversion) comme l'énalapril, largement prescrits contre l'hypertension, y compris au Cameroun.
\end{enumerate}
\end{corrige}

\begin{attention}[Point de vigilance]
Dans un graphique de Lineweaver-Burk, bien mémoriser les signatures : intersection sur l'axe des ordonnées = compétitive ($V_{\text{max}}$ constante, $K_m$ apparent augmenté) ; intersection sur l'axe des abscisses = non compétitive pure ($K_m$ constant, $V_{\text{max}}$ diminuée). Attention également au signe négatif de l'abscisse à l'origine ($-1/K_m$).
\end{attention}

\begin{corrige}[Exercice 11 — Synthèse : enzymes, environnement et santé publique au Cameroun]
\begin{enumerate}
    \item \textbf{(Doc 1) Mécanisme d'inhibition par Hg$^{2+}$ :} l'ion mercurique présente une très forte affinité pour les groupements thiol (-SH) des cystéines. Il y forme une liaison quasi \textbf{covalente} (Hg-S) stable, qui ne se rompt ni par dilution ni par dialyse : l'inhibition est donc \textbf{irréversible}. Si la cystéine bloquée appartient au site actif ou maintient la structure tertiaire, le site actif est détruit ou déformé $\Rightarrow$ perte d'activité. La fixation se faisant indifféremment sur l'enzyme libre ou le complexe ES, il s'agit d'une inhibition de type \textbf{non compétitive irréversible}. \textbf{Conséquences cinétiques :} la quantité d'enzyme active diminue de façon définitive, donc \textbf{$V_{\text{max}}$ diminue} ; en revanche les molécules d'enzyme restées intactes conservent leur affinité normale pour le substrat, donc \textbf{$K_m$ reste inchangé}.
    \item \textbf{(Doc 1) Plomb et anémie :} la $\delta$-ALA déshydratase catalyse une étape clé de la biosynthèse de l'\textbf{hème}, partie active de l'hémoglobine. Son inhibition par le plomb bloque la production d'hème $\Rightarrow$ les globules rouges sont produits avec moins d'hémoglobine : ils sont petits (\textbf{microcytaires}) et pâles (\textbf{hypochromes}) $\Rightarrow$ anémie microcytaire hypochrome, avec pâleur, fatigue et retard de développement. De plus, le $\delta$-ALA non transformé s'accumule et exerce une neurotoxicité. \textbf{Signe d'alerte chez l'enfant exposé :} la \textbf{ligne de Burton} (liseré bleu-noir au bord des gencives), des coliques abdominales ou des troubles neurologiques (irritabilité, troubles de l'attention, encéphalopathie dans les cas graves).
    \item \textbf{(Doc 2) Stratégie antivirale :} la protéase NS2B-NS3 est \emph{indispensable} à la maturation des protéines du flavivirus (clivage de la polyprotéine virale) et n'a \emph{pas d'équivalent} chez l'humain. Un inhibiteur spécifique bloque donc la réplication virale sans perturber les enzymes de l'hôte : c'est une cible thérapeutique idéale. La \textbf{fenêtre thérapeutique} (marge entre dose efficace et dose toxique) sera d'autant plus large que l'inhibiteur est sélectif de la protéase virale et épargne les protéases humaines. La \textbf{toxicité} éventuelle proviendrait d'une inhibition « hors cible » de protéases humaines (effets hépatiques ou rénaux) : d'où la nécessité d'essais cliniques rigoureux avant utilisation.
    \item \textbf{(Doc 3) Pectinases et clarification :} la \textbf{pectine} est un polysaccharide de la paroi végétale, formé de chaînes d'\textbf{acide galacturonique} reliées par des liaisons $\alpha$-1,4, partiellement méthylestérifiées ; elle retient l'eau et forme un gel qui rend le jus visqueux et trouble. Les pectinases \textbf{hydrolysent} ces liaisons glycosidiques $\alpha$-1,4 (action de la pectinestérase qui déméthyle, puis de la polygalacturonase qui fragmente la chaîne). La fragmentation du polymère provoque la \textbf{chute de la viscosité} et la précipitation des particules en suspension $\Rightarrow$ le jus devient limpide (\textbf{clarification}) et le rendement d'extraction augmente (application : jus d'ananas et de mangue des unités agroalimentaires camerounaises).
    \item \textbf{Mesures de prévention :}
    \begin{itemize}
        \item \textbf{Prévention primaire (population) :} (1) réglementer et contrôler les rejets miniers et industriels (exploitations artisanales de Bétaré-Oya, fonderies) pour limiter la contamination des eaux et des sols par Pb, Hg, Cd ; (2) organiser la collecte et le recyclage des batteries usagées et encadrer (voire interdire) les pesticides organophosphorés les plus toxiques, avec information des populations.
        \item \textbf{Prévention secondaire (individu) :} (1) dépistage biologique ciblé des populations exposées (plombémie chez les enfants des zones minières, activité cholinestérasique chez les agriculteurs utilisant des pesticides) ; (2) port obligatoire d'équipements de protection individuelle (gants, masques) et mesures d'hygiène (lavage des mains, ne pas manger ni fumer sur le lieu de travail, laver les vêtements de travail séparément).
    \end{itemize}
\end{enumerate}
\end{corrige}

\begin{attention}[Point de vigilance]
Pour les inhibiteurs irréversibles, la justification clé est : « la concentration d'enzyme \emph{active} diminue » $\Rightarrow$ $V_{\text{max}}$ diminue, $K_m$ inchangé. Ne pas écrire que le métal « entre en compétition avec le substrat » : il ne ressemble pas au substrat et agit sur les thiols, pas sur le site de fixation du substrat. En santé publique, bien distinguer prévention \emph{primaire} (agir sur la cause, avant l'exposition) et \emph{secondaire} (dépistage précoce et protection des exposés).
\end{attention}

\newpage

\coursec{Fiche bilan}

\begin{popfigure}
\begin{tikzpicture}[mindmap, grow cyclic, every node/.style={concept, font=\small},
root concept/.append style={concept color=popblue, font=\small\bfseries, text=white},
level 1/.append style={level distance=3.4cm, sibling angle=60},
level 2/.append style={level distance=2.2cm, font=\scriptsize}]
\node[root concept]{Enzymes et renouvellement moléculaire}
  child[concept color=poporange] { node {Notion d'enzyme}
    child { node {Catalyseur biologique (protéine)} }
    child { node {Substrat + produit} }
  }
  child[concept color=popgreen] { node {Spécificité}
    child { node {Site actif} }
    child { node {Modèle clé--serrure} }
  }
  child[concept color=popblue!70] { node {Température}
    child { node {Optimum $\approx$ \SI{37}{\celsius}} }
    child { node {Dénaturation à chaud} }
  }
  child[concept color=poppurple] { node {pH et concentrations}
    child { node {pH optimum propre à chaque enzyme} }
    child { node {Saturation en substrat} }
  }
  child[concept color=popgold] { node {Rôle métabolique}
    child { node {Synthèses (anabolisme)} }
    child { node {Dégradations (catabolisme)} }
  }
  child[concept color=poppink] { node {Déficits et pathologies}
    child { node {Intolérance au lactose} }
    child { node {Phénylcétonurie, albinisme} }
  };
\end{tikzpicture}
\legende{Carte mentale de la leçon : les six idées forces sur les enzymes, catalyseurs du renouvellement moléculaire de la cellule.}
\end{popfigure}

\begin{bilan}
\textbf{Définitions clés}
\begin{itemize}
  \item \cle{Enzyme} : protéine \cle{catalyseur biologique} qui accélère une réaction chimique de la cellule sans être consommée ni modifiée à la fin de la réaction.
  \item \cle{Substrat} : molécule sur laquelle agit l'enzyme ; le \cle{produit} est la molécule obtenue après la réaction.
  \item \cle{Site actif} : région de l'enzyme, de forme complémentaire du substrat, où se fixe le substrat et où se déroule la réaction.
  \item \cle{Spécificité enzymatique} : une enzyme donnée ne catalyse qu'une seule réaction (ou un seul type de réaction) sur un substrat précis.
  \item \cle{Dénaturation} : destruction irréversible de la structure spatiale de l'enzyme (par forte chaleur ou pH extrême) entraînant la perte de son activité.
\end{itemize}

\textbf{Ce qu'il faut connaître par c\oe ur}
\begin{center}
\begin{tabularx}{\linewidth}{|l|X|}\hline
\rowcolor{popblueL} \textbf{Notion} & \textbf{À retenir} \\ \hline
Bilan d'une réaction enzymatique & Enzyme $+$ Substrat $\rightarrow$ complexe enzyme--substrat $\rightarrow$ Enzyme $+$ Produit \\ \hline
Température & Activité nulle à froid (ralentie, réversible), maximale vers \SI{37}{\celsius} (enzymes humaines), nulle au-delà de $\sim$\SI{60}{\celsius} (dénaturation irréversible) \\ \hline
pH & Chaque enzyme possède un \cle{pH optimum} : pepsine $\approx 2$ (estomac), amylase salivaire $\approx 7$, trypsine $\approx 8$ \\ \hline
Concentration en substrat & La vitesse croît puis plafonne : \cle{saturation} de tous les sites actifs \\ \hline
Concentration en enzyme & À substrat en excès, la vitesse est proportionnelle à la quantité d'enzyme \\ \hline
Rôle métabolique & Les enzymes catalysent les réactions de \cle{dégradation} (catabolisme, ex. digestion) et de \cle{synthèse} (anabolisme, ex. synthèse des protéines) assurant le renouvellement moléculaire \\ \hline
\end{tabularx}
\end{center}

\textbf{Méthodes express}
\begin{enumerate}
  \item \textbf{Exploiter une courbe activité $=f(T)$ ou $f(\mathrm{pH})$} : repérer le maximum (optimum), justifier la montée (agitation moléculaire) et la chute (dénaturation ou déformation du site actif).
  \item \textbf{Prouver une spécificité} : comparer l'activité d'une même enzyme sur deux substrats différents (ex. amylase $+$ amidon vs amylase $+$ glucose : révélation à l'eau iodée ou à la liqueur de Fehling).
  \item \textbf{Relier déficit et pathologie} : déficit enzymatique $\rightarrow$ substrat non transformé qui s'accumule (ou produit manquant) $\rightarrow$ symptômes (ex. déficit en lactase $\rightarrow$ intolérance au lactose).
\end{enumerate}
\end{bilan}

\begin{aretenir}[Checklist avant le Probatoire]
\begin{itemize}
  \item[$\square$] Je sais définir enzyme, substrat, produit et site actif sans hésiter.
  \item[$\square$] Je sais qu'une enzyme est une protéine et qu'elle n'est pas consommée par la réaction.
  \item[$\square$] Je sais expliquer la spécificité enzymatique par le modèle clé--serrure (complémentarité de forme).
  \item[$\square$] Je sais décrire et interpréter la courbe de l'activité enzymatique en fonction de la température.
  \item[$\square$] Je distingue inhibition par le froid (réversible) et dénaturation par la chaleur (irréversible).
  \item[$\square$] Je connais le pH optimum de la pepsine ($\approx 2$), de l'amylase salivaire ($\approx 7$) et de la trypsine ($\approx 8$).
  \item[$\square$] Je sais interpréter l'effet de la concentration en substrat (saturation) et en enzyme.
  \item[$\square$] Je sais citer un exemple d'enzyme de dégradation (amylase, pepsine, lipase) et son rôle dans la digestion.
  \item[$\square$] Je sais relier un déficit enzymatique à une pathologie : lactase/intolérance au lactose, phénylalanine hydroxylase/phénylcétonurie, tyrosinase/albinisme.
  \item[$\square$] Je sais rédiger une interprétation d'expérience en quatre étapes : observation, hypothèse, résultat, conclusion.
  \item[$\square$] Je sais que les enzymes assurent le renouvellement permanent des molécules de la cellule (anabolisme et catabolisme).
  \item[$\square$] Je vérifie toujours les unités et je légende les axes lorsque j'exploite un graphique.
\end{itemize}
\end{aretenir}

\findecours

\end{document}
